\chapter{一元函数积分学}\label{ch:3}
\begin{tcolorbox}[title=章首知识地图,colframe=Ink]
原函数 $\to$ 换元与分部 $\to$ 有理化 $\to$ 定积分与反常积分 $\to$ 对称、递推、含参求导及不等式。
\end{tcolorbox}
\begin{examtip}考研主线：换元、分部、定积分性质、反常积分与几何应用。常见题型：凑微分、部分分式、区间再现及变限求导。竞赛拓展：双元代换、Wallis递推、含参求导和积分不等式。上述为学习任务分层，具体范围以所报考科目的当年大纲为准。\end{examtip}
\section{原函数与第一、第二换元}\label{sec:3:1}
\subsection{概念、定理与推导}

\begin{definition}{原函数与不定积分}区间$I$上$F\prime=f$时称$F$为$f$的原函数，$\int f\dd x=F+C$。同一区间内两原函数之差为常数；不连通定义域上的常数可分支选择。连续函数有原函数。\end{definition}
\begin{theorem}{换元公式}设$F\prime=f$，$\varphi$在所用区间可导且值落在$F$的定义区间。第一换元为$\int f(\varphi(x))\varphi\prime(x)\dd x=F(\varphi(x))+C$。第二换元$x=\psi(t)$后积分为$\int f(\psi(t))\psi\prime(t)\dd t$，回代需选定可逆分支。\end{theorem}
\begin{proof}第一换元来自链式法则；第二换元对所得原函数作复合求导验证。定积分换元时同时变上下限，可避免回代。\end{proof}
常见原函数：$\int x^\alpha\dd x=\frac{x^{\alpha+1}}{\alpha+1}+C$（$\alpha\ne-1$且在实值可导区间），$\int\frac{\dd x}{x}=\ln|x|+C$，$\int\frac{\dd x}{1+x^2}=\arctan x+C$。

\begin{method}{原函数与第一、第二换元}条件：代换可导且回代分支明确。步骤：先凑微分，不能凑时用根式或三角代换；最后求导验算。易错点：忘记$\dd x$的比例因子。\end{method}
\subsection{例题与方法比较}
\begin{example}{基础例题｜凑微分}{难度：$\star$}
\label{ex:3:1}\textbf{题目} 求$\int2x\cos(x^2)\dd x$。\par
\textbf{思路} 外层余弦的内层导数已经出现。\par
\textbf{解答} 方法一：令$u=x^2$，得$\int\cos u\dd u=\sin(x^2)+C$。方法二：直接识别$[\sin(x^2)]\prime=2x\cos(x^2)$。\par
\textbf{方法点评} 换元显示步骤，逆链式法则适合熟练计算。\par
\textbf{易错点或变式} 不能把$\cos(x^2)$写成$\cos^2x$。
\end{example}
\begin{example}{考研风格例题｜根式}{难度：$\star\star$}
\label{ex:3:2}\textbf{题目} 求$\int\frac{\dd x}{\sqrt{a^2-x^2}}$，$a>0,|x|<a$。\par
\textbf{思路} 用三角参数化圆。\par
\textbf{解答} 方法一：$x=a\sin t$，取$|t|<\frac{\pi}{2}$，根式为$a\cos t$，积分$t+C=\arcsin(\frac{x}{a})+C$。方法二：直接对$\arcsin(\frac{x}{a})$求导验证。\par
\textbf{方法点评} 代换说明反三角函数如何出现。\par
\textbf{易错点或变式} $\sqrt{\cos^2t}=|\cos t|$，分支选择不能省去。
\end{example}
\begin{example}{考研风格例题｜双元代换}{难度：$\star\star$}
\label{ex:3:3}\textbf{题目} 求$\int\frac{\dd x}{\sqrt{x^2+a^2}}$，$a>0$。\par
\textbf{思路} 利用两个相关量$x$与根式的和。\par
\textbf{解答} 方法一：令$u=x+\sqrt{x^2+a^2}>0$，则$\frac{\dd u}{u}=\frac{\dd x}{\sqrt{x^2+a^2}}$，答案$\ln(x+\sqrt{x^2+a^2})+C$。方法二：$x=a\sinh t$，积分变为$t+C=\operatorname{arsinh}(\frac{x}{a})+C$，与前式只差常数$\ln a$。\par
\textbf{方法点评} 双元思路利用共轭乘积恒定；双曲代换系统地消去平方和。\par
\textbf{易错点或变式} 使用“双元法”时，必须写出实际代换等式。
\end{example}
\section{分部积分与有理函数}\label{sec:3:2}

\begin{method}{部分分式的完整模板}
多项式除法后，对线性因子$(x-a)^m$配置$\sum_{k=1}^m \frac{A_k}{(x-a)^k}$；对不可约二次因子$(x^2+px+q)^m$配置$\sum_{k=1}^m\frac{B_kx+C_k}{(x^2+px+q)^k}$。比较多项式系数求未知数。二次分子的线性项先拆为分母导数的倍数与常数，分别产生对数与反正切；高次幂可用递推。
\end{method}

\subsection{概念、定理与推导}

\begin{theorem}{分部与部分分式}$\int u\dd v=uv-\int v\dd u$。有理函数先作多项式除法，真分式按实数域线性因子和不可约二次因子的幂分解，分子次数低于相应因子次数。\end{theorem}
\begin{proof}分部积分是乘积求导公式的积分。部分分式由多项式恒等式配系数，重根必须包含从一次到最高次的全部分母幂。\end{proof}
递推型积分应先定义$I_n$与初值，再由分部积分把阶数降低；循环型积分须把重新出现的原积分移至左侧。

\begin{method}{分部积分与有理函数}条件：分部两因子可导，分式远离极点。步骤：多项式乘指数三角通常让多项式求导；对数通常作$u$。易错点：重复不可约二次因子的分子仍须一次式。\end{method}
\subsection{例题与方法比较}
\begin{example}{考研风格例题｜分部}{难度：$\star\star$}
\label{ex:3:4}\textbf{题目} 求$\int x\e^x\dd x$。\par
\textbf{思路} 多项式求导后次数下降。\par
\textbf{解答} 方法一：$u=x,\dd v=\e^x\dd x$，得$x\e^x-\e^x+C$。方法二：设原函数为$(Ax+B)\e^x$，比较导数$(Ax+A+B)\e^x=x\e^x$，得$A=1,B=-1$。\par
\textbf{方法点评} 分部通用，待定系数利用封闭函数族。\par
\textbf{易错点或变式} 减号来自分部公式。
\end{example}
\begin{example}{考研风格例题｜部分分式}{难度：$\star\star$}
\label{ex:3:5}\textbf{题目} 求$\int\frac{\dd x}{x(x+1)}$。\par
\textbf{思路} 分母两个一次因子。\par
\textbf{解答} 方法一：$\frac{1}{x(x+1)}=\frac{1}{x}-\frac{1}{x+1}$，答案$\ln|x|-\ln|x+1|+C$。方法二：令$u=\frac{x}{x+1}$，则$\frac{\dd u}{u}=\frac{\dd x}{x(x+1)}$，得$\ln|u|+C$。\par
\textbf{方法点评} 部分分式机械可靠；比值代换揭示对数导数。\par
\textbf{易错点或变式} 每个避开$0,-1$的区间分别讨论常数。
\end{example}
\begin{example}{考研风格例题｜循环分部}{难度：$\star\star$}
\label{ex:3:6}\textbf{题目} 求$\int\e^x\cos x\dd x$。\par
\textbf{思路} 分部两次产生原积分。\par
\textbf{解答} 记$I$为该积分。一次分部得$I=\e^x\cos x+\int\e^x\sin x\dd x$；第二项等于$\e^x\sin x-I$。因此$2I=\e^x(\sin x+\cos x)$，$I=\frac{\e^x(\sin x+\cos x)}{2}+C$。另设$\e^x(A\cos x+B\sin x)$求导比较可得$A=B=\frac{1}{2}$。\par
\textbf{方法点评} 循环分部是解线性方程，待定系数更快。\par
\textbf{易错点或变式} 第二次分部的负号决定左边是$2I$。
\end{example}
\section{三角有理式与简单无理式}\label{sec:3:3}
\subsection{概念、定理与推导}

\begin{theorem}{万能代换与根式有理化}令$t=\tan(\frac{x}{2})$，则
\begin{equation}\sin x=\dfrac{2t}{1+t^2},\quad\cos x=\dfrac{1-t^2}{1+t^2},\quad\dd x=\dfrac{2\dd t}{1+t^2}.\end{equation}
故$R(\sin x,\cos x)$化为$t$的有理函数。对$R(x,\sqrt[m]{ax+b})$，当$a\ne0$、$m$为正整数时，可在所选实支设$t=\sqrt[m]{ax+b}$；$a=0$时根式为常量，另行简化。\end{theorem}
\begin{proof}半角公式由$\sin2u,\cos2u$除以$\cos^2u$得到，微分来自正切导数。根式代换给$x=\frac{t^m-b}{a}$，直接消去根式。\end{proof}
万能代换不是最短代换。若含单独的$\sin x\dd x$或$\cos x\dd x$，先考虑$u=\cos x$或$u=\sin x$。

\begin{method}{三角有理式与简单无理式}条件：半角坐标避开$\tan(\frac{x}{2})$极点。步骤：先找奇偶三角幂结构，再考虑万能代换；无理式选共同根次。易错点：跨越代换极点时要分段。\end{method}
\subsection{例题与方法比较}
\begin{example}{考研风格例题｜三角有理式}{难度：$\star\star$}
\label{ex:3:7}\textbf{题目} 求$\int\frac{\dd x}{2+\cos x}$。\par
\textbf{思路} 分母为正，半角化为二次有理式。\par
\textbf{解答} $t=\tan(\frac{x}{2})$给$\int\frac{2\dd t}{t^2+3}=\frac2{\sqrt3}\arctan(\frac{t}{\sqrt3})+C$。在任一避开半角极点的区间回代即可。跨越极点需调整积分常数以保持原函数连续。另一方法：令$u=\frac{x}{2}$，原积分为$2\int\frac{\dd u}{3\cos^2u+\sin^2u}$；把分子分母同除以$\cos^2u$，直接配成$\dd(\tan u)$并得到同一原函数。此法从倍角公式导出代换，而非直接套表。\par
\textbf{方法点评} 万能代换稳定，但反正切分支必须跟踪。\par
\textbf{易错点或变式} 原函数公式不应未经分段宣称在整个实轴单一分支成立。
\end{example}
\begin{example}{基础例题｜根式有理化}{难度：$\star$}
\label{ex:3:8}\textbf{题目} 求$\int\frac{\dd x}{1+\sqrt x}$，$x>0$。\par
\textbf{思路} 令根式成为新变量。\par
\textbf{解答} $t=\sqrt x$，$\dd x=2t\dd t$，积分$=2\int[1-\frac{1}{1+t}]\dd t=2\sqrt x-2\ln(1+\sqrt x)+C$。另一方法：设原函数形如$A\sqrt x+B\ln(1+\sqrt x)$，求导通分后分子为$A(1+\sqrt x)+B$，分母$2\sqrt x(1+\sqrt x)$。比较常数与$\sqrt x$系数得$A=2,B=-2$。\par
\textbf{方法点评} 换元后先做除法再积分。\par
\textbf{易错点或变式} 只代分母不代微分会丢失$2t$。
\end{example}
\begin{example}{考研风格例题｜三角幂}{难度：$\star\star$}
\label{ex:3:9}\textbf{题目} 求$\int\sin^3x\cos^2x\dd x$。\par
\textbf{思路} 留一个$\sin x$，其余化成余弦。\par
\textbf{解答} 令$u=\cos x$，$\dd u=-\sin x\dd x$，得$-\int(1-u^2)u^2\dd u=-\frac{u^3}{3}+\frac{u^5}{5}+C$。回代为$-\frac{\cos^3x}{3}+\frac{\cos^5x}{5}+C$。第二种方法：用积化和差写$\sin^3x\cos^2x=\frac{\sin x}{8}+\frac{\sin3x}{16}-\frac{\sin5x}{16}$，逐项积分得$-\frac{\cos x}{8}-\frac{\cos3x}{48}+\frac{\cos5x}{80}+C$。通过三倍角、五倍角展开可验证两原函数仅差常数。\par
\textbf{方法点评} 奇次三角幂通常比半角代换省算。\par
\textbf{易错点或变式} 变式：若两幂均偶，优先降幂公式。
\end{example}
\section{定积分、反常积分与应用}\label{sec:3:4}

\begin{theorem}{积分中值与变限求导}若$f$在$[a,b]$连续、$w$可积且不变号，则存在$\xi\in[a,b]$使$\int_a^bfw=f(\xi)\int_a^bw$。若$f$连续且上下限可导，则
\[\dfrac\dd{\dd x}\int_{u(x)}^{v(x)}f(t)\dd t=f(v(x))v\prime(x)-f(u(x))u\prime(x).\]
\end{theorem}
\begin{proof}当$w\ge0$时，用$f$的最小最大值夹住加权积分；若权积分为正，介值性提供相应$\xi$，为零时两边同为零。变限公式将积分写成两个变上限积分之差再用链式法则。\end{proof}

\subsection{概念、定理与推导}

\begin{definition}{Riemann积分与反常积分}有界函数的带标记分割和在网格趋零时趋于同一数，称该数为定积分。连续函数与有限间断点的有界函数可积。$\int_a^\infty f=\lim_{R\to\infty}\int_a^R f$；内点奇性须两侧积分分别收敛。\end{definition}
\begin{theorem}{微积分基本定理}连续$f$的变上限积分$F(x)=\int_a^x f(t)\dd t$满足$F\prime=f$；若$G\prime=f$，则$\int_a^bf=G(b)-G(a)$。\end{theorem}
\begin{proof}差商为$\frac1h\int_x^{x+h}f(t)\dd t$，连续性使其趋于$f(x)$；再用原函数只差常数得到Newton--Leibniz公式。\end{proof}
几何量：平面面积$\int|f-g|\dd x$，绕$x$轴的圆盘体积$\pi\int f^2\dd x$，弧长$\int\sqrt{1+(f\prime)^2}\dd x$，旋转曲面面积$2\pi\int|f|\sqrt{1+(f\prime)^2}\dd x$。正函数反常积分可用比较；$\int_1^\infty x^{-p}\dd x$当且仅当$p>1$收敛，$\int_0^1x^{-p}\dd x$当且仅当$p<1$收敛。

\begin{method}{定积分、反常积分与应用}条件：先判断普通或反常积分。步骤：有对称区间先做$x\mapsto a+b-x$；应用题先画截面。易错点：主值抵消不等于反常积分收敛。\end{method}
\subsection{例题与方法比较}
\begin{example}{考研风格例题｜区间再现}{难度：$\star\star$}
\label{ex:3:10}\textbf{题目} 求$I=\int_0^\pi x\sin x\dd x$。\par
\textbf{思路} 对称变换将$x$变成$\pi-x$。\par
\textbf{解答} 方法一：替换$x=\pi-t$得$I=\int_0^\pi(\pi-t)\sin t\dd t$，相加$2I=\pi\int_0^\pi\sin t\dd t=2\pi$，故$I=\pi$。方法二：分部得$[-x\cos x+\sin x]_0^\pi=\pi$。\par
\textbf{方法点评} 对称法可推广到任意关于中点对称的权函数。\par
\textbf{易错点或变式} 再现后积分限与微分负号要同步处理。
\end{example}
\begin{example}{考研风格例题｜反常判敛}{难度：$\star\star$}
\label{ex:3:11}\textbf{题目} 判断$\int_0^\infty\frac{\dd x}{\sqrt{x}(1+x)}$并计算。\par
\textbf{思路} 两端分别比较，再令$x=t^2$。\par
\textbf{解答} 在0附近像$x^{-\frac{1}{2}}$，无穷远像$x^{-\frac{3}{2}}$，均可积。设$x=t^2$，积分$2\int_0^\infty\frac{\dd t}{1+t^2}=\pi$。第二种方法：令$x=\tan^2u$，$u\in(0,\frac{\pi}{2})$，原积分直接变成$2\int_0^{\frac{\pi}{2}}\dd u=\pi$。前者化有理式，后者利用三角恒等式消去整个分母。\par
\textbf{方法点评} 先判敛保证变换的是合法极限。\par
\textbf{易错点或变式} 括号明确表示分母为$\sqrt x(1+x)$。
\end{example}
\begin{example}{基础例题｜旋转体}{难度：$\star$}
\label{ex:3:12}\textbf{题目} 区域$0\le y\le x,0\le x\le1$绕$x$轴旋转，求体积。\par
\textbf{思路} 圆盘或柱壳均可。\par
\textbf{解答} 圆盘法$V=\pi\int_0^1x^2\dd x=\frac{\pi}{3}$。柱壳法按$y$切片：半径$y$，壳长$1-y$，$V=2\pi\int_0^1y(1-y)\dd y=\frac{\pi}{3}$。\par
\textbf{方法点评} 两种切片方向验证同一几何体。\par
\textbf{易错点或变式} 壳长不是$y$。
\end{example}
\section{Wallis、含参积分与积分不等式}\label{sec:3:5}
\subsection{概念、定理与推导}

\begin{theorem}{Wallis递推}令$I_n=\int_0^{\frac{\pi}{2}}\sin^n x\dd x$。$I_0=\frac{\pi}{2},I_1=1$，$n\ge2$时$I_n=\frac{(n-1)I_{n-2}}{n}$。因此$I_{2m}=\frac{(2m-1)!!}{(2m)!!}\frac\pi2$，$I_{2m+1}=\frac{(2m)!!}{(2m+1)!!}$。\end{theorem}
\begin{proof}对$\sin^{n-1}x\sin x$分部，边界项为零，得到$I_n=(n-1)(I_{n-2}-I_n)$，移项即得。\end{proof}
\begin{theorem}{含参求导与Cauchy--Schwarz}若$f,\partial_tf$在紧矩形上连续，则$\frac{\dd}{\dd t}\int_a^bf(x,t)\dd x=\int_a^b\partial_tf(x,t)\dd x$。变限时另加$f(b(t),t)b\prime(t)-f(a(t),t)a\prime(t)$。若$f,g$在$[a,b]$上平方可积，则\[\left(\int_a^b f(x)g(x)\dd x\right)^2\le\left(\int_a^b f(x)^2\dd x\right)\left(\int_a^b g(x)^2\dd x\right).\]\end{theorem}
\begin{proof}连续偏导给差商在紧区间一致趋近，可交换积分与极限。Cauchy--Schwarz由$\int(f-\lambda g)^2\ge0$对$\lambda$的二次式判别式不正得到。反常含参积分还需可积支配等额外条件。\end{proof}

\begin{method}{Wallis、含参积分与积分不等式}条件：含参运算须验证一致控制，递推须有初值。步骤：对参数求导使被积式变简单，再积分并用特殊参数定常数。易错点：反常端点不能套紧矩形条件。\end{method}
\subsection{例题与方法比较}
\begin{example}{竞赛风格例题｜Wallis应用}{难度：$\star\star\star$}
\label{ex:3:13}\textbf{题目} 求$\int_0^{\frac{\pi}{2}}\sin^6x\dd x$。\par
\textbf{思路} 连续降二阶。\par
\textbf{解答} $I_6=\frac56\frac34\frac12\frac\pi2=\frac{5\pi}{32}$。另一法：降幂展开$\sin^6x=\frac{10-15\cos2x+6\cos4x-\cos6x}{32}$，余弦项积分均为零。\par
\textbf{方法点评} 递推适合任意整数幂，降幂适合小次数。\par
\textbf{易错点或变式} 偶数幂的初值是$\frac{\pi}{2}$。
\end{example}
\begin{example}{竞赛风格例题｜含参求导}{难度：$\star\star\star$}
\label{ex:3:14}\textbf{题目} $a>-1$，求$F(a)=\int_0^1\frac{x^a-1}{\ln x}\dd x$。\par
\textbf{思路} 对$a$求导消去对数分母。\par
\textbf{解答} 固定任意紧参数区间$[\alpha,\beta]\subset(-1,\infty)$，导数$x^a$被可积的$x^\alpha$控制；在$x=1$原式取连续极限$a$。因此$F\prime(a)=\int_0^1x^a\dd x=\frac{1}{a+1}$。$F(0)=0$给$F(a)=\ln(a+1)$。\par
\textbf{方法点评} 支配条件使反常含参求导严格成立。\par
\textbf{易错点或变式} 参数$a=-1$不在收敛范围。
\end{example}
\begin{example}{竞赛风格例题｜积分估计}{难度：$\star\star\star$}
\label{ex:3:15}\textbf{题目} 若$f$在$[0,1]$连续且$\int_0^1f=1$，证明$\int_0^1f^2\ge1$并判等。\par
\textbf{思路} 均值约束下常函数平方积分最小。\par
\textbf{解答} 方法一：Cauchy--Schwarz，$1=(\int f\cdot1)^2\le\int f^2$。方法二：$\int(f-1)^2=\int f^2-1\ge0$。等号要求连续非负函数$(f-1)^2$积分为零，即$f\equiv1$。\par
\textbf{方法点评} 平方配凑给出误差的精确表达。\par
\textbf{易错点或变式} 不需假设$f\ge0$。
\end{example}
\begin{warning}原函数不是所有函数都有的初等表达式；$\int\e^{-x^2}\dd x$不能强求初等形式。反常积分须分别处理每个奇点；两侧发散不能靠抵消定义为普通积分。\end{warning}


\section{不定积分的完整方法库}\label{mat:integrals_methods}
\begin{examtip}本节用可执行的代换、恒等式与递推组织积分技巧。每个原函数都在指定区间内理解；最后以求导核验。\end{examtip}
\subsection{原函数、常数与分支修正}\label{mat:integrals_methods:i01}
\begin{method}{原函数、常数与分支修正}\textbf{适用条件} 定义域可能不连通。\par\textbf{实施与依据} 每个连通区间保留独立常数；绝对值与反函数分支单独处理。\end{method}
\begin{example}{专题例题｜原函数、常数与分支修正}{难度：$\star\star$}
\textbf{题目} 说明$\int\frac{\dd x}{x}$在$\R\setminus\{0\}$上的全部原函数。\par\textbf{思路与解答} $x>0$为$\ln x+C_+$，$x<0$为$\ln(-x)+C_-$，$C_+,C_-$相互独立。不能因写一个$+C$就强迫跨越奇点的常数一致。\par\textbf{方法点评与边界} 考试书写应显式保留定义域、参数与分支条件，不把条件修正交给读者猜测。
\end{example}
\subsection{分项、重组与凑微分}\label{mat:integrals_methods:i02}
\begin{method}{分项、重组与凑微分}\textbf{适用条件} 分子可拆成分母导数与余项。\par\textbf{实施与依据} 先进行代数分解，再应用逆链式法则。\end{method}
\begin{example}{专题例题｜分项、重组与凑微分}{难度：$\star\star$}
\textbf{题目} 求$\int\frac{2x+3}{x^2+1}\dd x$。\par\textbf{思路与解答} 拆为$\ln(x^2+1)+3\arctan x+C$。求导后分子恢复$2x+3$。\par\textbf{方法点评与边界} 拆分依据线性，不允许把商的积分拆成两个积分的商。
\end{example}
\subsection{第一换元：内层导数}\label{mat:integrals_methods:i03}
\begin{method}{第一换元：内层导数}\textbf{适用条件} 有$f(g(x))g\prime(x)$。\par\textbf{实施与依据} 用$u=g(x)$连同微分整体替换。\end{method}
\begin{example}{专题例题｜第一换元：内层导数}{难度：$\star\star$}
\textbf{题目} 求$\int\frac{\ln x\,\dd x}{x}$，$x>0$。\par\textbf{思路与解答} $u=\ln x$，积分$\int u\dd u=\frac{u^2}{2}+C$。也可逆识别$\frac12(\ln x)^2$的导数。\par\textbf{方法点评与边界} 对数真数与换元后的区间同时记录。
\end{example}
\subsection{第二换元：主动参数化}\label{mat:integrals_methods:i04}
\begin{method}{第二换元：主动参数化}\textbf{适用条件} 根式可被三角或双曲恒等式消去。\par\textbf{实施与依据} 选择$x=a\sin t,a\tan t,a\sinh t$等适当分支。\end{method}
\begin{example}{专题例题｜第二换元：主动参数化}{难度：$\star\star$}
\textbf{题目} 求$\int\sqrt{a^2-x^2}\dd x$，$a>0,|x|<a$。\par\textbf{思路与解答} 令$x=a\sin t$，得$a^2\int\cos^2t\dd t$，结果$\frac12[x\sqrt{a^2-x^2}+a^2\arcsin(\frac{x}{a})]+C$。\par\textbf{方法点评与边界} 回代后求导可检验两个项的抵消。
\end{example}
\subsection{分部积分的选择}\label{mat:integrals_methods:i05}
\begin{method}{分部积分的选择}\textbf{适用条件} 乘积中有可降低复杂度的因子。\par\textbf{实施与依据} 对数或反三角常作求导因子，多项式反复降阶。\end{method}
\begin{example}{专题例题｜分部积分的选择}{难度：$\star\star$}
\textbf{题目} 求$\int\arctan x\dd x$。\par\textbf{思路与解答} 取$u=\arctan x,\dd v=\dd x$，得$x\arctan x-\int \frac{x}{1+x^2}\dd x=x\arctan x-\frac12\ln(1+x^2)+C$。\par\textbf{方法点评与边界} 经验顺序不能代替实际检查剩余积分是否更简单。
\end{example}
\subsection{多次分部与表格法}\label{mat:integrals_methods:i06}
\begin{method}{多次分部与表格法}\textbf{适用条件} 多项式乘指数、三角函数。\par\textbf{实施与依据} 交替符号排列多项式导数与另一因子的连续原函数。\end{method}
\begin{example}{专题例题｜多次分部与表格法}{难度：$\star\star$}
\textbf{题目} 求$\int x^2\e^{ax}\dd x$，$a\ne0$。\par\textbf{思路与解答} 两次分部给$\e^{ax}(\frac{x^2}{a}-\frac{2x}{a^2}+\frac{2}{a^3})+C$。对$a=0$另为$\frac{x^3}{3}+C$。\par\textbf{方法点评与边界} 常参数为零时必须回到未约化的表达式分类。
\end{example}
\subsection{循环积分与待定原函数}\label{mat:integrals_methods:i07}
\begin{method}{循环积分与待定原函数}\textbf{适用条件} 导数封闭的指数三角族。\par\textbf{实施与依据} 循环两次后解关于原积分的线性方程，或设同类试解。\end{method}
\begin{example}{专题例题｜循环积分与待定原函数}{难度：$\star\star$}
\textbf{题目} 求$\int\e^{ax}\sin bx\dd x$，$a^2+b^2>0$。\par\textbf{思路与解答} 设原函数$\e^{ax}(A\sin bx+B\cos bx)$，比较系数得$A=\frac{a}{a^2+b^2},B=-\frac{b}{a^2+b^2}$。\par\textbf{方法点评与边界} 该公式也由循环分部获得，两法都应回代验算。
\end{example}
\subsection{部分分式：一次与重因子}\label{mat:integrals_methods:i08}
\begin{method}{部分分式：一次与重因子}\textbf{适用条件} 真有理分式。\par\textbf{实施与依据} 一次重因子逐幂配置常数分子，通分比较系数。\end{method}
\begin{example}{专题例题｜部分分式：一次与重因子}{难度：$\star\star$}
\textbf{题目} 求$\int\frac{\dd x}{x(x-1)^2}$。\par\textbf{思路与解答} $\frac{1}{x(x-1)^2}=\frac{1}{x}-\frac{1}{x-1}+\frac{1}{(x-1)^2}$，故原函数$\ln|x|-\ln|x-1|-\frac{1}{x-1}+C$。\par\textbf{方法点评与边界} 重根不只保留最高次项。
\end{example}
\subsection{不可约二次因子}\label{mat:integrals_methods:i09}
\begin{method}{不可约二次因子}\textbf{适用条件} 分母为正定二次或其幂。\par\textbf{实施与依据} 配方，分子拆成分母导数与常数。\end{method}
\begin{example}{专题例题｜不可约二次因子}{难度：$\star\star$}
\textbf{题目} 求$\int\frac{x+1}{x^2+2x+5}\dd x$。\par\textbf{思路与解答} 分子恰为分母导数一半，结果$\frac12\ln(x^2+2x+5)+C$。若分子改成$x+2$，再加$\frac12\arctan(\frac{x+1}{2})$。\par\textbf{方法点评与边界} 二次分式的分子一般必须保留一次项。
\end{example}
\subsection{三角奇偶幂}\label{mat:integrals_methods:i10}
\begin{method}{三角奇偶幂}\textbf{适用条件} $\sin^mx\cos^nx$的某个指数为奇数。\par\textbf{实施与依据} 留一个奇函数微分，其余利用平方和改写。\end{method}
\begin{example}{专题例题｜三角奇偶幂}{难度：$\star\star$}
\textbf{题目} 求$\int\sin^2x\cos^3x\dd x$。\par\textbf{思路与解答} 令$u=\sin x$，积分为$\int u^2(1-u^2)\dd u=\frac{\sin^3x}{3}-\frac{\sin^5x}{5}+C$。\par\textbf{方法点评与边界} 若两指数均偶数，改用降幂公式。
\end{example}
\subsection{倍角、降幂与积化和差}\label{mat:integrals_methods:i11}
\begin{method}{倍角、降幂与积化和差}\textbf{适用条件} 偶次三角幂或不同频率乘积。\par\textbf{实施与依据} 用倍角与积化和差把乘积变为有限和。\end{method}
\begin{example}{专题例题｜倍角、降幂与积化和差}{难度：$\star\star$}
\textbf{题目} 求$\int\sin2x\cos3x\dd x$。\par\textbf{思路与解答} $\sin2x\cos3x=\frac{\sin5x-\sin x}{2}$，结果$-\frac{\cos5x}{10}+\frac{\cos x}{2}+C$。\par\textbf{方法点评与边界} 负频率项的奇偶性决定符号。
\end{example}
\subsection{三角万能代换}\label{mat:integrals_methods:i12}
\begin{method}{三角万能代换}\textbf{适用条件} 一般三角有理式。\par\textbf{实施与依据} $t=\tan(\frac{x}{2})$，将$\sin x,\cos x,\dd x$同时替换成有理式。\end{method}
\begin{example}{专题例题｜三角万能代换}{难度：$\star\star$}
\textbf{题目} 求$\int\frac{\dd x}{1+\sin x}$。\par\textbf{思路与解答} 可直接乘$1-\sin x$，化为$\sec^2x-\sec x\tan x$，原函数$\tan x-\sec x+C$。万能代换也给$2\int\frac{\dd t}{(t+1)^2}=-\frac{2}{t+1}+C$，在共同定义的连通区间上两式只差常数。为覆盖$\cos x=0,\sin x=1$这些原函数连续的点，应将$\tan x-\sec x$改写为$-\frac{\cos x}{1+\sin x}$；它在原分母非零的整个区间内可用。\par\textbf{方法点评与边界} 万能法系统，但恒等式法更短；均须避开原分母零点。
\end{example}
\subsection{正割、余割与递推}\label{mat:integrals_methods:i13}
\begin{method}{正割、余割与递推}\textbf{适用条件} 高次$\sec,\csc,\tan,\cot$。\par\textbf{实施与依据} 分部或$\tan^2x=\sec^2x-1$降低幂次。\end{method}
\begin{example}{专题例题｜正割、余割与递推}{难度：$\star\star$}
\textbf{题目} 求$J_n=\int\sec^nx\dd x$的递推。\par\textbf{思路与解答} 对$\sec^{n-2}x\cdot\sec^2x$分部，整理得$J_n=\frac{\sec^{n-2}x\tan x}{n-1}+\frac{(n-2)J_{n-2}}{n-1}$，$n>1$；初值$J_0=x,J_1=\ln|\sec x+\tan x|$。\par\textbf{方法点评与边界} 不定积分递推中的常数统一吸收到最终$C$。
\end{example}
\subsection{双曲函数有理式}\label{mat:integrals_methods:i14}
\begin{method}{双曲函数有理式}\textbf{适用条件} 平方差或指数成对出现。\par\textbf{实施与依据} 用$\cosh^2t-\sinh^2t=1$或$t=\e^x$。\end{method}
\begin{example}{专题例题｜双曲函数有理式}{难度：$\star\star$}
\textbf{题目} 求$\int\frac{\dd x}{\cosh x}$。\par\textbf{思路与解答} $t=\e^x>0$，$\cosh x=\frac{t+t^{-1}}{2}$，积分$2\arctan t+C=2\arctan\e^x+C$。\par\textbf{方法点评与边界} 双曲函数可统一还原为指数有理式。
\end{example}
\subsection{二次双元：圆型}\label{mat:integrals_methods:i15}
\begin{method}{二次双元：圆型}\textbf{适用条件} 两变量满足$u^2+v^2=c^2$。\par\textbf{实施与依据} 利用角变量或比值的微分，把两变量的微分关系闭合。\end{method}
\begin{example}{专题例题｜二次双元：圆型}{难度：$\star\star$}
\textbf{题目} 设$u^2+v^2=1,v>0$，积分$\int\frac{\dd u}{v}$。\par\textbf{思路与解答} 由$u=\sin t,v=\cos t$得到$t+C=\arcsin u+C$。若以$\frac{u}{v}$作参数，则$\dd(\frac{u}{v})=\frac{\dd u}{v^3}$，可反推其他双元积分。\par\textbf{方法点评与边界} “双元”是利用代数关系的组织方式，不是无条件公式。
\end{example}
\subsection{二次双元：双曲型}\label{mat:integrals_methods:i16}
\begin{method}{二次双元：双曲型}\textbf{适用条件} $v^2-u^2=a^2>0$。\par\textbf{实施与依据} 利用$(v+u)(v-u)=a^2$和对数微分。\end{method}
\begin{example}{专题例题｜二次双元：双曲型}{难度：$\star\star$}
\textbf{题目} 在$v=\sqrt{u^2+a^2}$下求$\int\frac{\dd u}{v}$。\par\textbf{思路与解答} $\dd\ln(u+v)=\frac{\dd u}{v}$，答案$\ln(u+v)+C$；$u=a\sinh t$给同一表达式加常数。\par\textbf{方法点评与边界} 平方根取正分支，因此$u+v>0$。
\end{example}
\subsection{双元降次与递推}\label{mat:integrals_methods:i17}
\begin{method}{双元降次与递推}\textbf{适用条件} $I_n=\int(x^2+a^2)^{-n}\dd x$。\par\textbf{实施与依据} 对$x(x^2+a^2)^{1-n}$求导，解出$I_n$。\end{method}
\begin{example}{专题例题｜双元降次与递推}{难度：$\star\star$}
\textbf{题目} 推导$I_n$递推，$a>0,n\ge2$。\par\textbf{思路与解答} 直接求导得$[\frac{x}{(x^2+a^2)^{n-1}}]\prime=\frac{3-2n}{(x^2+a^2)^{n-1}}+\frac{2(n-1)a^2}{(x^2+a^2)^n}$，故$I_n=\frac{x}{2(n-1)a^2(x^2+a^2)^{n-1}}+\frac{(2n-3)I_{n-1}}{2(n-1)a^2}$。\par\textbf{方法点评与边界} 用$n=2$可快速检验系数。
\end{example}
\subsection{分式构型与齐次比值}\label{mat:integrals_methods:i18}
\begin{method}{分式构型与齐次比值}\textbf{适用条件} 分子分母关于两元同次齐次。\par\textbf{实施与依据} 选$\frac{u}{v}$或$\frac{v}{u}$消除一个变量，微分由商法则推导。\end{method}
\begin{example}{专题例题｜分式构型与齐次比值}{难度：$\star\star$}
\textbf{题目} 求$\int\frac{\dd x}{x\sqrt{x^2-1}}$，$x>1$。\par\textbf{思路与解答} 令$t=\frac{\sqrt{x^2-1}}{x}$，$t\prime=\frac{1}{x^2\sqrt{x^2-1}}$，或更直接设$x=\sec\theta$，积分$\dd\theta$，得$\arccos(\frac{1}{x})+C$。\par\textbf{方法点评与边界} 先算换元微分，不能凭齐次外形省掉额外$x$因子。
\end{example}
\subsection{恒等式诱导与Möbius变换}\label{mat:integrals_methods:i19}
\begin{method}{恒等式诱导与Möbius变换}\textbf{适用条件} 互为共轭的和差，或分式线性变换。\par\textbf{实施与依据} 用$t=\frac{ax+b}{cx+d}$并反解，行列式$ad-bc\ne0$。\end{method}
\begin{example}{专题例题｜恒等式诱导与Möbius变换}{难度：$\star\star$}
\textbf{题目} 求$\int\frac{\dd x}{x^2-1}$。\par\textbf{思路与解答} 令$t=\frac{x-1}{x+1}$，有$\frac{\dd t}{t}=\frac{2\dd x}{x^2-1}$，故结果$\frac12\ln|t|+C$。部分分式得到同式。\par\textbf{方法点评与边界} 非零行列式和分母非零保证局部代换合法。
\end{example}
\subsection{双元化单元与共同根次}\label{mat:integrals_methods:i20}
\begin{method}{双元化单元与共同根次}\textbf{适用条件} 多重根号具有共同基底。\par\textbf{实施与依据} 用最小公倍根次一次消除多个分数指数。\end{method}
\begin{example}{专题例题｜双元化单元与共同根次}{难度：$\star\star$}
\textbf{题目} 求$\int\frac{\dd x}{\sqrt x+\sqrt[3]x}$，$x>0$。\par\textbf{思路与解答} 设$x=t^6$，积分$6\int \frac{t^3}{1+t}\dd t=6\int(t^2-t+1-\frac{1}{1+t})\dd t$。结果$2t^3-3t^2+6t-6\ln(1+t)+C$，回代$t=x^{\frac{1}{6}}$。\par\textbf{方法点评与边界} 统一根次比先分两次换元简洁。
\end{example}
\subsection{多重根号：和差参数}\label{mat:integrals_methods:i21}
\begin{method}{多重根号：和差参数}\textbf{适用条件} 出现$\sqrt{x+a}\pm\sqrt{x+b}$。\par\textbf{实施与依据} 利用两者和差的乘积为常数。\end{method}
\begin{example}{专题例题｜多重根号：和差参数}{难度：$\star\star$}
\textbf{题目} 求$\int\frac{\dd x}{\sqrt{x+1}+\sqrt x}$，$x>0$。\par\textbf{思路与解答} 乘共轭后 被积式为$\sqrt{x+1}-\sqrt x$，答案$\frac23[(x+1)^{\frac{3}{2}}-x^{\frac{3}{2}}]+C$。\par\textbf{方法点评与边界} 先做代数消根，通常优于引入复杂新变量。
\end{example}
\subsection{Euler根式替换}\label{mat:integrals_methods:i22}
\begin{method}{Euler根式替换}\textbf{适用条件} 一般$R(x,\sqrt{ax^2+bx+c})$。\par\textbf{实施与依据} $a>0$时令根式$=t+\sqrt a\,x$；$c>0$时也可令根式$=tx+\sqrt c$。平方后反解$x$。\end{method}
\begin{example}{专题例题｜Euler根式替换}{难度：$\star\star$}
\textbf{题目} 将$\sqrt{x^2+x+1}$有理化。\par\textbf{思路与解答} 设$\sqrt{x^2+x+1}=t+x$，平方得$x(1-2t)=t^2-1$，故$x=\frac{t^2-1}{1-2t}$，根式与$\dd x$均成为$t$的有理函数。\par\textbf{方法点评与边界} 平方可能引入分支，回代须确认原根式非负。
\end{example}
\subsection{指数单元法}\label{mat:integrals_methods:i23}
\begin{method}{指数单元法}\textbf{适用条件} $R(\e^{ax})$且$\dd x$可转成$\frac{\dd t}{t}$。\par\textbf{实施与依据} 取$t=\e^{ax}$，化有理积分。\end{method}
\begin{example}{专题例题｜指数单元法}{难度：$\star\star$}
\textbf{题目} 求$\int\frac{\dd x}{1+\e^x}$。\par\textbf{思路与解答} $t=\e^x$后$\int\frac{\dd t}{t(1+t)}=\ln t-\ln(1+t)+C=x-\ln(1+\e^x)+C$。\par\textbf{方法点评与边界} 若同时含$x$，换元后$x=\ln t$未必仍是有理式。
\end{example}
\subsection{有理积分中的倒数对称}\label{mat:integrals_methods:i24}
\begin{method}{有理积分中的倒数对称}\textbf{适用条件} 分母回文，多项式满足倒数对称。\par\textbf{实施与依据} 按$x\pm\frac{1}{x}$组合及其微分选变量。\end{method}
\begin{example}{专题例题｜有理积分中的倒数对称}{难度：$\star\star$}
\textbf{题目} 求$\int\frac{1-x^{-2}}{1+(x+x^{-1})^2}\dd x$，$x\ne0$。\par\textbf{思路与解答} 设$u=x+\frac{1}{x}$，$\dd u=(1-x^{-2})\dd x$，得$\arctan(x+\frac{1}{x})+C$。\par\textbf{方法点评与边界} 并非所有回文分式都直接可用，需要微分结构匹配。
\end{example}
\subsection{隐函数积分与微分方程}\label{mat:integrals_methods:i25}
\begin{method}{隐函数积分与微分方程}\textbf{适用条件} 曲线关系可给$x,y$参数式或导数。\par\textbf{实施与依据} 对约束求微分，将$y\dd x$与$\dd(xy)$结合。\end{method}
\begin{example}{专题例题｜隐函数积分与微分方程}{难度：$\star\star$}
\textbf{题目} 单位圆第一象限上求$\int y\dd x$。\par\textbf{思路与解答} $x=\cos t,y=\sin t$，$y\dd x=-\sin^2t\dd t$，积分$-\frac{t}{2}+\frac{\sin2t}{4}+C$，可回代成$\frac12[x\sqrt{1-x^2}+\arcsin x]+C$。\par\textbf{方法点评与边界} 不同参数方向产生的符号要与$\dd x$一致。
\end{example}
\subsection{配对与组合积分}\label{mat:integrals_methods:i26}
\begin{method}{配对与组合积分}\textbf{适用条件} 两式相加后出现可积结构。\par\textbf{实施与依据} 把未知积分与同族积分建立线性关系。\end{method}
\begin{example}{专题例题｜配对与组合积分}{难度：$\star\star$}
\textbf{题目} 令$I=\int\frac{\sin x}{\sin x+\cos x}\dd x$，求原函数。\par\textbf{思路与解答} 用$\sin x=\frac12[(\sin x+\cos x)-(\cos x-\sin x)]$，得$I=\frac{x}{2}-\frac12\ln|\sin x+\cos x|+C$。\par\textbf{方法点评与边界} 配对实际上是把分子拆成分母与分母导数。
\end{example}
\subsection{含参构造消去对数}\label{mat:integrals_methods:i27}
\begin{method}{含参构造消去对数}\textbf{适用条件} 对数幂可由参数导数产生。\par\textbf{实施与依据} 先求简单含参积分，再在有控制的参数区间求导。\end{method}
\begin{example}{专题例题｜含参构造消去对数}{难度：$\star\star$}
\textbf{题目} 求$\int_0^1x^a\ln x\dd x$，$a>-1$。\par\textbf{思路与解答} 从$F(a)=\frac{1}{a+1}$求导得$-\frac{1}{(a+1)^2}$；直接分部并取0端极限也得同值。\par\textbf{方法点评与边界} 反常端点处要有$x^\alpha|\ln x|$等可积控制。
\end{example}
\subsection{复数分解与取实部}\label{mat:integrals_methods:i28}
\begin{method}{复数分解与取实部}\textbf{适用条件} 指数三角组合或实不可约二次因子。\par\textbf{实施与依据} 先在复数域简化，再取实部虚部；对数须说明分支。\end{method}
\begin{example}{专题例题｜复数分解与取实部}{难度：$\star\star$}
\textbf{题目} 求$\int\e^x\cos2x\dd x$。\par\textbf{思路与解答} 取$\int\e^{(1+2i)x}\dd x$的实部，结果$\frac{\e^x(\cos2x+2\sin2x)}{5}+C$；待定系数亦可验证。\par\textbf{方法点评与边界} 复数技巧是实积分的简化工具，最终答案需实值且分支一致。
\end{example}
\subsection{伪椭圆构型与恰当代换}\label{mat:integrals_methods:i29}
\begin{method}{伪椭圆构型与恰当代换}\textbf{适用条件} 看似高次根式，但可化二次结构。\par\textbf{实施与依据} 先查$x^m$及其微分，避免仅凭四次根式认定非初等。\end{method}
\begin{example}{专题例题｜伪椭圆构型与恰当代换}{难度：$\star\star$}
\textbf{题目} 求$\int \frac{x\dd x}{\sqrt{1+x^4}}$。\par\textbf{思路与解答} 设$u=x^2$，得$\frac12\int\frac{\dd u}{\sqrt{1+u^2}}=\frac12\operatorname{arsinh}(x^2)+C$。真正的$\int\frac{\dd x}{\sqrt{1+x^4}}$一般属于椭圆积分。\par\textbf{方法点评与边界} 被积式的微分因子决定是否能降为初等积分。
\end{example}
\subsection{三元与椭圆积分}\label{mat:integrals_methods:i30}
\begin{method}{三元与椭圆积分}\textbf{适用条件} 根式对应不可降的三次或四次多项式。\par\textbf{实施与依据} 区分Legendre标准型的第一、第二类，不把特殊函数伪装为初等函数。\end{method}
\begin{example}{专题例题｜三元与椭圆积分}{难度：$\star\star$}
\textbf{题目} 写出$\int\frac{\dd x}{\sqrt{(1-x^2)(1-k^2x^2)}}$的标准表达，$0<k<1,|x|<1$。\par\textbf{思路与解答} 令$x=\sin\phi$，得到$F(\phi,k)=\int_0^\phi(1-k^2\sin^2t)^{-\frac{1}{2}}\dd t$，答案$F(\arcsin x,k)+C$。$k=0$退化为反正弦。\par\textbf{方法点评与边界} 这是完整的特殊函数表达，并不是省略计算。
\end{example}
\subsection{微分方程识别原函数}\label{mat:integrals_methods:i31}
\begin{method}{微分方程识别原函数}\textbf{适用条件} 原函数属于有限维函数族或其逆函数。\par\textbf{实施与依据} 令未知原函数满足$F\prime=f$，用函数族的微分关系解系数或降阶。\end{method}
\begin{example}{专题例题｜微分方程识别原函数}{难度：$\star\star$}
\textbf{题目} 求$\int\e^{x^2}(1+2x^2)\dd x$。\par\textbf{思路与解答} 注意$[x\e^{x^2}]\prime=\e^{x^2}(1+2x^2)$，故答案$x\e^{x^2}+C$。虽然$\int\e^{x^2}\dd x$非初等，当前组合却初等。\par\textbf{方法点评与边界} 不能从一个因子的非初等原函数推出乘积也非初等。
\end{example}

\section{特殊函数、复变积分与积分变换}\label{mat:special_integrals}
\begin{examtip}本节给出特殊函数的定义、微分关系、可计算算例与适用边界。\end{examtip}
\subsection{误差函数与Fresnel积分}\label{mat:special_integrals:erf}
\begin{method}{误差函数与Fresnel积分}\textbf{适用条件} 平方指数或二次相位振荡。\par\textbf{实施与依据} $\operatorname{erf}x=\frac2{\sqrt\pi}\int_0^x\e^{-t^2}\dd t$；$C(x)=\int_0^x\cos(\frac{\pi t^2}{2})\dd t,S(x)=\int_0^x\sin(\frac{\pi t^2}{2})\dd t$。\end{method}
\begin{example}{专题例题｜误差函数与Fresnel积分}{难度：$\star\star$}
\textbf{题目} 表达$\int\e^{-a x^2}\dd x$，$a>0$。\par\textbf{思路与解答} 令$u=\sqrt a\,x$，得到$\frac{\sqrt\pi\operatorname{erf}(\sqrt a\,x)}{2\sqrt a}+C$。导数定义直接检验。\par\textbf{方法点评与边界} 特殊函数的积分上下限可有不同约定，计算时须固定同一组定义。
\end{example}
\subsection{指数、正弦与余弦积分}\label{mat:special_integrals:sici}
\begin{method}{指数、正弦与余弦积分}\textbf{适用条件} 除以自变量的指数三角积分。\par\textbf{实施与依据} $\operatorname{Si}(x)=\int_0^x\frac{\sin t}{t}\dd t$；$\operatorname{Ci}(x)=\gamma+\ln x+\int_0^x\frac{\cos t-1}{t}\dd t$，$x>0$。实$\operatorname{Ei}$需避开或按主值处理0。\end{method}
\begin{example}{专题例题｜指数、正弦与余弦积分}{难度：$\star\star$}
\textbf{题目} 求$\int\frac{\sin(ax)}{x}\dd x$。\par\textbf{思路与解答} $a\ne0$时为$\operatorname{Si}(ax)+C$，因为链式求导给$\frac{\sin(ax)}{x}$；$a=0$时为常数。\par\textbf{方法点评与边界} 不能把$\int_0^x\frac{\cos t}{t}\dd t$当普通收敛积分。
\end{example}
\subsection{多重对数与指对双元}\label{mat:special_integrals:polylog}
\begin{method}{多重对数与指对双元}\textbf{适用条件} 反复出现$\frac{\ln(1-x)}{x}$。\par\textbf{实施与依据} 定义$\operatorname{Li}_s(z)=\sum_{n\ge1}\frac{z^n}{n^s}$，$|z|<1$；$\frac{\dd\operatorname{Li}_s(z)}{\dd z}=\frac{\operatorname{Li}_{s-1}(z)}{z}$。\end{method}
\begin{example}{专题例题｜多重对数与指对双元}{难度：$\star\star$}
\textbf{题目} 求$\int\frac{\ln(1-x)\dd x}{x}$，$|x|<1$。\par\textbf{思路与解答} $\operatorname{Li}_1(x)=-\ln(1-x)$，故答案$-\operatorname{Li}_2(x)+C$。从对数幂级数逐项积分也得到此表达。\par\textbf{方法点评与边界} 超出单位圆时需解析延拓及分支信息，不直接沿用实区间公式。
\end{example}
\subsection{Lambert W与隐式变量}\label{mat:special_integrals:lambert}
\begin{method}{Lambert W与隐式变量}\textbf{适用条件} 变量同时在线性因子及指数中。\par\textbf{实施与依据} $W(z)\e^{W(z)}=z$；在所选可微分支且$z\ne0,\,W(z)\ne-1$时，$W\prime=\frac{W}{z(1+W)}$。主实分支在0处由隐函数关系得到$W_0\prime(0)=1$，商式的$\frac{0}{0}$须按可去极限理解。\end{method}
\begin{example}{专题例题｜Lambert W与隐式变量}{难度：$\star\star$}
\textbf{题目} 求$\int W(x)\dd x$，在可微实分支区间。\par\textbf{思路与解答} 设$u=W(x)$，$x=u\e^u,\dd x=(u+1)\e^u\dd u$，积分$\e^u(u^2-u+1)+C$。求导回到$u$，再用隐式关系核验。\par\textbf{方法点评与边界} $x=-\frac{1}{\e}$是实分支的关键点；不同实分支需分别讨论。
\end{example}
\subsection{Gamma与Beta函数}\label{mat:special_integrals:gamma}
\begin{method}{Gamma与Beta函数}\textbf{适用条件} 无穷区间指数权、单位区间幂权。\par\textbf{实施与依据} $\Gamma(s)=\int_0^\infty t^{s-1}\e^{-t}\dd t$，$s>0$；$B(p,q)=\int_0^1t^{p-1}(1-t)^{q-1}\dd t$，实参数$p,q>0$；这两个条件分别保证0与1端点可积。\end{method}
\begin{example}{专题例题｜Gamma与Beta函数}{难度：$\star\star$}
\textbf{题目} 设$p,q>0$，证明$B(p,q)=\frac{\Gamma(p)\Gamma(q)}{\Gamma(p+q)}$。\par\textbf{思路与解答} 在正象限把Gamma乘积写成二重积分，令$x=ru,y=r(1-u)$，Jacobian为$r$。积分分离成$\Gamma(p+q)B(p,q)$，故结论成立；非负性保证换序合法。\par\textbf{方法点评与边界} 由分部得到$\Gamma(s+1)=s\Gamma(s)$、$\Gamma(\frac{1}{2})=\sqrt\pi$。
\end{example}
\subsection{Bessel函数及加权积分}\label{mat:special_integrals:bessel}
\begin{method}{Bessel函数及加权积分}\textbf{适用条件} 圆对称问题与Bessel方程。\par\textbf{实施与依据} 在$x>0,\nu\ge0$时，第一类Bessel函数由归一化级数
\[
J_\nu(x)=\sum_{k=0}^\infty
\dfrac{(-1)^k(\frac{x}{2})^{2k+\nu}}{k!\,\Gamma(k+\nu+1)}
\]
确定，而不仅由二阶方程的名称确定。它满足$x^2y\prime\prime+xy\prime+(x^2-\nu^2)y=0$；常用关系$[x^{\nu+1}J_{\nu+1}(x)]\prime=x^{\nu+1}J_\nu(x)$。\end{method}
\begin{example}{专题例题｜Bessel函数及加权积分}{难度：$\star\star$}
\textbf{题目} 求$\int xJ_0(x)\dd x$。\par\textbf{思路与解答} 取$\nu=0$，答案$xJ_1(x)+C$。关系可由$J_0=\sum\frac{(-1)^k(\frac{x}{2})^{2k}}{(k!)^2}$与$J_1$的级数逐项微分核验。\par\textbf{方法点评与边界} 第二类Bessel、虚宗量修正Bessel也由对应ODE与递推处理，不能不加符号调整照搬。
\end{example}
\subsection{Airy方程与积分}\label{mat:special_integrals:airy}
\begin{method}{Airy方程与积分}\textbf{适用条件} 线性转折点模型。\par\textbf{实施与依据} $\operatorname{Ai}$与$\operatorname{Bi}$满足$y\prime\prime-xy=0$，本册用以下初值固定标准归一化：
\[
\begin{aligned}
\operatorname{Ai}(0)&=\dfrac1{3^{\frac{2}{3}}\Gamma(\frac{2}{3})},&
\operatorname{Ai}'(0)&=-\dfrac1{3^{\frac{1}{3}}\Gamma(\frac{1}{3})},\\
\operatorname{Bi}(0)&=\sqrt3\,\operatorname{Ai}(0),&
\operatorname{Bi}'(0)&=-\sqrt3\,\operatorname{Ai}'(0).
\end{aligned}
\]
二阶初值问题唯一性确定这两个函数；仅写共同微分方程不能区分它们。\end{method}
\begin{example}{专题例题｜Airy方程与积分}{难度：$\star\star$}
\textbf{题目} 求$\int x\operatorname{Ai}(x)\dd x$。\par\textbf{思路与解答} 由定义方程$\operatorname{Ai}\prime\prime=x\operatorname{Ai}$，原函数为$\operatorname{Ai}\prime(x)+C$。\par\textbf{方法点评与边界} 特殊函数满足的微分方程可直接用于构造其原函数。
\end{example}
\subsection{Legendre多项式与正交积分}\label{mat:special_integrals:legendre}
\begin{method}{Legendre多项式与正交积分}\textbf{适用条件} 区间$[-1,1]$上的多项式谱。\par\textbf{实施与依据} $P_n(x)=\frac1{2^nn!}\frac{\dd^n}{\dd x^n}(x^2-1)^n$。分部$n$次可证明对低次多项式正交。\end{method}
\begin{example}{专题例题｜Legendre多项式与正交积分}{难度：$\star\star$}
\textbf{题目} 求$\int_{-1}^1xP_2(x)\dd x$。\par\textbf{思路与解答} $P_2=\frac{3x^2-1}{2}$为偶，乘$x$为奇，因此积分0。一般有$\int_{-1}^1P_nP_m=0$（$m\ne n$）。\par\textbf{方法点评与边界} 等次平方积分为$\frac{2}{2n+1}$，需要归一化时不可漏掉。
\end{example}
\subsection{超几何与Mathieu函数}\label{mat:special_integrals:hypergeom}
\begin{method}{超几何与Mathieu函数}\textbf{适用条件} 参数化二阶ODE及特殊函数积分。\par\textbf{实施与依据} Gauss超几何级数${}_2F_1(a,b;c;z)=\sum_{n=0}^\infty\frac{(a)_n(b)_n z^n}{(c)_nn!}$，先取$|z|<1$、$c$避开非正整数；其中$(a)_0=1$、$(a)_n=a(a+1)\cdots(a+n-1)$是升阶乘。Mathieu方程为$y\prime\prime+(a-2q\cos2x)y=0$。\end{method}
\begin{example}{专题例题｜超几何与Mathieu函数}{难度：$\star\star$}
\textbf{题目} 写出$\int(1-x^2)^{-\frac{1}{2}}\dd x$的超几何形式。\par\textbf{思路与解答} 二项式展开后逐项积分得$x\,{}_2F_1(\frac{1}{2},\frac{1}{2};\frac{3}{2};x^2)+C=\arcsin x+C$，$|x|<1$。Mathieu函数在$q=0$退化为常系数三角或指数函数。\par\textbf{方法点评与边界} 不为每个参数强行给初等形式；先定义函数及其微分方程，再利用微分关系积分。
\end{example}
\subsection{Cauchy积分与留数法}\label{mat:special_integrals:residue}
\begin{method}{Cauchy积分与留数法}\textbf{适用条件} 复平面上孤立奇点和可闭合路径。\par\textbf{实施与依据} 若函数在分段光滑简单闭曲线及其内部的一个开邻域内解析，则闭路积分为0；更一般地，在单连通解析区域内任意闭路积分为0。正向围道内部只有有限个孤立极点、边界无奇点时，积分为$2\pi i$乘留数和；还要估计闭合大弧。\end{method}
\begin{example}{专题例题｜Cauchy积分与留数法}{难度：$\star\star$}
\textbf{题目} 用留数求$\int_{-\infty}^\infty\frac{\dd x}{1+x^2}$。\par\textbf{思路与解答} 上半平面仅极点$i$，留数$\frac{1}{2i}$，围道积分为$\pi$；大半圆弧上的积分模不超过$\frac{\pi R}{R^2-1}\to0$，故实积分为$\pi$。\par\textbf{方法点评与边界} 弧积分消失必须给估计；路径上的极点需缩进或主值，不能直接套公式。
\end{example}
\subsection{Laplace与Fourier积分变换}\label{mat:special_integrals:laplace}
\begin{method}{Laplace与Fourier积分变换}\textbf{适用条件} 半轴增长可控的函数或全轴可积函数。\par\textbf{实施与依据} $\mathcal Lf(s)=\int_0^\infty\e^{-st}f(t)\dd t$，$\widehat f(\xi)=\int_\R f(x)\e^{-2\pi i x\xi}\dd x$。\end{method}
\begin{example}{专题例题｜Laplace与Fourier积分变换}{难度：$\star\star$}
\textbf{题目} 求$\mathcal L(t^n)(s)$，$s>0$、$n$为非负整数。\par\textbf{思路与解答} 令$u=st$，得到$\frac{\Gamma(n+1)}{s^{n+1}}=\frac{n!}{s^{n+1}}$。对变换求$s$导数也可从$\frac{1}{s}$递推得到。\par\textbf{方法点评与边界} 变换公式的收敛半平面是公式的一部分。
\end{example}

\section{代数约化、椭圆积分与特殊函数递推}\label{mat:advanced_integral_algebra}
\begin{examtip}补充有理和代数积分的算法视角，以及特殊函数公式的生成机制。\end{examtip}
\subsection{Hermite约化与代数恰当项}\label{mat:advanced_integral_algebra:hermite}
\begin{method}{Hermite约化与代数恰当项}\textbf{适用条件} 有理函数，或带代数根式的候选精确导数。\par\textbf{实施与依据} 先消高次极点，或设$\frac{P(x)}{\sqrt{Q(x)}}$、$P(x)\sqrt{Q(x)}$的导数匹配被积式，解多项式系数。\end{method}
\begin{example}{专题例题｜Hermite约化与代数恰当项}{难度：$\star\star$}
\textbf{题目} 求$\int\frac{2x^2+1}{\sqrt{x^2+1}}\dd x$。\par\textbf{思路与解答} 设$F=P\sqrt{x^2+1}$，则$F\prime=\frac{P\prime(x)(x^2+1)+xP(x)}{\sqrt{x^2+1}}$。取$P=x$恰好匹配，答案$x\sqrt{x^2+1}+C$。\par\textbf{方法点评与边界} 一般代数积分可能余下非初等部分；连分式与伪椭圆构造须按多项式恒等式核验，不能凭外观分类。
\end{example}
\subsection{第二、第三类椭圆积分}\label{mat:advanced_integral_algebra:elliptic23}
\begin{method}{第二、第三类椭圆积分}\textbf{适用条件} $0<k<1$且积分路径避开分母零点。\par\textbf{实施与依据} $E(\phi,k)=\int_0^\phi\sqrt{1-k^2\sin^2t}\dd t$；$\Pi(n;\phi,k)=\int_0^\phi\frac{\dd t}{(1-n\sin^2t)\sqrt{1-k^2\sin^2t}}$。\end{method}
\begin{example}{专题例题｜第二、第三类椭圆积分}{难度：$\star\star$}
\textbf{题目} 表达$\int\frac{\sqrt{1-k^2x^2}}{\sqrt{1-x^2}}\dd x$，$|x|<1$。\par\textbf{思路与解答} 令$x=\sin t$，积分变成$\int\sqrt{1-k^2\sin^2t}\dd t$，答案$E(\arcsin x,k)+C$。若再乘$\frac{1}{1-nx^2}$，不能直接写成$\Pi$。令$\phi=\arcsin x$，则对$n\ne0$，
\[
\begin{aligned}
&\int\dfrac{\sqrt{1-k^2x^2}}{(1-nx^2)\sqrt{1-x^2}}\dd x\\
&=\dfrac{k^2}{n}F(\phi,k)
+\left(1-\dfrac{k^2}{n}\right)\Pi(n;\phi,k)+C.
\end{aligned}
\]
依据是$\frac{1-k^2s^2}{1-ns^2}=\frac{k^2}{n}+\frac{1-\frac{k^2}{n}}{1-ns^2}$；$n=0$时仍为$E(\phi,k)+C$。第三类的纯标准核是第一类核另乘$(1-nx^2)^{-1}$。\par\textbf{方法点评与边界} 椭圆积分以振幅及模数为变量，不等同于椭圆函数的逆关系。
\end{example}
\subsection{第二类与修正Bessel积分}\label{mat:advanced_integral_algebra:modbessel}
\begin{method}{第二类与修正Bessel积分}\textbf{适用条件} 对应Bessel或修正Bessel方程的解。\par\textbf{实施与依据} 对$x>0$，$I_\nu$由$J_\nu$级数去掉$(-1)^k$而定义（本节用$\nu=0,1$）。第二类Bessel函数与修正第二类函数的标准归一化可写为
\[
Y_\nu=\dfrac{\cos(\pi\nu)J_\nu-J_{-\nu}}{\sin(\pi\nu)},
\qquad
K_\nu=\dfrac{\pi(I_{-\nu}-I_\nu)}{2\sin(\pi\nu)}
\quad(\nu\notin\mathbb Z).
\]
这里非整数负阶的$J_{-\nu},I_{-\nu}$仍用相应级数，Gamma函数按递推延拓；整数阶按$\nu$的可去极限定义。由此得到$[xY_1(x)]\prime=xY_0(x)$；$[xI_1(x)]\prime=xI_0(x)$；$[xK_1(x)]\prime=-xK_0(x)$。\end{method}
\begin{example}{专题例题｜第二类与修正Bessel积分}{难度：$\star\star$}
\textbf{题目} 求$\int xK_0(x)\dd x$，$x>0$。\par\textbf{思路与解答} 由修正Bessel递推，答案$-xK_1(x)+C$。这些公式可从定义ODE与相邻阶递推联立得到，负号是$K$与$I$的关键区别。\par\textbf{方法点评与边界} 在0附近第二类和$K$常有奇性，定积分还需检查端点。
\end{example}
\subsection{Mathieu积分与ODE还原}\label{mat:advanced_integral_algebra:mathieu}
\begin{method}{Mathieu积分与ODE还原}\textbf{适用条件} 周期系数方程$y\prime\prime+(a-2q\cos2x)y=0$。\par\textbf{实施与依据} 用微分方程把某些加权积分转为导数；$q=0$是常系数退化情形。\end{method}
\begin{example}{专题例题｜Mathieu积分与ODE还原}{难度：$\star\star$}
\textbf{题目} 若$y$为任一Mathieu解，求$\int(a-2q\cos2x)y(x)\dd x$。\par\textbf{思路与解答} 方程直接给被积式$=-y\prime\prime$，所以原函数$-y\prime+C$。$q=0,a=1,y=\cos x$时退化为$\int\cos x\dd x=\sin x+C$。\par\textbf{方法点评与边界} 特殊函数积分应写清采用的解与参数，不随意把不同归一化函数混同。
\end{example}

\section{积分理论、不等式与近似计算}\label{mat:integral_analysis}
\begin{examtip}把可积性、广义积分、含参极限及数值近似写成正文内容，连接计算技巧与严格条件。\end{examtip}
\subsection{上下和与可积准则}\label{mat:integral_analysis:riemann}
\begin{method}{上下和与可积准则}\textbf{适用条件} 闭区间上的有界函数。\par\textbf{实施与依据} 上和$U(P)$与下和$L(P)$之差可任意小等价于Riemann可积。\end{method}
\begin{example}{专题例题｜上下和与可积准则}{难度：$\star\star$}
\textbf{题目} 证明有界单调函数Riemann可积。\par\textbf{思路与解答} 等分网格宽$h$时，$U-L=h[f(b)-f(a)]$（递增情形），随$h\to0$趋0。递减同理。有限间断点也可用小总长度区间包围再控制误差。\par\textbf{方法点评与边界} 无界函数不能直接使用普通Riemann可积定义。
\end{example}
\subsection{原函数与可积性不是同一条件}\label{mat:integral_analysis:primitive}
\begin{method}{原函数与可积性不是同一条件}\textbf{适用条件} 区分连续、可积与存在原函数。\par\textbf{实施与依据} 用导数介值性排除跳跃函数的原函数。\end{method}
\begin{example}{专题例题｜原函数与可积性不是同一条件}{难度：$\star\star$}
\textbf{题目} 阶跃函数$f(x)=0\ (x<0),1\ (x\ge0)$在$[-1,1]$可积，是否有处处原函数？\par\textbf{思路与解答} 函数只有一个跳跃点，因此Riemann可积。但若存在$F\prime=f$，导数Darboux性质要求导数取得0与1之间的值，与阶跃取值矛盾。因此无处处原函数。\par\textbf{方法点评与边界} Newton--Leibniz不能只凭“可积”就随意指定原函数。
\end{example}
\subsection{Hölder与Minkowski}\label{mat:integral_analysis:holder}
\begin{method}{Hölder与Minkowski}\textbf{适用条件} $p,q>1,\frac{1}{p}+\frac{1}{q}=1$，相关幂可积。\par\textbf{实施与依据} Young不等式$ab\le \frac{a^p}{p}+\frac{b^q}{q}$经归一化积分给Hölder；再用Hölder证明三角不等式。\end{method}
\begin{example}{专题例题｜Hölder与Minkowski}{难度：$\star\star$}
\textbf{题目} 证明$\int_0^1|f|\le(\int_0^1|f|^p)^{\frac{1}{p}}$。\par\textbf{思路与解答} 令另一因子为1，Hölder右端第二个因子为$1$，结论成立。一般区间长度为$\ell$时额外乘$\ell^{\frac{1}{q}}$。\par\textbf{方法点评与边界} 把积分区间归一化后才能省去长度因子。
\end{example}
\subsection{Jensen、Chebyshev与单调权}\label{mat:integral_analysis:jensen}
\begin{method}{Jensen、Chebyshev与单调权}\textbf{适用条件} 凸函数或同向单调的函数对。\par\textbf{实施与依据} 凸函数用切线下界积分；同向单调用双重差乘积非负。\end{method}
\begin{example}{专题例题｜Jensen、Chebyshev与单调权}{难度：$\star\star$}
\textbf{题目} 若$f,g$在$[0,1]$同为递增，证明$\int fg\ge\int f\int g$。\par\textbf{思路与解答} $\int_0^1\int_0^1[f(x)-f(y)][g(x)-g(y)]\dd x\dd y\ge0$，展开为$2\int fg-2\int f\int g$，得到结论。\par\textbf{方法点评与边界} 反向单调时不等式反向。
\end{example}
\subsection{Gronwall积分不等式}\label{mat:integral_analysis:gronwall}
\begin{method}{Gronwall积分不等式}\textbf{适用条件} $u\ge0$连续，$u(t)\le A+B\int_0^tu(s)\dd s$，$A,B\ge0$。\par\textbf{实施与依据} 设$v=A+B\int_0^tu$，用$v\prime\le Bv$及积分因子。\end{method}
\begin{example}{专题例题｜Gronwall积分不等式}{难度：$\star\star$}
\textbf{题目} 证明$u(t)\le A\e^{Bt}$，$t\ge0$。\par\textbf{思路与解答} $v\prime=Bu\le Bv$，所以$(\e^{-Bt}v)\prime\le0$，得$v\le A\e^{Bt}$，再用$u\le v$。$A=0$时可用$A+\varepsilon$逼近，得$u\equiv0$。\par\textbf{方法点评与边界} 这是ODE唯一性与连续依赖的关键估计。
\end{example}
\subsection{反常积分的Dirichlet判别}\label{mat:integral_analysis:dirichlet}
\begin{method}{反常积分的Dirichlet判别}\textbf{适用条件} 原函数$F(x)=\int_a^xf$有界，$g$单调趋0。\par\textbf{实施与依据} 分部积分并利用$\int|g\prime|$的尾部控制；一般单调情形可用Stieltjes版本。\end{method}
\begin{example}{专题例题｜反常积分的Dirichlet判别}{难度：$\star\star$}
\textbf{题目} 证明$\int_1^\infty\frac{\sin x}{x}\dd x$收敛。\par\textbf{思路与解答} $\int_1^R\sin x\dd x$一致有界，$\frac{1}{x}\downarrow0$，Dirichlet判别成立；也可分部得到边界$-\frac{\cos R}{R}$加绝对收敛积分$-\int_1^R\frac{\cos x}{x^2}\dd x$。\par\textbf{方法点评与边界} 它不是绝对收敛：每个正弦半周期的绝对值积分与调和项同阶。
\end{example}
\subsection{Frullani积分}\label{mat:integral_analysis:frullani}
\begin{method}{Frullani积分}\textbf{适用条件} $a,b>0$且所用积分收敛，可由参数或截断控制。\par\textbf{实施与依据} 差分指数可写成参数积分，先证绝对或非负换序条件。\end{method}
\begin{example}{专题例题｜Frullani积分}{难度：$\star\star$}
\textbf{题目} 求$\int_0^\infty\frac{(\e^{-ax}-\e^{-bx})\dd x}{x}$。\par\textbf{思路与解答} 若$b>a$，分子$=\int_a^b x\e^{-tx}\dd t$，原积分为$\int_a^b\dd t\int_0^\infty\e^{-tx}\dd x=\int_a^b\frac{\dd t}{t}=\ln(\frac{b}{a})$。其余情形交换$a,b$。\par\textbf{方法点评与边界} 将两个各自发散的积分分开是不合法的，必须保留差。
\end{example}
\subsection{梯形与Simpson误差}\label{mat:integral_analysis:numerical}
\begin{method}{梯形与Simpson误差}\textbf{适用条件} 函数具有相应阶数连续导数。\par\textbf{实施与依据} 每小区间用低阶插值积分，累加余项。\end{method}
\begin{example}{专题例题｜梯形与Simpson误差}{难度：$\star\star$}
\textbf{题目} 写出等距复合梯形公式及误差阶。\par\textbf{思路与解答} $h=\frac{b-a}{n}$，$T_n=h[\frac{f(a)+f(b)}{2}+\sum_{k=1}^{n-1}f(a+kh)]$。若$|f\prime\prime|\le M$，线性插值余项积分后给$|I-T_n|\le\frac{(b-a)Mh^2}{12}$。复合Simpson在偶数分段且四阶导有界时误差为$O(h^4)$。\par\textbf{方法点评与边界} 更高阶规则需要更高阶正则性；奇点处不可盲用误差公式。
\end{example}
\subsection{含参积分的连续与求导}\label{mat:integral_analysis:parameter}
\begin{method}{含参积分的连续与求导}\textbf{适用条件} 参数位于紧区间，函数或偏导有可积支配。\par\textbf{实施与依据} 将连续、极限、求导分别匹配其支配条件。\end{method}
\begin{example}{专题例题｜含参积分的连续与求导}{难度：$\star\star$}
\textbf{题目} 证明$F(a)=\int_0^\infty\e^{-ax}\dd x=\frac{1}{a}$且$F\prime(a)=-\frac{1}{a^2}$，$a>0$。\par\textbf{思路与解答} 固定$a\in[\alpha,\beta]$，$\alpha>0$，被积式受$\e^{-\alpha x}$控制，参数导数受$x\e^{-\alpha x}$控制，两者可积。故可求导，得到$-\int_0^\infty x\e^{-ax}\dd x=-\frac{1}{a^2}$。\par\textbf{方法点评与边界} 当参数逼近0，统一支配失效，不能声称在含0区间一致。
\end{example}
\subsection{功、流体压力与消费者剩余}\label{mat:integral_analysis:applications}
\begin{method}{功、流体压力与消费者剩余}\textbf{适用条件} 连续力、密度或需求函数。\par\textbf{实施与依据} 先找微元物理量，再积分，记录单位与边界。\end{method}
\begin{example}{专题例题｜功、流体压力与消费者剩余}{难度：$\star\star$}
\textbf{题目} 弹簧力$F(x)=kx$，拉伸从0到$a$所作功。\par\textbf{思路与解答} 位移微元的功$\dd W=kx\dd x$，故$W=\frac{ka^2}{2}$。类似地静水压力为深度乘重度后对面积积分；需求曲线$p(q)$的消费者剩余为$\int_0^{q_0}[p(q)-p_0]\dd q$。\par\textbf{方法点评与边界} 经济和物理模型的假设与单位不能由纯积分公式代替。
\end{example}




\section{双元法、代数约化与复杂不定积分}\label{deep:integral_source_deep}

\begin{examtip}本节系统讨论二次双元、分式构型、单元化、高次双元、隐函数和伪椭圆积分。本节每次换元都说明微分关系，结果在原积分的连通定义区间理解。\end{examtip}

\subsection{双元基本恒等式的统一推导}\label{deep:integral_source_deep:duality}

\begin{method}{双元基本恒等式的统一推导}\textbf{适用条件} $u,v$可微且满足$v^2-u^2=a^2$或$u^2+v^2=a^2$。\par\textbf{步骤与依据} 先对约束求微分，再计算$\frac{u}{v}$、$u+v$或角变量的微分；选择能使分母次数降低的参数。\par\textbf{易错点} 代数关系与微分关系必须同时满足，不能只凭外形套公式。\end{method}

\begin{theorem}{双曲双元的基本微分}若$v^2-u^2=a^2>0,v>0$，则$\dd v=\frac{u\dd u}{v}$，$\dd\ln(u+v)=\frac{\dd u}{v}$，$\dd(\frac{u}{v})=\frac{a^2\dd u}{v^3}$。\end{theorem}\begin{proof}对$v^2-u^2=a^2$求微分得$v\dd v=u\dd u$。第二式由$\frac{\dd u+\dd v}{u+v}=\frac{\dd u}{v}$；第三式由$\frac{v\dd u-u\dd v}{v^2}=\frac{(v^2-u^2)\dd u}{v^3}$。这些关系分别对应对数型、降次型积分。\end{proof}

\begin{example}{考研风格｜双元与三角换元双解}{难度：$\star\star$}\label{deep:ex:integral_source_deep:dual1}\textbf{题目} 求$\int\frac{\dd x}{(x^2+a^2)^{\frac{3}{2}}}$，$a>0$。\par\textbf{思路} 直接识别$\frac{x}{\sqrt{x^2+a^2}}$的导数，或令$x=a\tan t$。\par\textbf{详细解答} \textbf{双元法：}取$u=x,v=\sqrt{x^2+a^2}$，上式给原函数$\frac{x}{a^2\sqrt{x^2+a^2}}+C$。\par\textbf{三角法：}$x=a\tan t$、$t\in(-\frac{\pi}{2},\frac{\pi}{2})$，积分为$a^{-2}\int\cos t\dd t=a^{-2}\sin t+C$，回代同上。\par\textbf{方法点评} 双元法避免显式反三角换元；三角法把分支选择呈现得更直观。\par\textbf{变式与易错点} 若$a=0$，需按$x$的正负重新积分$|x|^{-3}$。\end{example}

\begin{example}{竞赛风格｜分式构型与共轭因子}{难度：$\star\star\star$}\label{deep:ex:integral_source_deep:dual2}\textbf{题目} 求$\int\frac{\dd x}{\sqrt{x^2+1}(x+\sqrt{x^2+1})}$。\par\textbf{思路} 使用$(v+x)(v-x)=1$。\par\textbf{详细解答} 令$v=\sqrt{x^2+1}$，则$\frac{1}{x+v}=v-x$，被积函数$=1-\frac{x}{v}$，故答案$x-v+C$。另一法令$t=x+v>0$，由$\frac{\dd t}{t}=\frac{\dd x}{v}$，积分$\int t^{-2}\dd t=-\frac{1}{t}+C$；因$-\frac{1}{t}=x-v$，结果完全相同。\par\textbf{方法点评} 先共轭化只需一次初等积分；参数$t$则系统地消去整个根式族。\par\textbf{变式与易错点} 不要漏掉换元后分母中的另一个$t$。\end{example}

\begin{example}{竞赛风格｜有理化参数的一般公式}{难度：$\star\star\star$}\label{deep:ex:integral_source_deep:dual3}\textbf{题目} 证明含$x,\sqrt{x^2+a^2}$的任意有理式积分都能化为有理积分。\par\textbf{思路} 用双元和$t=x+\sqrt{x^2+a^2}$，反解两个变量。\par\textbf{详细解答} 令$v=\sqrt{x^2+a^2}$，则$v-x=\frac{a^2}{t}$，所以$x=\frac{t-\frac{a^2}{t}}{2},v=\frac{t+\frac{a^2}{t}}{2}$，$\dd x=\frac{(1+\frac{a^2}{t^2})\dd t}{2}$。因此$R(x,v)\dd x$变成$t$的有理式。例$\int\frac{\dd x}{v}$化为$\int\frac{\dd t}{t}=\ln t+C$。\par\textbf{方法点评} 这是双元化单元的结构证明，不局限于一个积分。\par\textbf{变式与易错点} $a>0$时$t>0$；平方差$v^2=x^2-a^2$须选择其可微实分支。\end{example}

\subsection{Euler三类替换与多重根式}\label{deep:integral_source_deep:euler}

\begin{method}{Euler三类替换与多重根式}\textbf{适用条件} 含$\sqrt{ax^2+bx+c}$的有理式，根式实且可微。\par\textbf{步骤与依据} $a>0$设根式$=\sqrt a\,x+t$；$c>0$可设根式$=tx+\sqrt c$；有实根$r$时设根式$=t(x-r)$。平方后解$x$。\par\textbf{易错点} 平方引入的支路需由原根式非负条件筛选。\end{method}

\begin{example}{考研风格｜Euler替换全过程}{难度：$\star\star\star$}\label{deep:ex:integral_source_deep:euler1}\textbf{题目} 求$\int\frac{\dd x}{x\sqrt{x^2+x+1}}$，$x\ne0$。\par\textbf{思路} 设$t=\frac{\sqrt{x^2+x+1}-1}{x}$，利用常数项是平方。\par\textbf{详细解答} 写$v=1+tx$，平方得$x^2+x+1=1+2tx+t^2x^2$，故$x=\frac{2t-1}{1-t^2}$。计算$\frac{\dd x}{xv}=\frac{2\dd t}{2t-1}$，所以原函数$\ln|2t-1|+C$。也可令$u=\frac{1}{x}$，在$x>0$积分变为$-\int\frac{\dd u}{\sqrt{u^2+u+1}}$，配方后为$-\ln|u+\frac{1}{2}+\sqrt{u^2+u+1}|+C$；两表达在每个正区间相差常数。\par\textbf{方法点评} Euler法统一适用；倒数换元在本题更短。\par\textbf{变式与易错点} $x<0$时$\sqrt{x^2+x+1}=|x|\sqrt{1+u+u^2}$，不能无条件把$|x|$写成$x$。\end{example}

\begin{example}{竞赛风格｜嵌套根式先降一层}{难度：$\star\star\star$}\label{deep:ex:integral_source_deep:nested}\textbf{题目} 求$\int\frac{\dd x}{\sqrt{x+\sqrt{x^2+1}}}$。\par\textbf{思路} 以双元和为变量，外层根式再平方根化。\par\textbf{详细解答} 设$t=\sqrt{x+\sqrt{x^2+1}}>0$。由$t^2=x+v$得$x=\frac{t^2-t^{-2}}{2}$，$\dd x=(t+t^{-3})\dd t$。积分$=\int(1+t^{-4})\dd t=t-\frac1{3t^3}+C$。另一法$x=\sinh u$使$x+v=\e^u$，积分$\int\e^{-\frac{u}{2}}\cosh u\dd u=\e^{\frac{u}{2}}-\frac{\e^{-\frac{3u}{2}}}{3}+C$，回代相同。\par\textbf{方法点评} 双元和把嵌套根式变成单一幂函数，双曲法解释其指数结构。\par\textbf{变式与易错点} 应使用$t>0$支路，不能给$t$任意符号。\end{example}

\begin{example}{竞赛风格｜根式比值参数}{难度：$\star\star\star$}\label{deep:ex:integral_source_deep:rootratio}\textbf{题目} 求$\int\sqrt{\frac{x-1}{x+1}}\dd x$，$x>1$。\par\textbf{思路} 用$t=\sqrt{\frac{x-1}{x+1}}$或$x=\cosh u$。\par\textbf{详细解答} 双曲法令$x=\cosh u,u>0$，根式$=\tanh(\frac{u}{2})$，乘$\dd x=\sinh u\dd u$得$(\cosh u-1)\dd u$，所以答案$\sqrt{x^2-1}-\operatorname{arcosh}x+C=\sqrt{x^2-1}-\ln(x+\sqrt{x^2-1})+C$。比值法$x=\frac{1+t^2}{1-t^2}$，积分为$\int\frac{4t^2}{(1-t^2)^2}\dd t$，部分分式给同值。\par\textbf{方法点评} 比值法对一般分式构型系统；双曲法最短。\par\textbf{变式与易错点} $x<-1$时根式和对数支路要另验，不能照抄本区间的推导。\end{example}

\subsection{有理积分、重因子与Hermite约化}\label{deep:integral_source_deep:rational}

\begin{method}{有理积分、重因子与Hermite约化}\textbf{适用条件} 被积函数是实系数有理式。\par\textbf{步骤与依据} 先长除，再因式分解；重因子用导数项降低幂数，剩余平方自由分母只产生对数和反正切。\par\textbf{易错点} 二次因子分子须设一次多项式，重因子各幂均须保留。\end{method}

\begin{example}{考研风格｜四次分母的实因式分解}{难度：$\star\star\star$}\label{deep:ex:integral_source_deep:quartic}\textbf{题目} 求$\int\frac{\dd x}{x^4+1}$。\par\textbf{思路} 分解为两个正定二次式，利用对称性配分子。\par\textbf{详细解答} $x^4+1=(x^2+\sqrt2x+1)(x^2-\sqrt2x+1)$，且$\frac1{x^4+1}=\frac{\frac{x}{2\sqrt2}+\frac{1}{2}}{x^2+\sqrt2x+1}+\frac{-\frac{x}{2\sqrt2}+\frac{1}{2}}{x^2-\sqrt2x+1}$。配成分母导数加常数后，原函数为$\frac1{4\sqrt2}\ln\frac{x^2+\sqrt2x+1}{x^2-\sqrt2x+1}+\frac1{2\sqrt2}[\arctan(\sqrt2x+1)+\arctan(\sqrt2x-1)]+C$。\par\textbf{方法点评} 实二次分解比复根法更适合直接书写实原函数；复根展开是第二条系统路线。\par\textbf{变式与易错点} 把两个反正切合并为一个时可能有分段常数，保留和式更稳妥。\end{example}

\begin{example}{竞赛风格｜Hermite约化的系数方程}{难度：$\star\star\star$}\label{deep:ex:integral_source_deep:hermite}\textbf{题目} 求$\int\frac{\dd x}{(x^2+1)^3}$。\par\textbf{思路} 设导数部分为$\frac{x}{(x^2+1)^2}$和$\frac{x}{x^2+1}$，余下反正切。\par\textbf{详细解答} 记$I_n=\int(1+x^2)^{-n}\dd x$。对$x(1+x^2)^{1-n}$求导得$[x(1+x^2)^{1-n}]'=(3-2n)(1+x^2)^{1-n}+2(n-1)(1+x^2)^{-n}$，因此$I_n=\frac{x}{2(n-1)(1+x^2)^{n-1}}+\frac{2n-3}{2(n-1)}I_{n-1}$。依次用$n=3,2$得$I_3=\frac{x}{4(1+x^2)^2}+\frac{3x}{8(1+x^2)}+\frac38\arctan x+C$。第二法$x=\tan t$后积分$\cos^4t$并降幂。\par\textbf{方法点评} 递推体现约化算法；三角法在低次时直观。\par\textbf{变式与易错点} 求导首项的符号和$2n-3$系数要逐项核验。\end{example}

\begin{example}{竞赛风格｜高次双元的可积信号}{难度：$\star\star\star$}\label{deep:ex:integral_source_deep:binomial}\textbf{题目} 求$\int x^2(1+x^3)^{-\frac{2}{3}}\dd x$，在$1+x^3>0$上。\par\textbf{思路} 观察内层导数；复杂分数幂不一定非初等。\par\textbf{详细解答} 令$u=1+x^3$，$\dd u=3x^2\dd x$，积分$=u^{\frac{1}{3}}+C=(1+x^3)^{\frac{1}{3}}+C$。另一法设$v=(1+x^3)^{\frac{1}{3}}$，直接微分$v^3-x^3=1$得$v^2\dd v=x^2\dd x$，所以原积分$\int\dd v=v+C$。\par\textbf{方法点评} 高次双元约束可以比二次根式更容易。\par\textbf{变式与易错点} 判断是否初等，不能仅看根式次数；还要看微分是否匹配。\end{example}

\subsection{伪椭圆积分与精确微分识别}\label{deep:integral_source_deep:pseudo}

\begin{method}{伪椭圆积分与精确微分识别}\textbf{适用条件} 出现三、四次根式，但分子含其导数或可构造有理函数的导数。\par\textbf{步骤与依据} 先试$u=P(x)$、$u=R(x)$；再检查$\sqrt{P(x)}$或有理函数与根式乘积的导数。\par\textbf{易错点} “三次或四次根式”并不自动意味着真正椭圆积分；反过来，也不能把任意此类积分都初等化。\end{method}

\begin{example}{竞赛风格｜外观椭圆而实际可换元}{难度：$\star\star\star$}\label{deep:ex:integral_source_deep:pseudo1}\textbf{题目} 求$\int\frac{(3x^2-1)\dd x}{\sqrt{x^3-x+2}}$，根式为正。\par\textbf{思路} 分子是根式内多项式的导数。\par\textbf{详细解答} 令$u=x^3-x+2$，积分$\int u^{-\frac{1}{2}}\dd u=2\sqrt{x^3-x+2}+C$。也可直接求导核验。与$\int\frac{\dd x}{\sqrt{x^3-x+2}}$相比，本题多出的导数因子改变了可积性质。\par\textbf{方法点评} 先检查精确微分可避免误引特殊函数。\par\textbf{变式与易错点} 不能在根式过零点处直接宣称整个实轴上的同一原函数可微。\end{example}

\begin{example}{竞赛风格｜有理映射诱导的高次根式}{难度：$\star\star\star$}\label{deep:ex:integral_source_deep:pseudo2}\textbf{题目} 求$\int \frac{(x^2-1)\dd x}{x\sqrt{x^4+3x^2+1}}$，$x>0$。\par\textbf{思路} 根式可提取$x^2$，用$u=x+\frac{1}{x}$。\par\textbf{详细解答} $x^4+3x^2+1=x^2[(x+\frac{1}{x})^2+1]$，$\dd u=(1-x^{-2})\dd x$，故原积分$=\int\frac{\dd u}{\sqrt{u^2+1}}=\ln(u+\sqrt{u^2+1})+C$。回代$u=x+\frac{1}{x}$即可。另一法先除分母中的$x^2$再找凑微分，得到同一转换。\par\textbf{方法点评} 这是“非齐次化齐次”的实质：让$x$与$x^{-1}$成对。\par\textbf{变式与易错点} $x<0$时提根式必须带$|x|$，符号需重新处理。\end{example}

\begin{example}{竞赛风格｜由导数反构造积分}{难度：$\star\star\star$}\label{deep:ex:integral_source_deep:exactalg}\textbf{题目} 判断$F=\frac{x\sqrt{x^4+x^2+1}}{x^2+1}$是否是$\frac{2x^4+3x^2+2}{(x^2+1)^2\sqrt{x^4+x^2+1}}$的原函数，并构造一个可初等积分的邻近题。\par\textbf{思路} 设$F=\frac{x\sqrt{x^4+x^2+1}}{x^2+1}$，直接检验其导数是否匹配。\par\textbf{详细解答} 令$P=x^4+x^2+1$，则$F'=\frac{(1-x^2)P+\frac{x(x^2+1)P'}{2}}{(x^2+1)^2\sqrt P}=\frac{x^6+3x^4+x^2+1}{(x^2+1)^2\sqrt P}$，并不匹配题中分子。因此该试探不能直接作答案。改取$G=\frac{x}{\sqrt P}$，有$G'=\frac{1-x^4}{P^{\frac{3}{2}}}$；这给出一个可初等化的邻近题$\int(1-x^4)P^{-\frac{3}{2}}\dd x=\frac{x}{\sqrt P}+C$。该邻近题的构造已完成；判错题不要求计算被否定候选对应的另一积分。\par\textbf{方法点评} 本例训练对“猜原函数”的否证与修正，展示复杂积分中求导验算的必要性。\par\textbf{变式与易错点} 不能把形式相近当作导数相同；本题的判错和变式构造均已完成。\end{example}

\subsection{三角混合、等分结构与指数锁}\label{deep:integral_source_deep:trig}

\begin{method}{三角混合、等分结构与指数锁}\textbf{适用条件} 三角有理式、指数三角混合或隐含导数因子。\par\textbf{步骤与依据} 先找半角有理化和复指数；若分母$Q$且分子为$Q'$，先对数；若两项形成闭合微分组，再设待定原函数。\par\textbf{易错点} “锁”只是对隐藏微分结构的简写，不能替代明确的可逆换元与验证。\end{method}

\begin{example}{考研风格｜一般线性三角分母}{难度：$\star\star\star$}\label{deep:ex:integral_source_deep:triglinear}\textbf{题目} 求$\int\frac{\dd x}{a+b\cos x}$，$a>|b|$。\par\textbf{思路} 半角代换后只有正定二次式。\par\textbf{详细解答} 令$t=\tan(\frac{x}{2})$，$\dd x=\frac{2\dd t}{1+t^2}$，分母$=\frac{(a+b)+(a-b)t^2}{1+t^2}$。所以答案$\frac2{\sqrt{a^2-b^2}}\arctan\!\left(\sqrt{\frac{a-b}{a+b}}\tan\frac x2\right)+C$，在不跨越半角极点的区间理解。第二法将积分区间分段后用复单位圆，可得到相应定积分$\frac{2\pi}{\sqrt{a^2-b^2}}$。\par\textbf{方法点评} 半角法直接产出原函数；围道法更适合整周期积分。\par\textbf{变式与易错点} $a^2<b^2$时二次式会有实根，原函数改为对数型。\end{example}

\begin{example}{竞赛风格｜三角混合与分部闭合}{难度：$\star\star$}\label{deep:ex:integral_source_deep:trigmix}\textbf{题目} 求$\int x\sec^2x\dd x$及$\int \e^x(\sin x+\cos x)\dd x$。\par\textbf{思路} 前者让$x$求导，后者识别积的导数。\par\textbf{详细解答} 第一题取$u=x,\dd v=\sec^2x\dd x$，得$x\tan x-\int\tan x\dd x=x\tan x+\ln|\cos x|+C$。第二题由$(\e^x\sin x)'=\e^x(\sin x+\cos x)$直接得$\e^x\sin x+C$；也可设$\e^x(A\sin x+B\cos x)$比较系数得$A=1,B=0$。\par\textbf{方法点评} 分部法与待定系数法都依赖求导后复杂度是否下降。\par\textbf{变式与易错点} 复指数法取实部或虚部时，须保留最后的实常数。\end{example}

\begin{example}{竞赛风格｜隐函数参数积分}{难度：$\star\star\star$}\label{deep:ex:integral_source_deep:implicit}\textbf{题目} 曲线$y^2=x^2+1,y>0$上，求$\int y\dd x$。\par\textbf{思路} 用约束微分或参数化，最后回到$x$。\par\textbf{详细解答} 分部得$\int y\dd x=xy-\int x\dd y=xy-\int \frac{x^2\dd x}{y}=xy-\int(y-\frac{1}{y})\dd x$。移项得$\frac12[xy+\ln(x+y)]+C$。第二法$x=\sinh t,y=\cosh t$，积分$\int\cosh^2t\dd t=\frac14\sinh2t+\frac{t}{2}+C$，回代相同。\par\textbf{方法点评} 双元分部自动把同一积分移回左侧，适合圆锥曲线约束。\par\textbf{变式与易错点} 隐函数表达必须声明所取分支。\end{example}

\begin{example}{竞赛风格｜由一阶方程生成原函数}{难度：$\star\star\star$}\label{deep:ex:integral_source_deep:odeint}\textbf{题目} 求$\int \e^x(x^3-3x^2+6x-6)\dd x$。\par\textbf{思路} 设原函数$\e^xP(x)$，将积分变成$P'+P=Q$。\par\textbf{详细解答} 比较次数，设$P=x^3+ax^2+bx+c$，则$P'+P=x^3+(a+3)x^2+(b+2a)x+(c+b)$。与$Q$比较得$a=-6,b=18,c=-24$，所以答案$\e^x(x^3-6x^2+18x-24)+C$。另一法连续分部，有限算子逆$(1+D)^{-1}Q=Q-Q'+Q''-Q'''$，也得该多项式。\par\textbf{方法点评} 连接微分方程思想与微分算子法。\par\textbf{变式与易错点} 输入多项式若变为一般函数，几何级数形式不再自动有限终止。\end{example}

\subsection{共同根次、指数单元与根分解通式}\label{deep:integral_source_deep:unitroot}

\begin{method}{共同根次、指数单元与根分解通式}\textbf{适用条件} 有理分数次幂可取公分母，指数有理式可取$\e^x$。\par\textbf{步骤与依据} 将所有变量和微分同时化为一个有理参数；一般有理式的简单根可用导数给部分分式系数。\par\textbf{易错点} 共同根次换元通常需限定$x>0$；复对数结果要成对还原实值。\end{method}

\begin{example}{考研风格｜共同根次化有理}{难度：$\star\star$}\label{deep:ex:integral_source_deep:lcmroot}\textbf{题目} 求$\int\frac{\dd x}{\sqrt x+\sqrt[3]x}$，$x>0$。\par\textbf{思路} 平方根与立方根的公分母是6。\par\textbf{详细解答} 令$x=t^6,t>0$，积分$=6\int \frac{t^3}{t+1}\dd t=6\int(t^2-t+1-\frac{1}{t+1})\dd t$，得$2t^3-3t^2+6t-6\ln(t+1)+C$。回代为$2\sqrt x-3x^{\frac{1}{3}}+6x^{\frac{1}{6}}-6\ln(1+x^{\frac{1}{6}})+C$。\par\textbf{方法点评} 共同根次换元把所有分数幂同时化为整数幂。\par\textbf{变式与易错点} 若只换根式而不换$\dd x$，次数会完全错误。\end{example}

\begin{example}{考研风格｜指数单元与部分分式}{难度：$\star\star$}\label{deep:ex:integral_source_deep:expunit}\textbf{题目} 求$\int\frac{\e^{2x}\dd x}{\e^{2x}+3\e^x+2}$。\par\textbf{思路} 令$t=\e^x$，分母可分解。\par\textbf{详细解答} 积分$=\int \frac{t\dd t}{(t+1)(t+2)}=\int[-\frac{1}{t+1}+\frac{2}{t+2}]\dd t$，答案$-\ln(1+\e^x)+2\ln(2+\e^x)+C$。可直接设原函数$A\ln(1+\e^x)+B\ln(2+\e^x)$，求导比较得$A=-1,B=2$，作为第二法。\par\textbf{方法点评} 指数类的复杂外观被一个正参数消去。\par\textbf{变式与易错点} 此处对数真数恒正，不必多加绝对值。\end{example}

\begin{theorem}{简单根有理积分的统一分解}若$P$次数大于$Q$且$P$在复数域有互异根$r_1,\ldots,r_m$，则$\frac{Q(x)}{P(x)}=\sum_{j=1}^m\frac{\frac{Q(r_j)}{P'(r_j)}}{x-r_j}$。\end{theorem}\begin{proof}每项与$\frac{Q}{P}$在$r_j$有相同留数，差函数无极点且在无穷处趋0，故为零。也可通分得到次数小于$m$的多项式，在$m$个互异根处为0，因此恒为0。积分后为复对数和；实系数时将共轭根项配对，得到实对数和反正切。重根先作Hermite约化。\end{proof}

\begin{example}{考研风格｜三次有理积分双解}{难度：$\star\star\star$}\label{deep:ex:integral_source_deep:cubic}\textbf{题目} 求$\int\frac{\dd x}{1+x^3}$，$x\ne-1$。\par\textbf{思路} 按实因式分解，或用三个单位根的通式。\par\textbf{详细解答} $\frac{1}{1+x^3}=\frac{1}{3(x+1)}+\frac{-x+2}{3(x^2-x+1)}$。把二次分子配成导数，得$\frac13\ln|x+1|-\frac16\ln(x^2-x+1)+\frac1{\sqrt3}\arctan(\frac{2x-1}{\sqrt3})+C$。根分解通式给三个复对数，配共轭项可还原同一实式。\par\textbf{方法点评} 统一有理积分可由因式分解与部分分式解释，所用函数及其定义域均须明确。\par\textbf{变式与易错点} 对数的系数和反正切的缩放同时由配方决定。\end{example}

\subsection{恒等式诱导、Pell型关系与算法边界}\label{deep:integral_source_deep:pell}

\begin{method}{恒等式诱导、Pell型关系与算法边界}\textbf{适用条件} 已找到多项式$P,Q,D$使$P^2-DQ^2=c\ne0$，且$Q\sqrt D\ne0$。\par\textbf{步骤与依据} 对代数恒等式求导，构造$\ln|P+Q\sqrt D|$；连分式算法是寻找这类恒等式的一种途径。\par\textbf{易错点} 算法发现候选后必须检验恒等式；不存在候选不等于已经证明不可初等积分。\end{method}

\begin{theorem}{Pell型恒等式的对数微分}设$P,Q,D$为实多项式，在所用开区间有$D>0,Q\ne0$。若$P^2-DQ^2=c\ne0$，则$\frac{\dd\ln|P+Q\sqrt D|}{\dd x}=\frac{P'}{Q\sqrt D}$，在实表达有定义的区间成立。\end{theorem}\begin{proof}微分恒等式得$PP'=DQQ'+\frac{D'Q^2}{2}$。计算对数导数的分子$P'+Q'\sqrt D+\frac{QD'}{2\sqrt D}$，将后两项写成$\frac{DQ'+\frac{QD'}{2}}{\sqrt D}=\frac{PP'}{Q\sqrt D}$，整体等于$\frac{P'(P+Q\sqrt D)}{Q\sqrt D}$，约去分母即得。\end{proof}

\begin{example}{竞赛风格｜高次恒等式核验}{难度：$\star\star\star$}\label{deep:ex:integral_source_deep:pellcase}\textbf{题目} 求$\int\frac{2x\dd x}{\sqrt{x^4+x^2+1}}$。\par\textbf{思路} 可用$u=x^2$后配方，也可用Pell型关系。\par\textbf{详细解答} 取$P=2x^2+1,Q=2,D=x^4+x^2+1$，有$P^2-DQ^2=-3$，所以原函数$\ln|2x^2+1+2\sqrt{x^4+x^2+1}|+C$。第二法$u=x^2$，积分$\int\frac{\dd u}{\sqrt{u^2+u+1}}=\ln(u+\frac{1}{2}+\sqrt{u^2+u+1})+C$，两式只差$\ln2$。\par\textbf{方法点评} 展示“恒等式诱导”背后的严格导数机制。\par\textbf{变式与易错点} 本例有$2x$因子；去掉它后不能凭同一恒等式得相同答案。\end{example}

\section{特殊函数积分与含参方法的细化}\label{deep:integral_special_deep}

\begin{examtip}定义与归一化在本册内固定，公式由微分、级数或微分方程推导。\end{examtip}

\subsection{含参积分求导的可检验条件}\label{deep:integral_special_deep:parameter}

\begin{method}{含参积分求导的可检验条件}\textbf{适用条件} 有限区间连续可微，或无穷区间有与参数无关的可积控制。\par\textbf{步骤与依据} 先给参数范围，再估计$\frac{\partial f}{\partial a}$；计算参数导数，最后用初值确定常数。\par\textbf{易错点} 仅被积函数逐点可微不足以换序。\end{method}

\begin{theorem}{含参积分求导的控制条件}设$f(x,a)$对$a$可微，在$a_0$邻域有$|\partial_af(x,a)|\le g(x)$，$g$可积，且$f(\cdot,a_0)$可积。则$F(a)=\int f(x,a)\dd x$可微，$F'(a)=\int\partial_af(x,a)\dd x$。\end{theorem}\begin{proof}差商由一元中值定理受$g(x)$控制，并逐点趋$\partial_af(x,a)$，控制收敛给差商积分的极限。有限区间连续偏导情形可取$g$为常数。\end{proof}

\begin{example}{竞赛风格｜Feynman参数法双解}{难度：$\star\star\star$}\label{deep:ex:integral_special_deep:arctanparam}\textbf{题目} 求$I(a)=\int_0^\infty\frac{\e^{-x}\sin(ax)\dd x}{x}$，$a\in\R$。\par\textbf{思路} 参数求导后变成指数三角积分。\par\textbf{详细解答} $|\partial_a(\frac{\e^{-x}\sin ax}{x})|=|\e^{-x}\cos ax|\le\e^{-x}$，可换序。故$I'(a)=\int_0^\infty\e^{-x}\cos ax\dd x=\frac{1}{1+a^2}$，$I(0)=0$，得$I(a)=\arctan a$。第二法$\frac{\sin(ax)}{x}=\int_0^a\cos(tx)\dd t$，用绝对可积换序，积分$\int_0^a\frac{\dd t}{1+t^2}$。\par\textbf{方法点评} 两法共同依赖可积控制，但第二法不需显式微分。\par\textbf{变式与易错点} 去掉$\e^{-x}$后必须按反常积分另证收敛，不能照搬同一控制函数。\end{example}

\begin{example}{竞赛风格｜Frullani公式的严格版本}{难度：$\star\star\star$}\label{deep:ex:integral_special_deep:frullani}\textbf{题目} 设$a,b>0$，求$\int_0^\infty\frac{(\e^{-ax}-\e^{-bx})\dd x}{x}$。\par\textbf{思路} 把分子写成对参数的积分，再换序。\par\textbf{详细解答} 假设$b>a$，$\e^{-ax}-\e^{-bx}=\int_a^b x\e^{-tx}\dd t$。被积函数非负，Tonelli换序得$\int_a^b\int_0^\infty\e^{-tx}\dd x\dd t=\int_a^b\frac{\dd t}{t}=\ln(\frac{b}{a})$。$b<a$由交换符号得到同式。另一法固定$a$对$b$求导得$\frac{1}{b}$，再用$b=a$为零。\par\textbf{方法点评} 差的整体可积；两个单独的$x^{-1}\e^{-ax}$在0都发散。\par\textbf{变式与易错点} 一般Frullani公式还需端点值和适当收敛条件，不能任意套用。\end{example}

\begin{example}{竞赛风格｜对数正弦积分双解}{难度：$\star\star\star$}\label{deep:ex:integral_special_deep:logsin}\textbf{题目} 求$\int_0^{\frac{\pi}{2}}\ln\sin x\dd x$。\par\textbf{思路} 对称性与倍角形成同一积分的线性方程。\par\textbf{详细解答} 记$I$，换$x\mapsto\frac{\pi}{2}-x$得$I=\int\ln\cos x$。相加$2I=\int\ln(\sin2x)\dd x-(\frac{\pi}{2})\ln2$。令$u=2x$并利用关于$\frac{\pi}{2}$的对称性，第一项为$I$，故$I=-(\frac{\pi}{2})\ln2$。第二法从$J(a)=\int_0^{\frac{\pi}{2}}\sin^ax\dd x=\frac{\sqrt\pi\Gamma(\frac{a+1}{2})}{2\Gamma(\frac{a+2}{2})}$在$a=0$求导，再用Gamma倍增公式的对数导数，也得同值。\par\textbf{方法点评} 第一法初等且短；Gamma法可系统给高阶对数矩。\par\textbf{变式与易错点} 端点$\ln x$可积，必须先确认反常积分存在。\end{example}

\subsection{多重对数、对数乘积与分支}\label{deep:integral_special_deep:polylog}

\begin{method}{多重对数、对数乘积与分支}\textbf{适用条件} $|x|<1$内幂级数可逐项求导积分，实延拓另记区间。\par\textbf{步骤与依据} 将$\ln(1-x)$与$\frac{1}{x}$配对；用反射恒等式转换自变量。\par\textbf{易错点} 同名特殊函数的归一化与复分支可能不同。\end{method}

\begin{example}{竞赛风格｜二重对数反射式}{难度：$\star\star\star$}\label{deep:ex:integral_special_deep:dilog}\textbf{题目} 证明$0<x<1$时$\operatorname{Li}_2(x)+\operatorname{Li}_2(1-x)=\frac{\pi^2}{6}-\ln x\ln(1-x)$。\par\textbf{思路} 两边求导后检查为常数，再由端点确定常数。\par\textbf{详细解答} 由$\operatorname{Li}_2'(x)=-\frac{\ln(1-x)}{x}$，左边导数$=-\frac{\ln(1-x)}{x}+\frac{\ln x}{1-x}$，恰等于$-\ln x\ln(1-x)$的导数。故和加$\ln x\ln(1-x)$为常数。令$x\downarrow0$，$\operatorname{Li}_2(0)=0,\operatorname{Li}_2(1)=\sum n^{-2}=\frac{\pi^2}{6}$，乘积对数趋0，结论成立。\par\textbf{方法点评} 同一方法可验证许多指对双元恒等式；常数由一个方便点确定。\par\textbf{变式与易错点} 把恒等式直接用于$x>1$时须处理复对数分支。\end{example}

\begin{example}{竞赛风格｜对数乘积的完整原函数}{难度：$\star\star\star$}\label{deep:ex:integral_special_deep:logproduct}\textbf{题目} 求$\int\ln x\ln(1-x)\dd x$，$0<x<1$。\par\textbf{思路} 分部一次，把剩余项拆成初等项和二重对数。\par\textbf{详细解答} 取$u=\ln x\ln(1-x),\dd v=\dd x$，得到$x\ln x\ln(1-x)-\int\ln(1-x)\dd x+\int \frac{x\ln x}{1-x}\dd x$。利用$\frac{x}{1-x}=-1+\frac{1}{1-x}$，以及$\int\frac{\ln x}{1-x}\dd x=-\ln x\ln(1-x)-\operatorname{Li}_2(x)$，整理为$(x-1)\ln x\ln(1-x)+(1-x)\ln(1-x)-x\ln x+2x-1-\operatorname{Li}_2(x)+C$。\par\textbf{方法点评} 每个非初等项都来自已定义的导数关系，可以逐项验算。\par\textbf{变式与易错点} 积分常数可吸收末尾$-1$；不能把$\operatorname{Li}_2$漏掉。\end{example}

\subsection{椭圆积分的标准形与退化}\label{deep:integral_special_deep:elliptic}

\begin{method}{椭圆积分的标准形与退化}\textbf{适用条件} 平方根下有不能消去的四次结构。\par\textbf{步骤与依据} 用$x=\sin\theta$化为Legendre标准形；识别参数$k=0$或$1$的初等退化。\par\textbf{易错点} 写出特殊函数名称后仍需给定义、参数和回代。\end{method}

\begin{definition}{Legendre三类不完全椭圆积分}$F(\phi,k)=\int_0^\phi(1-k^2\sin^2t)^{-\frac{1}{2}}\dd t$，$E(\phi,k)=\int_0^\phi(1-k^2\sin^2t)^{\frac{1}{2}}\dd t$，$\Pi(n;\phi,k)=\int_0^\phi[(1-n\sin^2t)\sqrt{1-k^2\sin^2t}]^{-1}\dd t$。先在$|k|<1$、实$\phi$及分母不为0的区间定义。\end{definition}

\begin{example}{竞赛风格｜三类椭圆标准积分}{难度：$\star\star\star$}\label{deep:ex:integral_special_deep:elliptic3}\textbf{题目} 把$\int\frac{\dd x}{\sqrt{(1-x^2)(1-k^2x^2)}}$、$\int\frac{\sqrt{1-k^2x^2}\dd x}{\sqrt{1-x^2}}$和带额外因子$(1-nx^2)^{-1}$的第一式化为标准函数，$|x|<1$。\par\textbf{思路} 共同代换$x=\sin\phi$。\par\textbf{详细解答} 在$\phi\in(-\frac{\pi}{2},\frac{\pi}{2})$，$\dd x=\cos\phi\dd\phi,\sqrt{1-x^2}=\cos\phi$，三式分别是$F(\arcsin x,k)+C$、$E(\arcsin x,k)+C$、$\Pi(n;\arcsin x,k)+C$。$k=0$时前两式退化为$\arcsin x$，第三式退化为三角有理式；$k=1$时第一式在$|x|<1$变为$\operatorname{artanh}x$。\par\textbf{方法点评} 把三元根式转成同一角变量，可统一三类椭圆积分的标准形式。\par\textbf{变式与易错点} 第三类须避开$1-nx^2=0$；完全积分对应$\phi=\frac{\pi}{2}$，要另检端点。\end{example}

\subsection{Bessel、修正Bessel与乘积积分}\label{deep:integral_special_deep:bessel}

\begin{method}{Bessel、修正Bessel与乘积积分}\textbf{适用条件} 函数满足相应二阶方程，参数使表达有定义。\par\textbf{步骤与依据} 利用递推或两个方程相减得到Lagrange恒等式。\par\textbf{易错点} 修正Bessel方程的势项符号不同，公式不能机械照搬。\end{method}

\begin{theorem}{Bessel乘积的Lagrange恒等式}设$u(x)=J_\nu(ax),v(x)=J_\nu(bx)$，则$[x(u'v-uv')]'=(b^2-a^2)xuv$。当$a^2\ne b^2$，据此可积出$\int xJ_\nu(ax)J_\nu(bx)\dd x$。\end{theorem}\begin{proof}两方程为$(xu')'+(a^2x-\frac{\nu^2}{x})u=0$与$v$的对应式。前式乘$v$，后式乘$u$相减，$\frac{\nu^2}{x}$项抵消，得导数恒等式。\end{proof}

\begin{example}{竞赛拓展｜双Bessel积分与正交}{难度：$\star\star\star$}\label{deep:ex:integral_special_deep:besselpair}\textbf{题目} 推导上式积分，并说明不同零点参数时在$[0,1]$的正交性。\par\textbf{思路} 积分边界项决定是否为零。\par\textbf{详细解答} $\int xJ_\nu(ax)J_\nu(bx)\dd x=\frac{x[u'v-uv']}{b^2-a^2}+C$。若$\nu\ge0$、$a,b$是$J_\nu$不同的正零点，则$x=1$处$u=v=0$；$x\downarrow0$时由$J_\nu(cx)=O(x^\nu)$及级数首项抵消，边界项趋0，故$\int_0^1xJ_\nu(ax)J_\nu(bx)\dd x=0$。\par\textbf{方法点评} 不用背乘积积分表，方程相减即可推导。\par\textbf{变式与易错点} $a=b$不能把分母0直接代入，应取参数极限或另用平方积分公式。\end{example}

\begin{example}{竞赛拓展｜修正Bessel与Airy积分}{难度：$\star\star\star$}\label{deep:ex:integral_special_deep:modified}\textbf{题目} 证明$\int xI_0(x)\dd x=xI_1(x)+C$，并求$\int x\operatorname{Ai}(x)^2\dd x$。\par\textbf{思路} 前者用递推，后者由Airy方程构造能量差。\par\textbf{详细解答} 级数$I_0=\sum\frac{(\frac{x}{2})^{2k}}{(k!)^2}$和$I_1=\sum\frac{(\frac{x}{2})^{2k+1}}{k!(k+1)!}$逐项微分给$(xI_1)'=xI_0$。对$y=\operatorname{Ai}$，$y''=xy$，故$(xy^2-y'^2)'=y^2$，先得到$\int y^2=xy^2-y'^2+C$；而$\int xy^2=\int yy''=yy'-\int y'^2$。为直接闭合，计算$[x^2y^2-xy'^2+yy']'=3xy^2$，所以所求为$\frac{x^2y^2-xy'^2+yy'}{3}+C$。\par\textbf{方法点评} 由方程构造组合导数比单独分部更系统。\par\textbf{变式与易错点} 最后组合的每项都须求导，不能将$\int y^2$误当$\int xy^2$。\end{example}

\subsection{Legendre、超几何与Mathieu积分}\label{deep:integral_special_deep:orthogonal}

\begin{method}{Legendre、超几何与Mathieu积分}\textbf{适用条件} 特殊函数由级数或自伴微分方程定义。\par\textbf{步骤与依据} 先确定归一化，再由自伴形式导出正交；超几何积分可从逐项积分识别。\par\textbf{易错点} 参数奇点、端点条件和正交权函数均属于结论条件。\end{method}

\begin{example}{竞赛拓展｜Legendre正交的证明}{难度：$\star\star\star$}\label{deep:ex:integral_special_deep:legendreproof}\textbf{题目} 用Rodrigues公式证明$m<n$时$\int_{-1}^1x^mP_n(x)\dd x=0$。\par\textbf{思路} 分部$n$次，边界项由$(x^2-1)^n$的高阶零点消失。\par\textbf{详细解答} $P_n=\frac{(2^nn!)^{-1}\dd^n(x^2-1)^n}{\dd x^n}$。将$n$阶导数逐次移到$x^m$上，每一步边界只涉及$(x^2-1)^n$的低于$n$阶导数，在$\pm1$皆0。最终$x^m$的$n$阶导数为0，积分为0。因$P_m$次数$m$，得到$\int P_mP_n=0$。归一化平方积分为$\frac{2}{2n+1}$，可用首项系数和同一分部法算出。\par\textbf{方法点评} 正交来自高阶零点和降阶，不只是奇偶性。\par\textbf{变式与易错点} 只有$m<n$时最后导数为0，不能由此断言平方积分也是0。\end{example}

\begin{example}{竞赛拓展｜二项式积分的超几何式}{难度：$\star\star\star$}\label{deep:ex:integral_special_deep:hypergeneral}\textbf{题目} 求$\int x^{\mu-1}(1-zx^\nu)^{-\alpha}\dd x$的级数表达，$x>0,\mu,\nu>0,|zx^\nu|<1$。\par\textbf{思路} 二项式展开，再把$\mu+\nu k$写成Pochhammer比值。\par\textbf{详细解答} 展开后积分$\sum_{k\ge0}\frac{(\alpha)_kz^kx^{\mu+\nu k}}{k!(\mu+\nu k)}$。用$\frac{(\frac{\mu}{\nu})_k}{(1+\frac{\mu}{\nu})_k}=\frac{\mu}{\mu+\nu k}$，得到$\frac{x^\mu}{\mu}{}_2F_1(\alpha,\frac{\mu}{\nu};1+\frac{\mu}{\nu};zx^\nu)+C$。一致收敛保证紧子区间内逐项积分。\par\textbf{方法点评} 许多特殊函数积分可由此通式特化，但每次须检查参数。\par\textbf{变式与易错点} 超出收敛圆需解析延拓；不能把级数证明无条件扩大到全实轴。\end{example}

\begin{example}{竞赛拓展｜Mathieu函数的积分恒等式}{难度：$\star\star\star$}\label{deep:ex:integral_special_deep:mathieu}\textbf{题目} 设$y_a,y_b$分别满足$y''+(a-2q\cos2x)y=0$与参数$b$的方程，$a\ne b$。求$\int y_ay_b\dd x$。\par\textbf{思路} 两式交叉相减，公共势函数抵消。\par\textbf{详细解答} $[y_a'y_b-y_ay_b']'=(b-a)y_ay_b$，所以原函数$\frac{y_a'y_b-y_ay_b'}{b-a}+C$。若两函数满足同一自伴周期或分离边界条件使端点Wronskian相同，则区间积分为0，得到不同特征值的正交性。\par\textbf{方法点评} 此Wronskian恒等式给出Mathieu函数乘积积分的统一构造。\par\textbf{变式与易错点} $a=b$时须另作极限；边界条件不同不保证正交。\end{example}

\subsection{不完全Beta、Fresnel与非初等边界}\label{deep:integral_special_deep:betaextra}

\begin{method}{不完全Beta、Fresnel与非初等边界}\textbf{适用条件} 保留参数取值和积分路径，先定义再求导核验。\par\textbf{步骤与依据} 有限区间幂权可化不完全Beta；二次相位用Fresnel；对数核用多重对数。\par\textbf{易错点} 命名特殊函数并不取消定义域和分支检查。\end{method}

\begin{example}{竞赛拓展｜无穷区间幂权的统一代换}{难度：$\star\star\star$}\label{deep:ex:integral_special_deep:incompletebeta}\textbf{题目} 把$\int x^{p-1}(1+x)^{-q}\dd x$化为不完全Beta形式，$x>0,p>0$。\par\textbf{思路} 令$t=\frac{x}{1+x}$，把半轴映成$(0,1)$。\par\textbf{详细解答} $x=\frac{t}{1-t}$、$\dd x=\frac{\dd t}{(1-t)^2}$，所以被积式$=t^{p-1}(1-t)^{q-p-1}\dd t$。定义$B_z(a,b)=\int_0^zt^{a-1}(1-t)^{b-1}\dd t$，得到$B_{\frac{x}{1+x}}(p,q-p)+C$；对$z<1$，第二参数可为负而不遇端点1。若求$0$到$\infty$的定积分，则需要$q>p>0$，结果$B(p,q-p)$。\par\textbf{方法点评} 同一代换同时给原函数和收敛条件。\par\textbf{变式与易错点} 不完全积分的参数条件比完全Beta积分宽，不能混写。\end{example}

\begin{example}{竞赛风格｜二重对数的另一类原函数}{难度：$\star\star\star$}\label{deep:ex:integral_special_deep:logplus}\textbf{题目} 求$\int\frac{\ln x\dd x}{1+x}$，$x>0$。\par\textbf{思路} 分部使$\frac{\ln(1+x)}{x}$出现。\par\textbf{详细解答} 取$u=\ln x,\dd v=\frac{\dd x}{1+x}$，得$\ln x\ln(1+x)-\int\frac{\ln(1+x)\dd x}{x}$。由$\frac{\dd}{\dd x}\operatorname{Li}_2(-x)=-\frac{\ln(1+x)}{x}$，结果$\ln x\ln(1+x)+\operatorname{Li}_2(-x)+C$。在$0<x<1$也可展开$\frac{1}{1+x}$逐项积分，并整理成同式；到$x>1$用实积分定义延拓。\par\textbf{方法点评} 符号由链式法则决定，不要仅靠背$\operatorname{Li}_2'$。\par\textbf{变式与易错点} 采用实轴负参数的连续实分支，不额外引入虚常数。\end{example}

\begin{example}{竞赛拓展｜Fresnel的尺度与乘积}{难度：$\star\star\star$}\label{deep:ex:integral_special_deep:fresnelprimitive}\textbf{题目} 求$\int\cos(ax^2)\dd x$，$a>0$，并求$\int x\cos(ax^2)\dd x$。\par\textbf{思路} 第一式是Fresnel定义，第二式有内层导数而初等。\par\textbf{详细解答} 固定$C(z)=\int_0^z\cos(\frac{\pi t^2}{2})\dd t$。令$t=\sqrt{\frac{2a}{\pi}}\,x$，第一式为$\sqrt{\frac{\pi}{2a}}\,C(\sqrt{\frac{2a}{\pi}}\,x)+C_0$。第二式令$u=ax^2$，得到$\frac{\sin(ax^2)}{2a}+C_0$。相同相位下，分子是否含$x$决定积分的初等性。\par\textbf{方法点评} 区别特殊函数定义和可用简单换元的邻近题。\par\textbf{变式与易错点} $a<0$时余弦偶性可用$|a|$；正弦情形还需保留符号。\end{example}

\section{复变积分与变换的可计算桥梁}\label{deep:integral_complex_deep}

\begin{examtip}本节先建立公式，再展示实积分、卷积与微分方程的联系。\end{examtip}

\subsection{解析性、Cauchy公式与留数}\label{deep:integral_complex_deep:complex}

\begin{method}{解析性、Cauchy公式与留数}\textbf{适用条件} 函数在围道及内部解析，或只有有限个孤立极点。\par\textbf{步骤与依据} 辨认区域、方向和奇点阶数；用局部Laurent系数计算留数；闭合实轴围道时估计大弧。\par\textbf{易错点} 解析性不能由连续性替代；路径穿过奇点时原定理不适用。\end{method}

\begin{definition}{复导数与Laurent系数}复导数要求商$\frac{f(z+h)-f(z)}{h}$在复方向$h\to0$有同一极限。解析函数局部有Taylor展开；环域解析函数有Laurent展开$\sum_{n\in\mathbb Z}a_n(z-z_0)^n$，$a_{-1}$称留数。若$f=\frac{g}{(z-z_0)^m}$且$g$解析，则留数$\frac{g^{(m-1)}(z_0)}{(m-1)!}$。\end{definition}

\begin{theorem}{Cauchy积分公式及导数式}正向简单闭曲线$\Gamma$及内部含于解析区域，$z_0$在内部，则$f^{(n)}(z_0)=\frac{n!}{2\pi i}\int_\Gamma \frac{f(z)}{(z-z_0)^{n+1}}\dd z$。\end{theorem}\begin{proof}先对$n=0$，挖去$z_0$的小圆，Cauchy定理把外边积分化为小圆积分。写$f(z)=f(z_0)+[f(z)-f(z_0)]$，常数部分积分$2\pi if(z_0)$，余项由连续性趋0。对$z_0$在紧子区域内逐次求导，分母与围道有正距离，故可在积分号内求导，得到$n$阶式。\end{proof}

\begin{example}{竞赛拓展｜高阶极点与Cauchy公式}{难度：$\star\star\star$}\label{deep:ex:integral_complex_deep:cauchy}\textbf{题目} 计算正向圆$|z|=2$上的$\int\frac{\e^z}{(z-1)^3}\dd z$。\par\textbf{思路} 分母三次对应二阶导数。\par\textbf{详细解答} 取$f(z)=\e^z,z_0=1,n=2$，积分$=\frac{2\pi i f''(1)}{2!}=\pi i\e$。第二法在$z=1$展开$\e^z=\e[1+(z-1)+\frac{(z-1)^2}{2}+\cdots]$，留数为$\frac{\e}{2}$，同得$\pi i\e$。\par\textbf{方法点评} 导数公式最快，Laurent法有利于多个相乘因子。\par\textbf{变式与易错点} 分母幂$m$对应$m-1$阶导数。\end{example}

\begin{example}{竞赛拓展｜四次有理积分的留数法}{难度：$\star\star\star$}\label{deep:ex:integral_complex_deep:residue4}\textbf{题目} 计算$\int_{-\infty}^{\infty}\frac{\dd x}{1+x^4}$。\par\textbf{思路} 上半平面两个极点，弧长乘最大值控制大弧。\par\textbf{详细解答} 极点$z_1=\e^{\frac{i\pi}{4}},z_2=\e^{\frac{3i\pi}{4}}$，留数$\frac{1}{4z_j^3}$，和$=-\frac{i}{2\sqrt2}$。半径$R$大弧积分绝对值$\le\frac{\pi R}{R^4-1}\to0$。留数定理给实积分$2\pi i[-\frac{i}{2\sqrt2}]=\frac{\pi}{\sqrt2}$。由偶性半轴积分$\frac{\pi}{2\sqrt2}$。也可用前节四次分母的初等原函数在无穷两端取极限核对。\par\textbf{方法点评} 留数法免去实部分分式的长计算，初等法验证实值和符号。\par\textbf{变式与易错点} 必须取上半平面全部两个极点。\end{example}

\subsection{Fourier与Laplace的归一化、卷积和初值}\label{deep:integral_complex_deep:transforms}

\begin{method}{Fourier与Laplace的归一化、卷积和初值}\textbf{适用条件} 先选变换约定并验证绝对可积或指数增长控制。\par\textbf{步骤与依据} 全书实轴Fourier统一采用$\widehat f(\xi)=\int_\R f(x)\e^{-2\pi i\xi x}\dd x$；在反演成立的函数类上，$f(x)=\int_\R\widehat f(\xi)\e^{2\pi i\xi x}\dd\xi$，无额外系数。Laplace采用$\mathcal Lf(s)=\int_0^\infty\e^{-st}f(t)\dd t$。\par\textbf{易错点} 公式中的$2\pi$与初值项由约定决定。\end{method}

\begin{example}{竞赛拓展｜指数核、余弦变换与卷积}{难度：$\star\star\star$}\label{deep:ex:integral_complex_deep:fourierexp}\textbf{题目} 计算$f(x)=\e^{-a|x|}$的Fourier变换，$a>0$，并求$f*f$。\par\textbf{思路} 偶性化为余弦积分；卷积按$t<0,0<t<x,t>x$分段。\par\textbf{详细解答} $\widehat f(\xi)=2\int_0^\infty\e^{-ax}\cos(2\pi\xi x)\dd x=\frac{2a}{a^2+4\pi^2\xi^2}$。当$x\ge0$，$\int_\R\e^{-a|t|}\e^{-a|x-t|}\dd t=\frac{\e^{-ax}}{2a}+x\e^{-ax}+\frac{\e^{-ax}}{2a}=\e^{-ax}(x+\frac{1}{a})$；偶延拓得$(f*f)(x)=\e^{-a|x|}(|x|+\frac{1}{a})$。其变换为$[\frac{2a}{a^2+4\pi^2\xi^2}]^2$，可作第二种校验。\par\textbf{方法点评} 连接Fourier余弦变换与卷积性质。\par\textbf{变式与易错点} 本约定下$\widehat{f*g}=\widehat f\,\widehat g$；对Schwartz函数等满足交换条件的函数，还有$\widehat{fg}=\widehat f*\widehat g$，两式均无额外$2\pi$。角频率$\omega=2\pi\xi$的另一约定会改变逆变换及乘积公式的系数。\end{example}

\begin{example}{考研风格｜Laplace解受迫初值问题}{难度：$\star\star\star$}\label{deep:ex:integral_complex_deep:laplaceode}\textbf{题目} 解$y''+y=t,\ y(0)=1,y'(0)=0$。\par\textbf{思路} 导数变换包含初值，拆成标准逆变换。\par\textbf{详细解答} 令$Y=\mathcal Ly$，则$(s^2Y-s)+Y=\frac{1}{s^2}$，故$Y=\frac{s}{s^2+1}+\frac{1}{s^2(s^2+1)}=\frac{s}{s^2+1}+\frac{1}{s^2}-\frac{1}{s^2+1}$。逆变换得$y=\cos t+t-\sin t$。第二法待定系数取特解$t$，齐次解$A\cos t+B\sin t$，初值得$A=1,B=-1$。\par\textbf{方法点评} 本题待定系数更短；Laplace对分段外力和积分方程更统一。\par\textbf{变式与易错点} 漏写$-sy(0)-y'(0)$会得到错误齐次部分。\end{example}


\section{不定积分的选法总图}\label{ic:atlas}
不定积分是在一个连通区间内寻找导数等于被积函数的函数族。先整理，再看结构，最后决定是否需要特殊函数。下列分类覆盖本册使用的初等、代数、构造与标准特殊函数路线，每类方法均配有公式依据和可检验例题。
\begin{center}\begin{tikzpicture}[every node/.style={draw=Teal,rounded corners,align=center,font=\small,minimum height=9mm},>=stealth]
\node (a) at (0,0) {确定定义域\\整理被积式};
\node (b) at (-4,-1.6) {导数匹配\\换元与凑微分};
\node (c) at (0,-1.6) {乘积或封闭函数族\\分部与待定原函数};
\node (d) at (4,-1.6) {代数或对称约束\\有理化、双元与递推};
\node (e) at (0,-3.2) {约化后剩标准核\\级数、含参或特殊函数};
\node (f) at (0,-4.7) {回代并求导验算\\检查分支与区间常数};
\draw[->](a)--(b);\draw[->](a)--(c);\draw[->](a)--(d);
\draw[->](b)--(e);\draw[->](c)--(e);\draw[->](d)--(e);\draw[->](e)--(f);
\end{tikzpicture}\end{center}
\begin{warning}疑难点一：相同外观不代表相同微分，先算$\dd t$再判断是否匹配。疑难点二：换元的极点可能只是坐标奇点。疑难点三：跨原奇点的区间各有积分常数。疑难点四：找不到初等表达不等于没有原函数。\end{warning}
\subsection{基本原函数的校准表}
\begin{longtable}{p{49mm}p{67mm}}\toprule 被积式与条件&原函数（另加常数）\\\midrule\endhead
$x^\alpha,\ \alpha\ne-1$&$\frac{x^{\alpha+1}}{\alpha+1}$，实幂先固定定义区间。\\
$\frac{1}{x},\ x\ne0$&$\ln|x|$。\\
$\e^{ax},\ a\ne0$&$\frac{\e^{ax}}{a}$。\\
$\sin ax,\cos ax,\ a\ne0$&$-\frac{\cos ax}{a},\ \frac{\sin ax}{a}$。\\
$\frac{1}{a^2+x^2},\ a>0$&$a^{-1}\arctan(\frac{x}{a})$。\\
$\frac{1}{\sqrt{a^2-x^2}},\ |x|<a$&$\arcsin(\frac{x}{a})$。\\
$\frac{1}{\sqrt{x^2+a^2}},\ a>0$&$\ln(x+\sqrt{x^2+a^2})$。\\
$\frac{1}{\sqrt{x^2-a^2}},\ a>0,|x|>a$&$\ln|x+\sqrt{x^2-a^2}|$。\\
$\sec^2x,\csc^2x$&$\tan x,\ -\cot x$。\\
$\sec x\tan x,\csc x\cot x$&$\sec x,\ -\csc x$。\\
$\sec x,\csc x$&$\ln|\sec x+\tan x|,\ \ln|\csc x-\cot x|$。\\
\bottomrule\end{longtable}
每式均由对应求导公式反向验证。含常参数的公式须另检参数为零的退化情形。
\subsection{详细方法导航}
\begin{longtable}{p{9mm}p{91mm}p{23mm}}\toprule 序号&方法&正文位置\\\midrule\endhead
1&先整理：除法、配方、拆项与逆链式&\ref{ic:indefinite_classical_complete:normalize}\\
2&对数、反函数与整体换元&\ref{ic:indefinite_classical_complete:logsub}\\
3&分部积分：降阶、升阶与反向选择&\ref{ic:indefinite_classical_complete:parts}\\
4&循环分部与有限维待定原函数&\ref{ic:indefinite_classical_complete:cyclic}\\
5&部分分式：重根、二次因子与不完全分解&\ref{ic:indefinite_classical_complete:partial}\\
6&三角有理式的奇偶捷径&\ref{ic:indefinite_classical_complete:trigparity}\\
7&三角幂：留一个微分、降幂与递推&\ref{ic:indefinite_classical_complete:trigpowers}\\
8&双曲有理式与幂递推&\ref{ic:indefinite_classical_complete:hyperbolic}\\
9&Euler三替换的反解公式&\ref{ic:indefinite_classical_complete:eulerthree}\\
10&二次根式正负幂的递推&\ref{ic:indefinite_classical_complete:powerrec}\\
11&圆型、双曲型与齐次比值&\ref{ic:indefinite_dual_complete:differentials}\\
12&恒等式诱导与莫比乌斯变换的构造&\ref{ic:indefinite_dual_complete:mobius}\\
13&对勾三元与回文多项式&\ref{ic:indefinite_dual_complete:reciprocal}\\
14&双元化单元与根式和差参数&\ref{ic:indefinite_dual_complete:single}\\
15&高次双元：统一微分与有理参数&\ref{ic:indefinite_dual_complete:highdual}\\
16&统一有理积分族与公因数降阶&\ref{ic:indefinite_dual_complete:ufamily}\\
17&隐函数：参数化与约束微分&\ref{ic:indefinite_dual_complete:implicit}\\
18&二项微分的三条有理化路线&\ref{ic:indefinite_algebra_complete:binomial}\\
19&Hermite约化：模多项式算法&\ref{ic:indefinite_algebra_complete:hermitemod}\\
20&有理积分的留数式系数法&\ref{ic:indefinite_algebra_complete:residue}\\
21&平方分母的商导数降阶&\ref{ic:indefinite_algebra_complete:quotient}\\
22&反三角混合锁与有限维闭合&\ref{ic:indefinite_algebra_complete:triglock}\\
23&指数锁与归一化变量&\ref{ic:indefinite_algebra_complete:explock}\\
24&等分锁与对数微分配对&\ref{ic:indefinite_algebra_complete:equalparts}\\
25&代数原函数待定与伪椭圆判别&\ref{ic:indefinite_algebra_complete:exactalg}\\
26&四次根式的多项式连分式算法&\ref{ic:indefinite_algebra_complete:continuedfraction}\\
27&微分方程引领积分：Lagrange恒等式&\ref{ic:indefinite_algebra_complete:odeidentity}\\
28&含参不定积分须固定积分常数&\ref{ic:indefinite_special_complete:parameter}\\
29&幂级数逐项积分与命名函数&\ref{ic:indefinite_special_complete:series}\\
30&复指数与实原函数的还原&\ref{ic:indefinite_special_complete:complex}\\
31&Ei、li、Si、Ci的识别核&\ref{ic:indefinite_special_complete:ei}\\
32&Gaussian、erf、erfi与Fresnel&\ref{ic:indefinite_special_complete:gauss}\\
33&多重对数的逐阶降阶&\ref{ic:indefinite_special_complete:polylog}\\
34&Lambert W：用隐式关系换回单元&\ref{ic:indefinite_special_complete:lambert}\\
35&不完全Beta与Gamma的统一幂权&\ref{ic:indefinite_special_complete:incomplete}\\
36&三类椭圆标准化与实参数选择&\ref{ic:indefinite_special_complete:elliptic}\\
37&初等表达的边界与最终验算&\ref{ic:indefinite_special_complete:boundary}\\
\bottomrule\end{longtable}

\section{初等不定积分：从识别到算法}\label{ic:indefinite_classical_complete}

\begin{examtip}以下按被积函数的结构展开计算步骤；同类方法优先用最短换元，不把万能代换当作每题的第一步。\end{examtip}

\subsection{先整理：除法、配方、拆项与逆链式}\label{ic:indefinite_classical_complete:normalize}\begin{method}{先整理：除法、配方、拆项与逆链式}\textbf{识别信号} 分式分子次数不低于分母；分子接近分母导数。\par \textbf{操作步骤} 先长除；二次式配方；把分子拆成$A Q'(x)+B$；最后确认是否为$\frac{f'(x)}{f(x)}$或$f'(x)f(x)^r$。\par \textbf{疑难点与条件} 线性只适用于和及常数倍；积分不能分配到商或乘积。\end{method}

\begin{example}{考研风格｜长除与凑微分}{难度：$\star\star$}\label{ic:ex:indefinite_classical_complete:division}\textbf{题目} 求$\int\frac{x^3+2x^2+1}{x^2+1}\dd x$。\par \textbf{思路} 先除掉高次项，再拆导数项。\par \textbf{详细解答} 长除得$x+2+\frac{-x-1}{x^2+1}$，故原函数$\frac{x^2}{2}+2x-\frac12\ln(x^2+1)-\arctan x+C$。另一法先把分子写成$(x+2)(x^2+1)-(x+1)$，再逐项积分。\par \textbf{求导验算} 求导为$x+2-\frac{x+1}{x^2+1}$，通分恢复原分子。\par \textbf{疑难点与方法比较} 先作代数整理通常比立即换元省力。\end{example}

\begin{example}{考研风格｜导数配分母}{难度：$\star\star$}\label{ic:ex:indefinite_classical_complete:affine}\textbf{题目} 求$\int\frac{3x+4}{x^2+2x+5}\dd x$。\par \textbf{思路} $3x+4=\frac32(2x+2)+1$。\par \textbf{详细解答} 配方$x^2+2x+5=(x+1)^2+4$，答案$\frac32\ln(x^2+2x+5)+\frac12\arctan(\frac{x+1}{2})+C$。\par \textbf{求导验算} 导数$\frac{3x+3}{Q}+\frac1Q=\frac{3x+4}{Q}$。\par \textbf{疑难点与方法比较} 对数项与反正切项的系数分别由配导数和配方决定。\end{example}

\subsection{对数、反函数与整体换元}\label{ic:indefinite_classical_complete:logsub}\begin{method}{对数、反函数与整体换元}\textbf{识别信号} 出现$\frac{\dd x}{x}$配$\ln x$，或$\frac{\dd x}{1+x^2}$配$\arctan x$。\par \textbf{操作步骤} 将内层整体命名为$t$；同时替换自变量因子与微分；若换元后仍含$x$，用反函数回写。\par \textbf{疑难点与条件} 换元是局部可逆变换；不跨越对数、反三角或原分母的奇点。\end{method}

\begin{example}{考研风格｜反三角幂整体换元}{难度：$\star\star$}\label{ic:ex:indefinite_classical_complete:atanpower}\textbf{题目} 求$\int\frac{(\arctan x)^4\dd x}{1+x^2}$。\par \textbf{思路} 内层导数完整出现。\par \textbf{详细解答} 令$t=\arctan x$，$\dd t=\frac{\dd x}{1+x^2}$，积分$\frac{t^5}{5}+C=\frac{(\arctan x)^5}{5}+C$。\par \textbf{求导验算} 链式求导立即还原被积函数。\par \textbf{疑难点与方法比较} 反三角函数不总要分部；出现其导数时先换元。\end{example}

\begin{example}{考研风格｜对数换元后再分部}{难度：$\star\star\star$}\label{ic:ex:indefinite_classical_complete:xlog}\textbf{题目} 求$\int(\ln x)^2\dd x$，$x>0$。\par \textbf{思路} 可直接分部，也可令$t=\ln x$化指数多项式。\par \textbf{详细解答} 分部两次得$x(\ln x)^2-2\int\ln x\dd x=x[(\ln x)^2-2\ln x+2]+C$。另一法$x=\e^t$，积分$\int t^2\e^t\dd t=\e^t(t^2-2t+2)+C$，回代相同。\par \textbf{求导验算} 导数为$(\ln x)^2-2\ln x+2+2\ln x-2=(\ln x)^2$。\par \textbf{疑难点与方法比较} 换元可以把反函数积分转为指数多项式族。\end{example}

\subsection{分部积分：降阶、升阶与反向选择}\label{ic:indefinite_classical_complete:parts}\begin{method}{分部积分：降阶、升阶与反向选择}\textbf{识别信号} 乘积中一个因子求导会简化；或另一个因子积分会引入可用结构。\par \textbf{操作步骤} 先预测剩余积分，再选择$u,\dd v$；多项式反复求导；反三角或对数常求导；若遇循环，将未知积分移到同一侧。\par \textbf{疑难点与条件} 经验顺序不是定理，必须比较剩余积分是否变简单。\end{method}

\begin{example}{考研风格｜反向分部选择}{难度：$\star\star\star$}\label{ic:ex:indefinite_classical_complete:logsec}\textbf{题目} 求$\int x\ln(1+x^2)\dd x$。\par \textbf{思路} 令$t=1+x^2$再分部，比把$x$作为求导因子更短。\par \textbf{详细解答} $t=1+x^2$给$\frac12\int\ln t\dd t=\frac12[(1+x^2)\ln(1+x^2)-(1+x^2)]+C$。直接分部取$u=\ln(1+x^2),\dd v=x\dd x$也得$\frac{x^2}2\ln(1+x^2)-\int \frac{x^3}{1+x^2}\dd x$，长除后等价。\par \textbf{求导验算} 对第一式求导为$x\ln(1+x^2)$。\par \textbf{疑难点与方法比较} 先换元再分部可避免高次有理分式。\end{example}

\begin{example}{考研风格｜反三角函数分部}{难度：$\star\star\star$}\label{ic:ex:indefinite_classical_complete:arcsin}\textbf{题目} 求$\int\arcsin x\dd x$，$|x|<1$。\par \textbf{思路} 把隐含因子1积分成$x$。\par \textbf{详细解答} 取$u=\arcsin x,\dd v=\dd x$，积分$x\arcsin x-\int \frac{x}{\sqrt{1-x^2}}\dd x=x\arcsin x+\sqrt{1-x^2}+C$。第二法$x=\sin t$，对$\int t\cos t\dd t$分部，也得同式。\par \textbf{求导验算} 两项中的$\frac{x}{\sqrt{1-x^2}}$相消。\par \textbf{疑难点与方法比较} 最后根式前为正号；区间端点的可导性单独讨论。\end{example}

\subsection{循环分部与有限维待定原函数}\label{ic:indefinite_classical_complete:cyclic}\begin{method}{循环分部与有限维待定原函数}\textbf{识别信号} 指数乘三角、双曲，或求导后封闭于有限函数族。\par \textbf{操作步骤} 写联立积分关系，或设$F=\sum A_jf_j$并比较导数系数；解线性方程；最后检验共振或参数退化。\par \textbf{疑难点与条件} 循环回到原积分时必须移项，不能将中间积分误当结果。\end{method}

\begin{example}{考研风格｜循环分部与复指数三解}{难度：$\star\star\star$}\label{ic:ex:indefinite_classical_complete:expcycle}\textbf{题目} 求$\int\e^{2x}\sin3x\dd x$。\par \textbf{思路} 二次循环、待定系数或复指数。\par \textbf{详细解答} 设$I=\int\e^{2x}\sin3x$、$J=\int\e^{2x}\cos3x$。分部给$I=\frac{\e^{2x}\sin3x}{2}-\frac{3J}{2}$、$J=\frac{\e^{2x}\cos3x}{2}+\frac{3I}{2}$，解得$I=\frac{\e^{2x}(2\sin3x-3\cos3x)}{13}+C$。设$\e^{2x}(A\sin3x+B\cos3x)$比较也得$A=\frac{2}{13},B=-\frac{3}{13}$；复指数法取$\frac{\e^{(2+3i)x}}{2+3i}$虚部。\par \textbf{求导验算} 求导后正弦系数$\frac{4+9}{13}=1$，余弦系数抵消。\par \textbf{疑难点与方法比较} 待定原函数不是猜最终系数，而是解确定的线性方程组。\end{example}

\subsection{部分分式：重根、二次因子与不完全分解}\label{ic:indefinite_classical_complete:partial}\begin{method}{部分分式：重根、二次因子与不完全分解}\textbf{识别信号} 实有理式先化真分式。\par \textbf{操作步骤} 对$(x-a)^m$列全幂；对不可约二次$Q^m$每幂配置一次分子；可先分离特定因子，避免不必要的全分解。\par \textbf{疑难点与条件} 分解恒等式只在原函数定义的区间使用，不能跨极点合并常数。\end{method}

\begin{example}{考研风格｜一次重根的系统配置}{难度：$\star\star$}\label{ic:ex:indefinite_classical_complete:repeated}\textbf{题目} 求$\int\frac{\dd x}{x^2(x+1)}$。\par \textbf{思路} 重根$x^2$须保留$\frac{A}{x}+\frac{B}{x^2}$。\par \textbf{详细解答} $\frac{1}{x^2(x+1)}=-\frac{1}{x}+\frac{1}{x^2}+\frac{1}{x+1}$，故答案$-\ln|x|-\frac{1}{x}+\ln|x+1|+C$。\par \textbf{求导验算} 通分后的分子$-x(x+1)+(x+1)+x^2=1$。\par \textbf{疑难点与方法比较} 漏掉$\frac{1}{x}$项会使系数方程无解。\end{example}

\begin{example}{竞赛风格｜不可约二次重因子}{难度：$\star\star\star$}\label{ic:ex:indefinite_classical_complete:quadraticrepeated}\textbf{题目} 求$\int \frac{x^3\dd x}{(1+x^2)^2}$。\par \textbf{思路} 可用$t=1+x^2$，也可在原变量约化。\par \textbf{详细解答} $t=1+x^2$给$\frac12\int(t-1)t^{-2}\dd t=\frac12\ln(1+x^2)+\frac1{2(1+x^2)}+C$。另一法$\frac{x^3}{(1+x^2)^2}=\frac{x}{1+x^2}-\frac{x}{(1+x^2)^2}$，分别凑微分。\par \textbf{求导验算} 导数$\frac{x}{1+x^2}-\frac{x}{(1+x^2)^2}=\frac{x^3}{(1+x^2)^2}$。\par \textbf{疑难点与方法比较} 二次重因子遇奇次分子通常优先平方变量换元。\end{example}

\subsection{三角有理式的奇偶捷径}\label{ic:indefinite_classical_complete:trigparity}\begin{method}{三角有理式的奇偶捷径}\textbf{识别信号} 被积函数为$R(\sin x,\cos x)$。\par \textbf{操作步骤} 若对$\sin x$为奇函数，取$t=\cos x$；对$\cos x$为奇函数，取$t=\sin x$；同时变号不变，取$t=\tan x$；其余用$t=\tan(\frac{x}{2})$。\par \textbf{疑难点与条件} “奇”指整个有理表达式的符号变换性质；代换后用$\sin^2+\cos^2=1$化有理。\end{method}

\begin{example}{考研风格｜对正弦为奇的有理式}{难度：$\star\star$}\label{ic:ex:indefinite_classical_complete:sinodd}\textbf{题目} 求$\int\frac{\sin x\dd x}{1+\cos^2x}$。\par \textbf{思路} $t=\cos x$直接消去正弦微分。\par \textbf{详细解答} 积分$=-\int\frac{\dd t}{1+t^2}=-\arctan(\cos x)+C$。半角法也可有理化，但分母次数更高。\par \textbf{求导验算} 导数为$\frac{\sin x}{1+\cos^2x}$。\par \textbf{疑难点与方法比较} 先用符号对称性选变量，避免无条件使用万能代换。\end{example}

\begin{example}{考研风格｜同时变号不变}{难度：$\star\star\star$}\label{ic:ex:indefinite_classical_complete:simultaneous}\textbf{题目} 求$\int\frac{\dd x}{2\sin^2x+3\cos^2x}$。\par \textbf{思路} 除以$\cos^2x$后取$t=\tan x$。\par \textbf{详细解答} 在$\cos x\ne0$的区间，积分$=\int\frac{\dd t}{2t^2+3}=\frac1{\sqrt6}\arctan(\sqrt{\frac{2}{3}}\tan x)+C$。如需覆盖$\cos x=0$处，应在相邻区间调整常数，使原函数连续衔接。\par \textbf{求导验算} 导数为$\frac{\sec^2x}{3+2\tan^2x}$，乘$\cos^2x$恢复原式。\par \textbf{疑难点与方法比较} 原被积函数全实轴光滑，代换的极点不是原函数必有的奇点。\end{example}

\subsection{三角幂：留一个微分、降幂与递推}\label{ic:indefinite_classical_complete:trigpowers}\begin{method}{三角幂：留一个微分、降幂与递推}\textbf{识别信号} $\sin^mx\cos^nx$，或$\tan,\sec,\cot,\csc$的整数幂。\par \textbf{操作步骤} 奇幂留一个$\sin$或$\cos$；全偶用倍角；$\sec$偶幂留$\sec^2$，$\tan$奇幂配$\sec\tan$；无直接微分时分部递推。\par \textbf{疑难点与条件} 负幂也能递推，但分母零点必须从定义域剔除。\end{method}

\begin{example}{考研风格｜正切正割混合幂双解}{难度：$\star\star\star$}\label{ic:ex:indefinite_classical_complete:tansecpower}\textbf{题目} 求$\int\tan^3x\sec^4x\dd x$。\par \textbf{思路} 可留$\sec^2x$用$\tan x$，或留$\sec x\tan x$用$\sec x$。\par \textbf{详细解答} 取$t=\tan x$，得到$\int t^3(1+t^2)\dd t=\frac{t^4}{4}+\frac{t^6}{6}+C$。另一法$u=\sec x$，$\tan^2x=u^2-1$，积分$\int(u^2-1)u^3\dd u=\frac{u^6}{6}-\frac{u^4}{4}+C$，两式相差常数。\par \textbf{求导验算} 第一式求导为$(\tan^3+\tan^5)\sec^2=\tan^3\sec^4$。\par \textbf{疑难点与方法比较} 两条换元均可行时比较多项式次数。\end{example}

\begin{example}{考研风格｜不定积分降二阶递推}{难度：$\star\star\star$}\label{ic:ex:indefinite_classical_complete:sinrec}\textbf{题目} 推导$S_n=\int\sin^nx\dd x$并求$S_5$。\par \textbf{思路} 对$\sin^{n-1}x\sin x$分部。\par \textbf{详细解答} $S_n=-\sin^{n-1}x\cos x+(n-1)\int\sin^{n-2}x\cos^2x\dd x$，故$S_n=-\frac{\sin^{n-1}x\cos x}{n}+\frac{(n-1)S_{n-2}}{n}$。初值$S_0=x,S_1=-\cos x$。也可留$\sin x$令$u=\cos x$，直接得$S_5=-\cos x+\frac23\cos^3x-\frac15\cos^5x+C$。\par \textbf{求导验算} 多项式式求导为$\sin x(1-2\cos^2x+\cos^4x)=\sin^5x$。\par \textbf{疑难点与方法比较} 定积分Wallis递推是边界项消失后的特例，不定积分不能丢边界项。\end{example}

\subsection{双曲有理式与幂递推}\label{ic:indefinite_classical_complete:hyperbolic}\begin{method}{双曲有理式与幂递推}\textbf{识别信号} 由$\sinh,\cosh,\tanh$构成的有理式。\par \textbf{操作步骤} 通法$t=\e^x$；半角$t=\tanh(\frac{x}{2})$也可；幂函数用$\cosh^2-\sinh^2=1$，注意平方恒等式符号。\par \textbf{疑难点与条件} 实$\cosh x>0$，但$\sinh x=0$会产生分母奇点；$\sqrt{\sinh^2x}=|\sinh x|$。\end{method}

\begin{example}{考研风格｜双曲半角有理化}{难度：$\star\star\star$}\label{ic:ex:indefinite_classical_complete:hyperbolicrat}\textbf{题目} 求$\int\frac{\dd x}{2+\cosh x}$。\par \textbf{思路} 取$t=\tanh(\frac{x}{2})$，$|t|<1$。\par \textbf{详细解答} $\cosh x=\frac{1+t^2}{1-t^2}$、$\dd x=\frac{2\dd t}{1-t^2}$，积分$=2\int\frac{\dd t}{3-t^2}=\frac2{\sqrt3}\operatorname{artanh}(\frac{t}{\sqrt3})+C$。也可$t=\e^x>0$，积分化为$2\int\frac{\dd t}{t^2+4t+1}$，配方给对数形式。\par \textbf{求导验算} $\frac{\dd}{\dd x}\operatorname{artanh}(\frac{t}{\sqrt3})=\frac{t'}{\sqrt3(1-\frac{t^2}{3})}$，代$t'=\frac{1-t^2}{2}$还原原式。\par \textbf{疑难点与方法比较} 双曲型二次分母可能导出反双曲函数或实对数，不能误写成反正切。\end{example}

\begin{example}{竞赛风格｜双曲幂递推的符号}{难度：$\star\star\star$}\label{ic:ex:indefinite_classical_complete:hyperbolicpower}\textbf{题目} 令$C_n=\int\cosh^nx\dd x$，推导递推并求$C_3$。\par \textbf{思路} 分部并用$\sinh^2=\cosh^2-1$。\par \textbf{详细解答} $C_n=\cosh^{n-1}x\sinh x-(n-1)\int\cosh^{n-2}x\sinh^2x\dd x$，所以$C_n=\frac{\cosh^{n-1}x\sinh x}{n}+\frac{(n-1)C_{n-2}}{n}$。$C_3=\sinh x+\frac{\sinh^3x}{3}+C$，也由$t=\sinh x$得到。\par \textbf{求导验算} 导数$\cosh x(1+\sinh^2x)=\cosh^3x$。\par \textbf{疑难点与方法比较} 与正弦递推形式相近，但边界项符号不同。\end{example}

\subsection{Euler三替换的反解公式}\label{ic:indefinite_classical_complete:eulerthree}\begin{method}{Euler三替换的反解公式}\textbf{识别信号} $Q=ax^2+bx+c$为实二次式，$\sqrt Q$所在支路可微。\par \textbf{操作步骤} 若$a>0$设$\sqrt Q=\sqrt a\,x+t$，得$x=\frac{t^2-c}{b-2\sqrt a\,t}$；若$c>0$设$\sqrt Q=tx+\sqrt c$，得$x=\frac{2\sqrt c\,t-b}{a-t^2}$；若$Q=a(x-r)(x-s)$有实根，设$\sqrt Q=t(x-r)$，得$x=\frac{t^2r-as}{t^2-a}$。\par \textbf{疑难点与条件} 反解分母为0的参数需单独检查；平方后必须代回原式检验根式符号。三类方法至少一类可用于非空实正则区间。\end{method}

\begin{example}{考研风格｜第一Euler替换的完整计算}{难度：$\star\star\star$}\label{ic:ex:indefinite_classical_complete:eulerfirst}\textbf{题目} 求$\int\frac{\dd x}{\sqrt{x^2+2x+2}}$。\par \textbf{思路} 第一类替换，或先平移配方。\par \textbf{详细解答} 设$t=\sqrt{x^2+2x+2}-x$，则$x=\frac{t^2-2}{2-2t}$。由$\sqrt Q=x+t$微分，$\frac{(x+1)\dd x}{\sqrt Q}=\dd x+\dd t$，结合$x(t)$可化为$\frac{\dd x}{\sqrt Q}=-\frac{\dd t}{t-1}$，所以答案$-\ln|t-1|+C$。更短法$u=x+1$，结果$\ln(u+\sqrt{1+u^2})+C$。因$(\sqrt Q-x-1)(\sqrt Q+x+1)=1$，两式相同。\par \textbf{求导验算} $\frac{\dd\ln(x+1+\sqrt Q)}{\dd x}=\frac{1}{\sqrt Q}$。\par \textbf{疑难点与方法比较} 通法解释为何可有理化，配方通常是本题的优先方法。\end{example}

\begin{example}{考研风格｜第三Euler替换与比值根式}{难度：$\star\star\star$}\label{ic:ex:indefinite_classical_complete:eulerthird}\textbf{题目} 求$\int\frac{\dd x}{\sqrt{(x-1)(x+2)}}$，$x>1$。\par \textbf{思路} 以实根1为基点，设$t=\frac{\sqrt Q}{x-1}$。\par \textbf{详细解答} $t^2=\frac{x+2}{x-1}$，得$x=\frac{t^2+2}{t^2-1}$、$\dd x=-\frac{6t\dd t}{(t^2-1)^2}$、$\sqrt Q=\frac{3t}{t^2-1}$（$t>1$），故积分$=-2\int\frac{\dd t}{t^2-1}=\ln|\frac{t+1}{t-1}|+C$。也可先配方$Q=(x+\frac{1}{2})^2-(\frac{3}{2})^2$，得到$\ln(x+\frac{1}{2}+\sqrt Q)+C$。\par \textbf{求导验算} 对配方结果求导为$\frac{1}{\sqrt Q}$；与参数式只差常数。\par \textbf{疑难点与方法比较} 在$x<-2$支路须重新选择符号，不能照抄$t>1$。\end{example}

\subsection{二次根式正负幂的递推}\label{ic:indefinite_classical_complete:powerrec}\begin{method}{二次根式正负幂的递推}\textbf{识别信号} $\int(x^2+a^2)^{\frac{m}{2}}\dd x$，指数可逐步降到基本积分。\par \textbf{操作步骤} 对$x(x^2+a^2)^{\frac{m}{2}}$求导并重组；正整数与负整数分别从低阶初值递推；$m=-1$等退化位置改用基本对数式。\par \textbf{疑难点与条件} 递推公式的分母为0时不是积分不存在，而是该公式不适用。\end{method}

\begin{theorem}{统一幂递推}令$J_p=\int(x^2+a^2)^p\dd x$，则$p\ne-\frac{1}{2}$时$J_p=\frac{x(x^2+a^2)^p}{2p+1}+\frac{2pa^2}{2p+1}J_{p-1}$。\end{theorem}\begin{proof}$[x(x^2+a^2)^p]'=(2p+1)(x^2+a^2)^p-2pa^2(x^2+a^2)^{p-1}$，移项积分。$p=-\frac{1}{2}$的初值为$\ln(x+\sqrt{x^2+a^2})$（$a>0$）。\end{proof}

\begin{example}{竞赛风格｜根式三次幂的两解}{难度：$\star\star\star$}\label{ic:ex:indefinite_classical_complete:positivehalf}\textbf{题目} 求$\int(1+x^2)^{\frac{3}{2}}\dd x$。\par \textbf{思路} 统一递推或双曲换元。\par \textbf{详细解答} 递推$p=\frac{3}{2}$给$J_{\frac{3}{2}}=\frac{x(1+x^2)^{\frac{3}{2}}}{4}+\frac{3J_{\frac{1}{2}}}{4}$，而$J_{\frac{1}{2}}=\frac{x\sqrt{1+x^2}+\ln(x+\sqrt{1+x^2})}{2}$。故结果$\frac{x(2x^2+5)\sqrt{1+x^2}}8+\frac38\ln(x+\sqrt{1+x^2})+C$。另一法$x=\sinh t$，积分$\cosh^4t$后降幂。\par \textbf{求导验算} 对两项求导、通分后分子为$8(1+x^2)^2$，恢复$(1+x^2)^{\frac{3}{2}}$。\par \textbf{疑难点与方法比较} 正幂递推降低指数，负幂则需把恒等式反向求解。\end{example}

\section{双元、单元、高次关系与隐函数积分全流程}\label{ic:indefinite_dual_complete}

\begin{examtip}本节把双元名称还原为约束微分、齐次化与有理参数。有理积分采用记号$\mathcal U$，各系数由微分推导。\end{examtip}

\subsection{圆型、双曲型与齐次比值}\label{ic:indefinite_dual_complete:differentials}\begin{method}{圆型、双曲型与齐次比值}\textbf{识别信号} 两个可微量$u,v$满足$u^2+v^2=a^2$或$v^2-u^2=a^2$。\par \textbf{操作步骤} 圆型取角或比值；双曲型取$u+v$；同时记录约束的微分。商导数给$\dd(\frac{u}{v})$，乘积导数给分部闭合。\par \textbf{疑难点与条件} 角变量公式局部成立，跨象限须补连续分支；根式符号不能任意取消。\end{method}

\begin{theorem}{角微分与双曲对数微分}设$a>0$，并在$v\ne0$的可微实支上讨论。若$u^2+v^2=a^2$，则$\dd\arctan(\frac{u}{v})=\frac{v\dd u-u\dd v}{a^2}=\frac{\dd u}{v}$，在$v\ne0$的局部角支成立。若$v^2-u^2=a^2$，则$\dd\ln|u+v|=\frac{\dd u}{v}$。\end{theorem}\begin{proof}分别用反正切导数和约束$v\dd v=\mp u\dd u$计算。圆型中$v\dd u-u\dd v=\frac{(v^2+u^2)\dd u}{v}$；双曲型$\frac{\dd u+\dd v}{u+v}=\frac{\dd u}{v}$。\end{proof}

\begin{example}{考研风格｜圆型非齐次分式}{难度：$\star\star\star$}\label{ic:ex:indefinite_dual_complete:circlefraction}\textbf{题目} 求$\int\frac{\dd x}{\sqrt{1-x^2}(1+\sqrt{1-x^2})}$，$|x|<1$。\par \textbf{思路} 先取圆的角参数，再半角。\par \textbf{详细解答} $x=\sin t$，$t\in(-\frac{\pi}{2},\frac{\pi}{2})$，积分$\int\frac{\dd t}{1+\cos t}=\tan(\frac{t}{2})+C=\frac{x}{1+\sqrt{1-x^2}}+C$。另一法直接对商求导，利用$x^2+(1-x^2)=1$整理。\par \textbf{求导验算} 若$v=\sqrt{1-x^2}$，$[\frac{x}{1+v}]'=\frac{1+v+\frac{x^2}{v}}{(1+v)^2}=\frac{1}{v(1+v)}$。\par \textbf{疑难点与方法比较} 先除去一个圆型微分，再处理常数加双元，比分母全体有理化更短。\end{example}

\begin{example}{竞赛风格｜双元原函数再分部}{难度：$\star\star\star$}\label{ic:ex:indefinite_dual_complete:dual_log}\textbf{题目} 求$\int\ln(x+\sqrt{1+x^2})\dd x$。\par \textbf{思路} 对数的导数已由双元公式确定。\par \textbf{详细解答} 令$v=\sqrt{1+x^2}$，分部得$x\ln(x+v)-\int \frac{x\dd x}{v}=x\ln(x+v)-v+C$。也可$x=\sinh t$，积分$\int t\cosh t\dd t=t\sinh t-\cosh t+C$。\par \textbf{求导验算} 导数为$\ln(x+v)+\frac{x}{v}-\frac{x}{v}$。\par \textbf{疑难点与方法比较} $x+v>0$，此处实对数无需绝对值。\end{example}

\subsection{恒等式诱导与莫比乌斯变换的构造}\label{ic:indefinite_dual_complete:mobius}\begin{method}{恒等式诱导与莫比乌斯变换的构造}\textbf{识别信号} 分母含一次式与二次根式的乘积。\par \textbf{操作步骤} 把二次式写成两个一次式平方和或差，分别除以同一个一次式平方；所产生的两个比值满足圆或双曲关系；计算比值微分。\par \textbf{疑难点与条件} 不能只写“注意到恒等式”，须给平方展开或系数比较。\end{method}

\begin{example}{竞赛风格｜一次分母乘二次根式双解}{难度：$\star\star\star$}\label{ic:ex:indefinite_dual_complete:mobiusroot}\textbf{题目} 求$\int\frac{\dd x}{(x+1)\sqrt{x^2+1}}$，$x\ne-1$。\par \textbf{思路} 用$2(x^2+1)=(x+1)^2+(x-1)^2$产生双曲双元。\par \textbf{详细解答} 设$u=\frac{x-1}{x+1},v=\frac{\sqrt2\sqrt{x^2+1}}{x+1}$，有$v^2-u^2=1$，且$\frac{\dd u}{v}=\frac{\sqrt2\,\dd x}{(x+1)\sqrt{x^2+1}}$。所以原函数$\frac1{\sqrt2}\ln\left|\frac{x-1+\sqrt2\sqrt{x^2+1}}{x+1}\right|+C$。第二法$x=\sinh t$，积分$\int\frac{\dd t}{1+\sinh t}$，再$\e^t$化有理式，可得等价对数。\par \textbf{求导验算} 由$\dd\ln|u+v|=\frac{\dd u}{v}$直接验证，$v$可正可负但不为0。\par \textbf{疑难点与方法比较} 以带符号的$v$保留分支，比无条件替成正平方根安全。\end{example}

\begin{example}{竞赛风格｜四次根式的平方分解}{难度：$\star\star\star$}\label{ic:ex:indefinite_dual_complete:quarticidentity}\textbf{题目} 求$\int\frac{\dd x}{x\sqrt{1+x^2+3x^4}}$，$x>0$。\par \textbf{思路} 除以$x^4$后配方，或先倒数平方换元。\par \textbf{详细解答} 令$t=x^{-2}$，$\dd t=-2x^{-3}\dd x$，原式$=-\frac12\int\frac{\dd t}{\sqrt{t^2+t+3}}$。配方得$-\frac12\ln[t+\frac{1}{2}+\sqrt{t^2+t+3}]+C$。等价地利用$4(1+x^2+3x^4)=(x^2+2)^2+11x^4$构造双曲双元，得到相同对数形式。\par \textbf{求导验算} 对$t$的原函数导数为$-\frac{1}{2\sqrt{t^2+t+3}}$，乘$t'$恢复被积式。\par \textbf{疑难点与方法比较} 分母的$x$是倒数平方换元的信号；去掉它则不能沿用本解。\end{example}

\subsection{对勾三元与回文多项式}\label{ic:indefinite_dual_complete:reciprocal}\begin{method}{对勾三元与回文多项式}\textbf{识别信号} 同时出现$x^2$与$x^{-2}$，或四次回文多项式。\par \textbf{操作步骤} 设$p=x+\frac{a}{x},q=x-\frac{a}{x}$，则$p^2-q^2=4a$；根据微分中的$1\mp \frac{a}{x^2}$选变量；根式除以$|x|$后识别二次式。\par \textbf{疑难点与条件} 回文结构不保证任意分子都能直接换元；原微分的匹配决定是否初等化。\end{method}

\begin{example}{竞赛风格｜同一四次根式的两种信号}{难度：$\star\star\star$}\label{ic:ex:indefinite_dual_complete:reciprocalplus}\textbf{题目} 求$\int\frac{(x^2+1)\dd x}{x\sqrt{x^4+5x^2+1}}$，$x>0$。\par \textbf{思路} 选$q=x-\frac{1}{x}$，因为$q'=1+x^{-2}$。\par \textbf{详细解答} $\sqrt{x^4+5x^2+1}=x\sqrt{q^2+7}$，故积分$\int\frac{\dd q}{\sqrt{q^2+7}}=\ln(q+\sqrt{q^2+7})+C$。如果分子改为$x^2-1$，改取$p=x+\frac{1}{x}$，得到$\ln(p+\sqrt{p^2+3})+C$。\par \textbf{求导验算} 分别用$q'=1+x^{-2},p'=1-x^{-2}$验证两题。\par \textbf{疑难点与方法比较} 分子正负决定用和或差；不能凭分母回文就固定一种换元。\end{example}

\begin{example}{竞赛风格｜对勾结构导出反正切}{难度：$\star\star\star$}\label{ic:ex:indefinite_dual_complete:reciprocalangle}\textbf{题目} 求$\int\frac{(1-x^2)\dd x}{(1+x^2)\sqrt{x^4+x^2+1}}$。\par \textbf{思路} 设$r=\sqrt{x^4+x^2+1}$，试角变量$\frac{x}{r}$。\par \textbf{详细解答} 计算$(\frac{x}{r})'=\frac{1-x^4}{r^3}$，而$1+(\frac{x}{r})^2=\frac{(x^2+1)^2}{r^2}$。所以被积函数等于$[\arctan(\frac{x}{r})]'$，答案$\arctan(\frac{x}{\sqrt{x^4+x^2+1}})+C$。也可先以$x\pm\frac{1}{x}$构造三元，约成$\int\frac{\dd u}{1+u^2}$。\par \textbf{求导验算} 上面两条导数比值化简为$\frac{1-x^2}{(1+x^2)r}$，在$x=0$也可直接连续验证。\par \textbf{疑难点与方法比较} 把直接验证看作第二层证据，不能省略识别出角变量的恒等式。\end{example}

\subsection{双元化单元与根式和差参数}\label{ic:indefinite_dual_complete:single}\begin{method}{双元化单元与根式和差参数}\textbf{识别信号} 双元和差乘积为常数，或多个有理分数幂共用基底。\par \textbf{操作步骤} 用和$t=u+v$并由乘积反解差；嵌套根式再取合适根次；多个分数幂取分母最小公倍数。\par \textbf{疑难点与条件} 参数须选与原根式一致的实支；幂换元对负自变量需另检。\end{method}

\begin{example}{竞赛风格｜根式和差同时参数化}{难度：$\star\star\star$}\label{ic:ex:indefinite_dual_complete:nestedtwo}\textbf{题目} 求$\int[\sqrt{x+1}-\sqrt x]^{-2}\dd x$，$x>0$。\par \textbf{思路} 令$t=\sqrt{x+1}+\sqrt x$，其倒数是两根之差。\par \textbf{详细解答} $\sqrt{x+1}=\frac{t+t^{-1}}{2},\sqrt x=\frac{t-t^{-1}}{2}$，所以$x=\frac{t^2-2+t^{-2}}{4}$、$\dd x=\frac{(t-t^{-3})\dd t}{2}$。原式$\frac12\int(t^3-t^{-1})\dd t=\frac{t^4}{8}-\frac12\ln t+C$。另一法先将被积式有理化为$(\sqrt{x+1}+\sqrt x)^2$，再分别积分，结果相差常数。\par \textbf{求导验算} 乘参数导数可回到$t^2\dd x$，而$t^2$正是原被积函数。\par \textbf{疑难点与方法比较} 同一和差参数同时消掉两根，避免逐层换元。\end{example}

\subsection{高次双元：统一微分与有理参数}\label{ic:indefinite_dual_complete:highdual}\begin{method}{高次双元：统一微分与有理参数}\textbf{识别信号} $y^n=\varepsilon x^n+c$，$n\ge2,\varepsilon=\pm1,c\ne0$。\par \textbf{操作步骤} 取$t=\frac{x}{y}$，用约束微分计算$\dd t$；把$y^n$消成$\frac{c}{1-\varepsilon t^n}$；所需积分归到有理函数族$\mathcal U$。\par \textbf{疑难点与条件} $c=0$退化为幂函数关系；不能代入含$\frac{1}{c}$的公式。\end{method}

\begin{theorem}{高次双元的基本式与降次式}在可微实支上，$\dd(\frac{x}{y})=\frac{c\,\dd x}{y^{n+1}}$，$\int\frac{\dd x}{y}=\int\frac{\dd t}{1-\varepsilon t^n}$。另记$J_m=\int y^m\dd x$，$m\ne-1$时$J_m=\frac{xy^m}{m+1}+\frac{mc\,J_{m-n}}{m+1}$。\end{theorem}\begin{proof}由$y'=\frac{\varepsilon x^{n-1}}{y^{n-1}}$得商导数$\frac{c}{y^{n+1}}$，且$y^n(1-\varepsilon t^n)=c$。再算$(xy^m)'=y^m+m\varepsilon x^ny^{m-n}=(m+1)y^m-mc\,y^{m-n}$，移项得到递推。\end{proof}

\begin{example}{竞赛风格｜四次双元的初等化}{难度：$\star\star\star$}\label{ic:ex:indefinite_dual_complete:fourthdual}\textbf{题目} 求$\int\frac{\dd x}{(1+x^4)^{\frac{5}{4}}}$及$\int\frac{\dd x}{(1+x^4)^{\frac{1}{4}}}$。\par \textbf{思路} 两者用同一$t=\frac{x}{(1+x^4)^{\frac{1}{4}}}$，但微分系数不同。\par \textbf{详细解答} 取$y=(1+x^4)^{\frac{1}{4}}>0$，第一题由$\dd t=\frac{\dd x}{y^5}$直接得$t+C$。第二题化为$\int\frac{\dd t}{1-t^4}$；部分分式为$\frac12(1-t^2)^{-1}+\frac12(1+t^2)^{-1}$，答案$\frac14\ln|\frac{1+t}{1-t}|+\frac12\arctan t+C$。此处$|t|<1$。\par \textbf{求导验算} 第二式对$t$导数为$\frac{1}{1-t^4}$，再乘$\frac{1}{y^5}$，因$1-t^4=\frac{1}{y^4}$得$\frac{1}{y}$。\par \textbf{疑难点与方法比较} 高次根式也可初等；关键是指数与约束的匹配。\end{example}

\begin{example}{竞赛风格｜高次双元递推的系数}{难度：$\star\star\star$}\label{ic:ex:indefinite_dual_complete:cubicrecur}\textbf{题目} 求$\int(x^3+2)^{\frac{2}{3}}\dd x$，取$x>0$。\par \textbf{思路} $n=3,c=2,m=2$，降到$J_{-1}$。\par \textbf{详细解答} 令$y=(x^3+2)^{\frac{1}{3}},t=\frac{x}{y}$。递推给$J_2=\frac{xy^2}{3}+\frac{4J_{-1}}{3}$。而$J_{-1}=\int\frac{\dd t}{1-t^3}=-\frac13\ln(1-t)+\frac16\ln(t^2+t+1)+\frac1{\sqrt3}\arctan(\frac{2t+1}{\sqrt3})$。代入即为完整原函数。\par \textbf{求导验算} 由已证递推和$t'=\frac{2}{y^4}$核验；$1-t^3=\frac{2}{y^3}$消去全部余项。\par \textbf{疑难点与方法比较} 递推中的$mc$不能漏掉；各项系数应由求导核验，避免幂次和下标混淆。\end{example}

\subsection{统一有理积分族与公因数降阶}\label{ic:indefinite_dual_complete:ufamily}\begin{method}{统一有理积分族与公因数降阶}\textbf{识别信号} $\mathcal U^\varepsilon_{m,n}(x)=\int \frac{x^{m-1}}{1+\varepsilon x^n}\dd x$，$m,n$为正整数。\par \textbf{操作步骤} $m\ge n$先长除；有公因数先$x^k$降阶；剩余真分式按单位根配共轭，转成实对数及反正切。\par \textbf{疑难点与条件} 降阶有$\frac{1}{k}$系数；该记号表示可计算有理积分，不是新的非初等函数。\end{method}

\begin{theorem}{公因数降阶的比例因子}$\mathcal U^\varepsilon_{km,kn}(x)=k^{-1}\mathcal U^\varepsilon_{m,n}(x^k)+C$。若$P=1+\varepsilon x^n$有互异复根$r_j$，则$\frac{x^{m-1}}{P}=\sum \frac{r_j^{m-1}}{P'(r_j)(x-r_j)}$（$m\le n$）。\end{theorem}\begin{proof}第一式令$u=x^k$，$x^{km-1}\dd x=\frac{u^{m-1}\dd u}{k}$。第二式由各简单极点的系数及真分式唯一性得到；实系数下共轭项合并为实表达。\end{proof}

\begin{example}{考研风格｜公因数与长除联用}{难度：$\star\star$}\label{ic:ex:indefinite_dual_complete:gcd}\textbf{题目} 求$\int \frac{x^5\dd x}{1+x^4}$。\par \textbf{思路} 先$x^2=t$，再除法。\par \textbf{详细解答} $x^5\dd x=\frac{t^2\dd t}{2}$，积分$\frac12\int[1-\frac{1}{1+t^2}]\dd t=\frac{x^2}2-\frac12\arctan(x^2)+C$。也可$\frac{x^5}{1+x^4}=x-\frac{x}{1+x^4}$后用$x^2$凑微分。\par \textbf{求导验算} 导数$x-\frac{x}{1+x^4}=\frac{x^5}{1+x^4}$。\par \textbf{疑难点与方法比较} 公因数降阶必须同时变换分子幂和微分。\end{example}

\subsection{隐函数：参数化与约束微分}\label{ic:indefinite_dual_complete:implicit}\begin{method}{隐函数：参数化与约束微分}\textbf{识别信号} $F(x,y)=0$且所取支路可微。\par \textbf{操作步骤} 先找$t=x-y,\frac{y}{x}$等使曲线可有理参数化的组合；或由$\dd F=0$把原微分改成$\dd y$，利用分部和原约束消元。\par \textbf{疑难点与条件} 若$F_y=0$，不能自动视为可微$y(x)$；原被积函数的分母零点也须排除。\end{method}

\begin{example}{竞赛风格｜同一隐函数积分的两解}{难度：$\star\star\star$}\label{ic:ex:indefinite_dual_complete:implicitcurve}\textbf{题目} 若$y(x-y)^2=x$，求$\int\frac{\dd x}{x-3y}$，在正则且分母非零的支路上。\par \textbf{思路} 参数$t=x-y$能同时解出$x,y$。\par \textbf{详细解答} 由$y(t^2-1)=t$得$y=\frac{t}{t^2-1},x=\frac{t^3}{t^2-1}$，因此$\frac{\dd x}{x-3y}=\frac{t\dd t}{t^2-1}$，原函数$\frac12\ln|(x-y)^2-1|+C$。第二式由约束可得$\frac14[\ln|x|-3\ln|y|]+C$，因$\frac{x}{y^3}=(t^2-1)^2$与前式相同。\par \textbf{求导验算} 参数导数$x'=\frac{t^2(t^2-3)}{(t^2-1)^2}$与$x-3y=\frac{t(t^2-3)}{t^2-1}$约分验证。\par \textbf{疑难点与方法比较} 参数式的$t=0,\pm1,\pm\sqrt3$需按原曲线与被积式逐一检查，不能凭约分把奇点补成同一道题。\end{example}

\begin{example}{竞赛风格｜隐式三次关系的两种表示}{难度：$\star\star\star$}\label{ic:ex:indefinite_dual_complete:implicitcubic}\textbf{题目} $z^3+xz=8$定义可微支$z(x)$，求$\int z^2\dd x$。\par \textbf{思路} 把$x$用$z$表示，或构造含$xz^2$的分部关系。\par \textbf{详细解答} 因$z\ne0$，$x=\frac{8}{z}-z^2$，$\dd x=(-\frac{8}{z^2}-2z)\dd z$，故积分$=-8z-\frac{z^4}{2}+C$。等价形式$-xz^2-\frac32z^4+C$由原约束相乘得到。\par \textbf{求导验算} 由$(3z^2+x)z'=-z$验证第一式导数为$z^2$；只在$3z^2+x\ne0$的正则段使用。\par \textbf{疑难点与方法比较} 不要把隐式变量当常数；先变参数可省去复杂根式解。\end{example}

\section{复杂构型、锁型、伪椭圆与连分式算法}\label{ic:indefinite_algebra_complete}

\begin{examtip}本节将隐藏的微分结构称为锁型，并用复合微分、商导数、线性微分方程和代数恒等式解释。连分式部分给出有限步骤及最后的多项式验算。\end{examtip}

\subsection{二项微分的三条有理化路线}\label{ic:indefinite_algebra_complete:binomial}\begin{method}{二项微分的三条有理化路线}\textbf{识别信号} $\int x^m(a+bx^n)^p\dd x$，指数有理，$n\ne0$；先在$x>0,a>0,a+bx^n>0$的实支讨论。\par \textbf{操作步骤} 若$p$为整数，用共同根次有理化；若$q=\frac{m+1}{n}$为整数，令$u=a+bx^n$后取分数幂公分母；若$p+q$为整数，令$u=b+ax^{-n}$后同样消去分数幂。\par \textbf{疑难点与条件} 这里证明的是三条可执行的充分路线；条件均不满足时，不凭这一张卡就断言非初等。\end{method}

\begin{theorem}{第三条路线的变换公式}令$q=\frac{m+1}{n}$、$k=p+q$。换$u=b+ax^{-n}$后，$x^m(a+bx^n)^p\dd x=-\frac{a^k u^p(u-b)^{-k-1}\dd u}{n}$。若$k\in\mathbb Z$，再令$u=t^d$，$d$为$p$的分母，就得到有理式。\end{theorem}\begin{proof}写$x^n=\frac{a}{u-b}$，于是$x^m\dd x=-(\frac{a^q}{n})(u-b)^{-q-1}\dd u$；又$(a+bx^n)^p=a^pu^p(u-b)^{-p}$，相乘即得。第二条路线类似：令$u=a+bx^n$后，积分是常数倍$(u-a)^{q-1}u^p\dd u$，$q$整数时取公分母即可有理化。\end{proof}

\begin{example}{竞赛风格｜整数$q$的非平方根积分}{难度：$\star\star\star$}\label{ic:ex:indefinite_algebra_complete:binomialq}\textbf{题目} 求$\int x^5(1+x^3)^{\frac{1}{3}}\dd x$，$x>0$。\par \textbf{思路} $q=\frac{5+1}{3}=2$，先取$u=1+x^3$。\par \textbf{详细解答} $x^5\dd x=\frac{(u-1)\dd u}{3}$，所以积分$\frac13\int(u^{\frac{4}{3}}-u^{\frac{1}{3}})\dd u=\frac{u^{\frac{7}{3}}}{7}-\frac{u^{\frac{4}{3}}}{4}+C$。也可直接设$t=(1+x^3)^{\frac{1}{3}}$，积分化为$(t^6-t^3)\dd t$。\par \textbf{求导验算} 求导为$x^2[(1+x^3)^{\frac{4}{3}}-(1+x^3)^{\frac{1}{3}}]=x^5(1+x^3)^{\frac{1}{3}}$。\par \textbf{疑难点与方法比较} 不必看到立方根就引入特殊函数。\end{example}

\begin{example}{竞赛风格｜整数$p+q$路线双解}{难度：$\star\star\star$}\label{ic:ex:indefinite_algebra_complete:binomialsum}\textbf{题目} 求$\int\frac{\sqrt{1+x^2}\dd x}{x^2}$，$x>0$。\par \textbf{思路} $p=\frac{1}{2},q=-\frac{1}{2}$，其和为0。\par \textbf{详细解答} 令$t=\sqrt{1+x^{-2}}=\frac{\sqrt{1+x^2}}{x}>1$，$\dd t=-\frac{\dd x}{x^2\sqrt{1+x^2}}$，原式$=-\int \frac{t^2}{t^2-1}\dd t=-t+\frac12\ln|\frac{t+1}{t-1}|+C$。另一法分部取$u=\sqrt{1+x^2},\dd v=x^{-2}\dd x$，得$-\frac{\sqrt{1+x^2}}{x}+\ln(x+\sqrt{1+x^2})+C$。\par \textbf{求导验算} 第二式求导为$\frac{1}{\sqrt{1+x^2}}+\frac{1}{x^2\sqrt{1+x^2}}=\frac{\sqrt{1+x^2}}{x^2}$。\par \textbf{疑难点与方法比较} 二项路线给系统换元，分部法在低次时更短。\end{example}

\subsection{Hermite约化：模多项式算法}\label{ic:indefinite_algebra_complete:hermitemod}\begin{method}{Hermite约化：模多项式算法}\textbf{识别信号} $\frac{A}{V^kW}$，$k\ge2$，$V$平方自由且与$W$互素。\par \textbf{操作步骤} 求次数小于$\deg V$的$B$，使$-(k-1)WV'B\equiv A\pmod V$；由扩展Euclid求逆。令$C=\frac{A-WVB'+(k-1)WV'B}{V}$，减去$(\frac{B}{V^{k-1}})'$，分母幂数降1。重复直到平方自由。\par \textbf{疑难点与条件} 模逆存在依赖$\gcd(V,WV')=1$；不满足时先因式分解及平方自由分解。\end{method}

\begin{theorem}{Hermite单步约化的证明}按上法定义$B,C$，则$\frac A{V^kW}=\left(\frac B{V^{k-1}}\right)'+\frac C{V^{k-1}W}$。\end{theorem}\begin{proof}右侧通分后分子为$WVB'-(k-1)WV'B+CV=A$。模同余使$C$确为多项式，分母中$V$的次数下降，因而重因子部分有限步终止。最后的平方自由真分式通过实或复根分解产生对数与反正切。\end{proof}

\begin{example}{竞赛风格｜一次完整模法约化}{难度：$\star\star\star$}\label{ic:ex:indefinite_algebra_complete:hermitefull}\textbf{题目} 求$\int\frac{\dd x}{(x^2+1)^2(x+1)}$。\par \textbf{思路} $V=x^2+1,W=x+1,k=2$。\par \textbf{详细解答} 模$V$有$-WV'\equiv2-2x$，其逆为$\frac{1+x}{4}$，故$B=\frac{1+x}{4}$。计算$C=\frac{x+3}{4}$，于是原式导数部分为$[\frac{x+1}{4V}]'$，剩余$\frac{x+3}{4V(x+1)}=\frac{1}{4(x+1)}+\frac{-x+2}{4V}$。积分得$\frac{x+1}{4(x^2+1)}+\frac14\ln|x+1|-\frac18\ln(x^2+1)+\frac12\arctan x+C_0$。\par \textbf{求导验算} 将原函数求导通分到$4V^2(x+1)$，分子化为4。\par \textbf{疑难点与方法比较} 本法无需一次设全部重因子的系数；只有模逆和低次部分分式。\end{example}

\subsection{有理积分的留数式系数法}\label{ic:indefinite_algebra_complete:residue}\begin{method}{有理积分的留数式系数法}\textbf{识别信号} 分母根已知，或只想先分离某几个极点。\par \textbf{操作步骤} 简单根$a$的系数为$\frac{N(a)}{D'(a)}$；$k$重根展开$\frac{(x-a)^kN}{D}$的前$k$项Taylor系数；共轭根成对还原实式。\par \textbf{疑难点与条件} 这里只计算局部部分分式系数，不涉及围道；和定积分的留数定理用途应区分。\end{method}

\begin{example}{竞赛风格｜重极点的局部展开}{难度：$\star\star\star$}\label{ic:ex:indefinite_algebra_complete:localcoeff}\textbf{题目} 分解并积分$\frac{1}{(x-1)^2(x+2)}$。\par \textbf{思路} 先在$x=1$展开剩余因子$\frac{1}{x+2}$。\par \textbf{详细解答} 写$x=1+h$，$\frac{1}{x+2}=\frac{1}{3}-\frac{h}{9}+O(h^2)$，故重极点贡献$\frac{1}{3(x-1)^2}-\frac{1}{9(x-1)}$。$x=-2$的简单项系数$\frac{1}{9}$，完整分解再加$\frac{1}{9(x+2)}$。原函数$-\frac{1}{3(x-1)}-\frac{\ln|x-1|}{9}+\frac{\ln|x+2|}{9}+C$。\par \textbf{求导验算} 求导后按上述完整分解通分为原式。\par \textbf{疑难点与方法比较} 局部展开避免大型待定系数方程，仍须检验全局真分式无余项。\end{example}

\subsection{平方分母的商导数降阶}\label{ic:indefinite_algebra_complete:quotient}\begin{method}{平方分母的商导数降阶}\textbf{识别信号} $\int \frac{f}{g^2}$且$\frac{f}{g'}$比原式简单。\par \textbf{操作步骤} 用$\dd(-\frac{1}{g})=\frac{g'\dd x}{g^2}$分部，得到$\int \frac{f}{g^2}=-\frac{f}{gg'}+\int\frac{f'g'-fg''}{g(g')^2}\dd x$；若$f=u'v-uv'$且$g=v$，直接识别商导数。\par \textbf{疑难点与条件} 须在$g,g'\ne0$的区间使用该变形；它是分部恒等式，不是极限的L’Hôpital法则。\end{method}

\begin{example}{竞赛风格｜藏在平方分母中的商}{难度：$\star\star\star$}\label{ic:ex:indefinite_algebra_complete:wronskian}\textbf{题目} 求$\int\frac{[x\cos x-\sin x+x^2\sin x]\dd x}{(x\cos x-\sin x)^2}$。\par \textbf{思路} 分子是$x$与分母的商导数交叉乘积。\par \textbf{详细解答} 设$g=x\cos x-\sin x$，则$g'=-x\sin x$，分子恰为$g-xg'$，因此原函数为$\frac{x}{g}+C=\frac{x}{x\cos x-\sin x}+C$。另一法把被积式拆成$\frac{1}{g}-\frac{xg'}{g^2}$，用分部积分$\int x\dd(\frac{1}{g})=\frac{x}{g}-\int\frac{\dd x}{g}$，两个剩余积分相消。\par \textbf{求导验算} $[\frac{x}{g}]'=\frac{g-xg'}{g^2}$，直接恢复原式。\par \textbf{疑难点与方法比较} 只在$g\ne0$的连通区间写同一个常数；商导数的交叉项符号必须按顺序核对。\end{example}

\subsection{反三角混合锁与有限维闭合}\label{ic:indefinite_algebra_complete:triglock}\begin{method}{反三角混合锁与有限维闭合}\textbf{识别信号} 指数的内层是$\arctan x$，分母含$(1+x^2)$的半整数幂。\par \textbf{操作步骤} 先令$t=\arctan x$，根式变为$\sec t$，再检查是否化指数三角积分；或以双元分部构造两个积分的线性关系。\par \textbf{疑难点与条件} 选择$t\in(-\frac{\pi}{2},\frac{\pi}{2})$保证$\sec t>0$；不要直接把平方根换成带任意符号的正割。\end{method}

\begin{example}{竞赛风格｜混合锁型的两解}{难度：$\star\star\star$}\label{ic:ex:indefinite_algebra_complete:atanlock}\textbf{题目} 求$\int \frac{x\e^{\arctan x}\dd x}{(1+x^2)^{\frac{3}{2}}}$。\par \textbf{思路} 换成角变量后是$\e^t\sin t$。\par \textbf{详细解答} $t=\arctan x$给$\int\e^t\sin t\dd t=\frac{\e^t(\sin t-\cos t)}{2}+C$，回代为$\frac{\e^{\arctan x}(x-1)}{2\sqrt{1+x^2}}+C$。另一法令$v=\sqrt{1+x^2}$，对$\frac{\e^{\arctan x}}{v}$分部；再对$\frac{x}{v}$用$(\frac{x}{v})'=\frac{1}{v^3}$，两式联立可解出同结果。\par \textbf{求导验算} 导数分子为$(x-1)+(1+x)$，除以$2v^3$后为$\frac{x}{v^3}$并乘$\e^{\arctan x}$。\par \textbf{疑难点与方法比较} 混合锁的难点是先看见角变量，而不是展开复杂高阶导数。\end{example}

\subsection{指数锁与归一化变量}\label{ic:indefinite_algebra_complete:explock}\begin{method}{指数锁与归一化变量}\textbf{识别信号} 根式中有$\e^x-x^2$，或分母为指数与多项式之和。\par \textbf{操作步骤} 先提指数的正因子，再检查$x\e^{-\frac{x}{2}}$的微分；对$g=\e^x+P$使用$g'-g=P'-P$将多项式与指数拆开。\par \textbf{疑难点与条件} 先写实定义域；指数提根式没有符号问题，但剩余根式仍须非负。\end{method}

\begin{example}{竞赛风格｜指数根式锁}{难度：$\star\star\star$}\label{ic:ex:indefinite_algebra_complete:exprootlock}\textbf{题目} 求$\int\frac{(x-2)\dd x}{\sqrt{\e^x-x^2}}$，在$\e^x>x^2$的区间。\par \textbf{思路} 提$\e^{\frac{x}{2}}$后，$t=x\e^{-\frac{x}{2}}$的导数完整出现。\par \textbf{详细解答} $\sqrt{\e^x-x^2}=\e^{\frac{x}{2}}\sqrt{1-t^2}$，$\dd t=-\frac12(x-2)\e^{-\frac{x}{2}}\dd x$，故答案$-2\arcsin(x\e^{-\frac{x}{2}})+C$。\par \textbf{求导验算} 链式求导给$\frac{(x-2)\e^{-\frac{x}{2}}}{\sqrt{1-x^2\e^{-x}}}$，化回原式。\par \textbf{疑难点与方法比较} 看似含指数和代数根式的混合题，可以由一个归一化变量解锁。\end{example}

\begin{example}{竞赛风格｜指数平方分母的闭合}{难度：$\star\star\star$}\label{ic:ex:indefinite_algebra_complete:expsquarelock}\textbf{题目} 求$\int \frac{x(x+1)\dd x}{(\e^x+x+1)^2}$。\par \textbf{思路} 记$g=\e^x+x+1$，有$g-g'=x$。\par \textbf{详细解答} $[\frac{x+1}{g}]'=-\frac{x\e^x}{g^2}$，而$\frac{x(x+1)}{g^2}=\frac{x}{g}-\frac{x\e^x}{g^2}$。又$\frac{x}{g}=1-\frac{g'}{g}$，故原函数$x-\ln|g|+\frac{x+1}{g}+C$。\par \textbf{求导验算} 导数$1-\frac{g'}{g}-\frac{x\e^x}{g^2}=\frac{x(g-\e^x)}{g^2}=\frac{x(x+1)}{g^2}$。\par \textbf{疑难点与方法比较} 先利用$g-g'$消掉指数，比对平方分母硬设全部系数更有针对性。\end{example}

\subsection{等分锁与对数微分配对}\label{ic:indefinite_algebra_complete:equalparts}\begin{method}{等分锁与对数微分配对}\textbf{识别信号} 分子是两个量的相对增长率之和，如$\frac{f'}{f}+\frac{g'}{g}$。\par \textbf{操作步骤} 先设$t=fg$或$\frac{f}{g}$，把分子看成$\dd\ln|t|$；再将分母中的同一$t$替换。\par \textbf{疑难点与条件} 相对微分只在非零区间成立；最终答案仍需检查原式定义域。\end{method}

\begin{example}{竞赛风格｜等分性转化为一元有理积分}{难度：$\star\star\star$}\label{ic:ex:indefinite_algebra_complete:equalpartlock}\textbf{题目} 求$\int\frac{(1+\cot x)\dd x}{1+\e^x\sin x}$。\par \textbf{思路} $t=\e^x\sin x$满足$\frac{\dd t}{t}=(1+\cot x)\dd x$。\par \textbf{详细解答} 积分$\int\frac{\dd t}{t(1+t)}=\ln|t|-\ln|1+t|+C=x+\ln|\sin x|-\ln|1+\e^x\sin x|+C$。\par \textbf{求导验算} 对两对数求导，公因子$1+\cot x$乘$1-\frac{t}{1+t}$，得到原式。\par \textbf{疑难点与方法比较} 等分性是对数链式法则；不需要把它当作没有条件的新积分规则。\end{example}

\subsection{代数原函数待定与伪椭圆判别}\label{ic:indefinite_algebra_complete:exactalg}\begin{method}{代数原函数待定与伪椭圆判别}\textbf{识别信号} 被积式$\frac{R(x)}{\sqrt{D(x)}}$，但有可能是$A(x)\sqrt D$或对数的导数。\par \textbf{操作步骤} 先设有理函数$A$，解$DA'+\frac{D'A}{2}=R$；对极点和无穷处比较次数确定候选形式；剩余项再试Pell型对数或真正椭圆标准形。\par \textbf{疑难点与条件} 某个候选失败不能证明没有初等原函数；伪椭圆指可初等化的具体微分，不是只看根式次数。\end{method}

\begin{example}{竞赛风格｜代数恰当微分的完整配系数}{难度：$\star\star\star$}\label{ic:ex:indefinite_algebra_complete:algebraansatz}\textbf{题目} 求$\int\frac{(1-x^4)\dd x}{(x^4+x^2+1)^{\frac{3}{2}}}$。\par \textbf{思路} 分母三次半幂提示原函数为低次多项式除以$\sqrt D$。\par \textbf{详细解答} 设$F=\frac{ax+b}{\sqrt D}$，则$F'=\frac{aD-\frac{(ax+b)D'}{2}}{D^{\frac{3}{2}}}$。比较奇偶项得$b=0$，再比较常数与$x^4$得$a=1$，故$F=\frac{x}{\sqrt{x^4+x^2+1}}+C$。\par \textbf{求导验算} 分子$D-\frac{xD'}{2}=1-x^4$。\par \textbf{疑难点与方法比较} 先用奇偶性和次数把待定式限制到最小，再解系数；不随意猜高次函数。\end{example}

\subsection{四次根式的多项式连分式算法}\label{ic:indefinite_algebra_complete:continuedfraction}\begin{method}{四次根式的多项式连分式算法}\textbf{识别信号} $D$为首项系数1的四次多项式；在无穷远作形式Laurent展开寻找对数型原函数。\par \textbf{操作步骤} 令$P_0=0,Q_0=1$，$a_k$取$\frac{\sqrt D+P_k}{Q_k}$的多项式部分。递推$P_{k+1}=a_kQ_k-P_k$、$Q_{k+1}=\frac{D-P_{k+1}^2}{Q_k}$。收敛分式$\frac{p_k}{q_k}$按$p_k=a_kp_{k-1}+p_{k-2}$、$q_k=a_kq_{k-1}+q_{k-2}$计算。每步检验$p_k^2-Dq_k^2$是否为非零常数。\par \textbf{疑难点与条件} 初值$p_{-2}=0,p_{-1}=1,q_{-2}=1,q_{-1}=0$。有限步未找到恒等式时，不能凭此宣布不可积；形式展开不依赖原函数的实定义区间。\end{method}

\begin{theorem}{Pell恒等式的最终证书}设$P,Q,D$为实多项式，在所用区间有$D>0,Q\ne0$。若$P^2-DQ^2=c\ne0$，则$[\ln|P+Q\sqrt D|]'=\frac{P'}{Q\sqrt D}$。若进一步$\frac{P'}{Q}=\alpha x+\beta$，就得出该线性分子与四次根式的原函数。\end{theorem}\begin{proof}微分恒等式得$PP'=DQQ'+\frac{D'Q^2}{2}$。对数导数的分子$P'+Q'\sqrt D+\frac{QD'}{2\sqrt D}$等于$\frac{P'(P+Q\sqrt D)}{Q\sqrt D}$，约分得到结论。代数证书必须最后独立展开核对。\end{proof}

\begin{example}{竞赛拓展｜连分式积分的五步计算}{难度：$\star\star\star$}\label{ic:ex:indefinite_algebra_complete:continuedfull}\textbf{题目} 求$\int \frac{x\dd x}{\sqrt{(x+1)^4-12x^2}}$，在根式为正的区间。\par \textbf{思路} 用多项式连分式寻找常数Pell余项；本例可在第4号收敛分式找到。\par \textbf{详细解答} 记$D=(x+1)^4-12x^2$。按上述普通连分式递推，得到下表。\par \noindent\begin{tabular}{ccl}\toprule $k$&$a_k$&$p_k^2-Dq_k^2$\\\midrule 0&$x^2+2x-5$&$-24(x-1)$\\1&$\frac{x+3}{12}$&$-\frac{x}{3}$\\2&$-6(x+2)$&$-36x$\\3&$\frac{x+2}{18}$&$-\frac{2(x-1)}{9}$\\4&$-9(x+3)$&$108$\\\bottomrule\end{tabular}\par 将最后分式同时乘$-4$，得\begin{align*}P&=x^6+12x^5+45x^4+44x^3-33x^2+43,\\Q&=x^4+10x^3+30x^2+22x-11.\end{align*}独立展开验证$P^2-DQ^2=1728$、$P'=6xQ$，所以所求为$\frac16\ln|P+Q\sqrt D|+C$。\par \textbf{求导验算} Pell对数微分给$\frac{P'}{6Q\sqrt D}=\frac{x}{\sqrt D}$。若$Q$在某点为0，最后对数式可按原被积函数的正则性连续延拓，而中间商式仅在$Q\ne0$时使用。\par \textbf{疑难点与方法比较} 换元后的根式必须重新通分核验。本例以显示的$D$和两条多项式恒等式独立验证结果，所有幂次均由展开确定。\end{example}

\subsection{微分方程引领积分：Lagrange恒等式}\label{ic:indefinite_algebra_complete:odeidentity}\begin{method}{微分方程引领积分：Lagrange恒等式}\textbf{识别信号} 已知$y''+py'+qy=0$，被积函数与$y$及其导数有关。\par \textbf{操作步骤} 选简单$h$，计算$h''+ph'+qh$；将目标凑成这个组合乘$y\e^{\int p}$，再写出边界组合$h'y-hy'$。\par \textbf{疑难点与条件} $h$须有二阶导数，系数和解在同一区间正则；不是任意选$h$都能匹配目标。\end{method}

\begin{theorem}{二阶方程积分公式}若$y''+py'+qy=0$，则$\int \e^{\int p}(h''+ph'+qh)y\dd x=\e^{\int p}(h'y-hy')+C$。\end{theorem}\begin{proof}对右端求导，得到$\e^{\int p}[(h''+ph')y-h(y''+py')]$，用原方程$y''+py'=-qy$即为左端被积式。\end{proof}

\begin{example}{竞赛风格｜ODE恒等式与分部互证}{难度：$\star\star\star$}\label{ic:ex:indefinite_algebra_complete:ode_sine}\textbf{题目} 求$\int(x^2+2)\sin x\dd x$。\par \textbf{思路} $\sin x$满足$y''+y=0$，取$h=x^2$。\par \textbf{详细解答} $h''+h=x^2+2$，故原函数$h'y-hy'=2x\sin x-x^2\cos x+C$。两次分部也给$\int x^2\sin x=-x^2\cos x+2x\sin x+2\cos x$，加$2\int\sin x$后余项消去。\par \textbf{求导验算} 导数中的$2x\cos x$互相抵消，余下$(x^2+2)\sin x$。\par \textbf{疑难点与方法比较} ODE恒等式可先识别整组外力，避免逐项积分后再相消。\end{example}

\section{含参、级数、复数与非初等原函数的完整处理}\label{ic:indefinite_special_complete}

\begin{examtip}这里的“求出”包括给出定义明确、可求导验证的特殊函数表达。对于存在初等原函数的题，先用前述方法约化；特殊函数不是掩盖未完成计算的占位名称。\end{examtip}

\subsection{含参不定积分须固定积分常数}\label{ic:indefinite_special_complete:parameter}\begin{method}{含参不定积分须固定积分常数}\textbf{识别信号} 参数微分能产生对数幂，或消去复杂因子。\par \textbf{操作步骤} 用$F(x,a)=\int_{x_0}^xf(t,a)\dd t$归一化；在紧矩形上验证连续偏导；对参数求导；最后加一个与$x$无关的常数。\par \textbf{疑难点与条件} 未经归一化的原函数可多出$C(a)$，直接对参数求导会漏$C'(a)$。\end{method}

\begin{theorem}{归一化含参原函数}若$f,\partial_af$连续，则$\partial_a\int_{x_0}^xf(t,a)\dd t=\int_{x_0}^x\partial_af(t,a)\dd t$。对$a$多次求导的版本要求相应偏导连续。\end{theorem}\begin{proof}差商在紧积分区间上一致趋向偏导，故积分与极限可交换。下限$x_0$固定，使所有原函数在$x_0$值为0，不产生未知的参数函数。\end{proof}

\begin{example}{考研风格｜参数微分与退化参数}{难度：$\star\star\star$}\label{ic:ex:indefinite_special_complete:parampower}\textbf{题目} 求$\int x^a\ln x\dd x$，$x>0$，并处理$a=-1$。\par \textbf{思路} 从$\frac{x^{a+1}-1}{a+1}$的参数导数开始。\par \textbf{详细解答} $a\ne-1$时，归一化原函数为$\frac{x^{a+1}\ln x}{a+1}-\frac{x^{a+1}-1}{(a+1)^2}$。若只求普通原函数，可把末尾常数吸收成$\frac{x^{a+1}\ln x}{a+1}-\frac{x^{a+1}}{(a+1)^2}+C$。$a=-1$时令$t=\ln x$得$\frac{(\ln x)^2}{2}+C$，也等于归一化公式的参数极限。\par \textbf{求导验算} 对$x$求导，两项中的$\frac{x^a}{a+1}$相消，余下$x^a\ln x$。\par \textbf{疑难点与方法比较} 先固定常数再对参数求导，退化参数可通过极限或单独换元完成。\end{example}

\begin{example}{竞赛风格｜参数产生$x$因子的积分}{难度：$\star\star\star$}\label{ic:ex:indefinite_special_complete:paramtrig}\textbf{题目} 求$\int x\cos(ax)\dd x$。\par \textbf{思路} 对$\int\sin(ax)\dd x$的归一化式求参数导数。\par \textbf{详细解答} $a\ne0$时$\int_0^x\sin(at)\dd t=\frac{1-\cos ax}{a}$，参数微分得到$\int_0^xt\cos(at)\dd t=\frac{x\sin(ax)}{a}+\frac{\cos ax-1}{a^2}$。故普通原函数$\frac{x\sin(ax)}{a}+\frac{\cos(ax)}{a^2}+C$。$a=0$为$\frac{x^2}{2}+C$。直接分部也得到相同公式。\par \textbf{求导验算} 对$x$求导为$\frac{\sin(ax)}{a}+x\cos(ax)-\frac{\sin(ax)}{a}$。\par \textbf{疑难点与方法比较} 未归一化式中$a^{-2}$的发散可被参数常数抵消，不能据此误判$a=0$的原函数不存在。\end{example}

\subsection{幂级数逐项积分与命名函数}\label{ic:indefinite_special_complete:series}\begin{method}{幂级数逐项积分与命名函数}\textbf{识别信号} 已知幂级数在$|x|<R$收敛。\par \textbf{操作步骤} 先在紧子区间用一致收敛逐项积分；写出原函数的收敛半径；若识别为标准函数，给出定义和归一化；端点另判。\par \textbf{疑难点与条件} 数值拟合的级数不等于精确原函数；不能在收敛圆外直接沿用未延拓的级数。\end{method}

\begin{example}{竞赛风格｜反正切积分函数}{难度：$\star\star\star$}\label{ic:ex:indefinite_special_complete:invlogseries}\textbf{题目} 表达$\int\frac{\arctan x\dd x}{x}$，并给$|x|<1$的级数。\par \textbf{思路} 分子级数除以$x$后逐项积分。\par \textbf{详细解答} 定义\[\operatorname{Ti}_2(x)=\int_0^x\dfrac{\arctan t}{t}\dd t,\qquad \left.\dfrac{\arctan t}{t}\right|_{t=0}:=1.\]由反正切展开，在$|x|<1$有\[\arctan x=\sum_{n\ge0}\dfrac{(-1)^nx^{2n+1}}{2n+1},\qquad \operatorname{Ti}_2(x)=\sum_{n\ge0}\dfrac{(-1)^nx^{2n+1}}{(2n+1)^2}.\]所以答案$\operatorname{Ti}_2(x)+C$；在$x=1$该绝对收敛和是Catalan常数。\par \textbf{求导验算} 级数在紧子区间逐项求导恢复$\frac{\arctan x}{x}$，实轴上直接由积分定义验证。\par \textbf{疑难点与方法比较} 本题用精确定义完成表达，不强求不存在的简单初等形式。\end{example}

\begin{example}{考研风格｜有对数也可能完全初等}{难度：$\star\star\star$}\label{ic:ex:indefinite_special_complete:logelementary}\textbf{题目} 求$\int\frac{\ln(1-x)\dd x}{x^2}$，$0<x<1$。\par \textbf{思路} 分部后只剩有理积分。\par \textbf{详细解答} 取$u=\ln(1-x),\dd v=x^{-2}\dd x$，得$-\frac{\ln(1-x)}{x}-\int\frac{\dd x}{x(1-x)}$，故答案$-\frac{\ln(1-x)}{x}-\ln x+\ln(1-x)+C$。\par \textbf{求导验算} 导数中的$\frac{1}{x(1-x)}-\frac{1}{x}-\frac{1}{1-x}$抵消。\par \textbf{疑难点与方法比较} 是否需要多重对数取决于分部后出现的核，不取决于题目中是否有对数。\end{example}

\subsection{复指数与实原函数的还原}\label{ic:indefinite_special_complete:complex}\begin{method}{复指数与实原函数的还原}\textbf{识别信号} 指数三角乘积、三角有理式或高次有理分母。\par \textbf{操作步骤} 利用$\e^{ibx}=\cos bx+i\sin bx$合并；积分后取实部虚部；复对数按共轭成对消去虚常数。\par \textbf{疑难点与条件} 复对数分支变化会产生分段常数；只比较导数不能强迫跨奇点常数相同。\end{method}

\begin{example}{考研风格｜多项式乘复指数}{难度：$\star\star\star$}\label{ic:ex:indefinite_special_complete:complexpolynomial}\textbf{题目} 求$\int x\e^x\cos x\dd x$。\par \textbf{思路} 对$\int x\e^{(1+i)x}\dd x$一次分部。\par \textbf{详细解答} 复原函数$\e^{(1+i)x}[\frac{x}{1+i}-\frac{1}{(1+i)^2}]$。用$\frac{1}{1+i}=\frac{1-i}{2}$、$\frac{1}{(1+i)^2}=-\frac{i}{2}$，取实部得$\frac{\e^x}{2}[x\cos x+(x-1)\sin x]+C$。另一法先求$\int\e^x\cos x$再对$x$分部。\par \textbf{求导验算} 求导后正弦项系数为0，余弦项系数为$x$。\par \textbf{疑难点与方法比较} 复方法只是把两条实线性关系合成一条，不改变最终实常数的自由度。\end{example}

\subsection{Ei、li、Si、Ci的识别核}\label{ic:indefinite_special_complete:ei}\begin{method}{Ei、li、Si、Ci的识别核}\textbf{识别信号} 指数、正弦、余弦分别除以自变量；或$\ln x$在分母。\par \textbf{操作步骤} 固定$\operatorname{Si}(x)=\int_0^x\frac{\sin t\,\dd t}{t}$；$\operatorname{Ci}(x)=\gamma+\ln x+\int_0^x\frac{(\cos t-1)\dd t}{t}$（$x>0$）；实$\operatorname{Ei}(x)=\gamma+\ln|x|+\sum_{k\ge1}\frac{x^k}{k\,k!}$（$x\ne0$）。\par \textbf{疑难点与条件} $\operatorname{Ci}'=\frac{\cos x}{x}$，不是$-\frac{\sin x}{x}$；$\int_0^x\frac{\cos t}{t}$不作为普通收敛定义。\end{method}

\begin{example}{竞赛拓展｜指数积分的复合与分部}{难度：$\star\star\star$}\label{ic:ex:indefinite_special_complete:ei_scaling}\textbf{题目} 求$\int\frac{\e^{2x}\dd x}{x}$与$\int\e^x\ln x\dd x$（第二式$x>0$）。\par \textbf{思路} 第一式识别Ei，第二式分部。\par \textbf{详细解答} 第一式$\operatorname{Ei}(2x)+C$，$x$在不含0的区间。第二式取$u=\ln x,\dd v=\e^x\dd x$，得$\e^x\ln x-\operatorname{Ei}(x)+C$。\par \textbf{求导验算} $\operatorname{Ei}'(2x)\cdot2=\frac{\e^{2x}}{x}$；第二式导数中的$\frac{\e^x}{x}$抵消。\par \textbf{疑难点与方法比较} 链式法则的参数2与分母2x相消。\end{example}

\begin{example}{竞赛拓展｜对数积分的反向复合}{难度：$\star\star\star$}\label{ic:ex:indefinite_special_complete:li}\textbf{题目} 求$\int\frac{\dd x}{\ln x}$，$x>0,x\ne1$。\par \textbf{思路} 设$t=\ln x$，分子变成$\e^t$。\par \textbf{详细解答} 积分$\int\frac{\e^t\dd t}{t}=\operatorname{Ei}(\ln x)+C$。可将该实分支记为$\operatorname{li}(x)$，但跨$x=1$时它有对数奇点，两个区间的常数分别选择。\par \textbf{求导验算} 链式求导$\frac{\frac{\e^{\ln x}}{\ln x}\cdot1}{x}=\frac{1}{\ln x}$。\par \textbf{疑难点与方法比较} 不能从0跨过1直接写成普通收敛积分；若用主值必须明确标记。\end{example}

\subsection{Gaussian、erf、erfi与Fresnel}\label{ic:indefinite_special_complete:gauss}\begin{method}{Gaussian、erf、erfi与Fresnel}\textbf{识别信号} 平方指数或二次相位且无内层导数因子。\par \textbf{操作步骤} 缩放成标准核；多项式权重先分部递推，留下基本特殊函数；若分子已有$x$，优先直接平方变量换元。\par \textbf{疑难点与条件} $\operatorname{erfi}$与$\operatorname{erf}$的指数符号相反；Fresnel常见归一化有不同$\frac{\pi}{2}$因子。\end{method}

\begin{example}{竞赛拓展｜Gaussian二阶矩原函数}{难度：$\star\star\star$}\label{ic:ex:indefinite_special_complete:gaussianmoment}\textbf{题目} 求$\int x^2\e^{-x^2}\dd x$。\par \textbf{思路} 用$x\e^{-x^2}\dd x=-\frac{\dd(\e^{-x^2})}{2}$分部。\par \textbf{详细解答} 积分$=-\frac{x\e^{-x^2}}{2}+\frac12\int\e^{-x^2}\dd x=-\frac{x\e^{-x^2}}{2}+\frac{\sqrt\pi\,\operatorname{erf}x}{4}+C$，其中$\operatorname{erf}x=2\pi^{-\frac{1}{2}}\int_0^x\e^{-t^2}\dd t$。\par \textbf{求导验算} 导数$-\frac{\e^{-x^2}}{2}+x^2\e^{-x^2}+\frac{\e^{-x^2}}{2}$。\par \textbf{疑难点与方法比较} 先消掉可分部的多项式权重，再保留基本核；不必每次定义新函数。\end{example}

\begin{example}{竞赛拓展｜Fresnel归一化}{难度：$\star\star\star$}\label{ic:ex:indefinite_special_complete:fresnelscale}\textbf{题目} 求$\int\sin(3x^2)\dd x$。\par \textbf{思路} 将相位缩放为$\frac{\pi t^2}{2}$。\par \textbf{详细解答} 定义$S(z)=\int_0^z\sin(\frac{\pi t^2}{2})\dd t$，则原函数$\sqrt{\frac{\pi}{6}}\,S(\sqrt{\frac{6}{\pi}}\,x)+C$。如果题目改为$\int x\sin(3x^2)\dd x$，则为$-\frac{\cos(3x^2)}{6}+C$。\par \textbf{求导验算} 前式链式求导的两个缩放系数乘积为1，后式由内层导数$6x$核验。\par \textbf{疑难点与方法比较} 命名函数前必须给出采用的归一化。\end{example}

实变量还采用$\operatorname{erfi}(x)=\frac2{\sqrt\pi}\int_0^x\e^{t^2}\dd t$，所以$a>0$时$\int\e^{ax^2}\dd x=\frac{\sqrt\pi}{2\sqrt a}\operatorname{erfi}(\sqrt a\,x)+C$；负指数相应使用$\operatorname{erf}$。两式均由基本定理与链式法则验证。

\subsection{多重对数的逐阶降阶}\label{ic:indefinite_special_complete:polylog}\begin{method}{多重对数的逐阶降阶}\textbf{识别信号} 出现$\frac{\ln(1-x)}{x}$、$\frac{\operatorname{Li}_s(x)}{x}$或对数乘积。\par \textbf{操作步骤} 在$|x|<1$定义$\operatorname{Li}_s(x)=\sum_{n\ge1}\frac{x^n}{n^s}$，用$\operatorname{Li}_1=-\ln(1-x)$、$\operatorname{Li}_s'=\frac{\operatorname{Li}_{s-1}}{x}$；分部使对数幂下降。\par \textbf{疑难点与条件} 超出单位圆或跨分支切线时需延拓约定；不能把复数恒等式不加条件地当成实式。\end{method}

\begin{example}{竞赛拓展｜三重对数完成对数平方积分}{难度：$\star\star\star$}\label{ic:ex:indefinite_special_complete:polylog3}\textbf{题目} 求$\int\frac{(\ln x)^2\dd x}{1-x}$，$0<x<1$。\par \textbf{思路} 先用$\ln(1-x)$作原函数，再分部两次。\par \textbf{详细解答} 第一次得$-(\ln x)^2\ln(1-x)+2\int\frac{\ln x\ln(1-x)\dd x}{x}$。利用$\frac{\ln(1-x)\dd x}{x}=-\dd\operatorname{Li}_2(x)$，剩余式$=-2\ln x\operatorname{Li}_2(x)+2\operatorname{Li}_3(x)$。总和加$C$即答案。\par \textbf{求导验算} 三项导数中的$\frac{\ln x\ln(1-x)}{x}$和$\frac{\operatorname{Li}_2(x)}{x}$分别成对抵消，余下原式。\par \textbf{疑难点与方法比较} 分部沿$\operatorname{Li}_2\to\operatorname{Li}_3$构成可追踪的层次，不是任意命名。\end{example}

\subsection{Lambert W：用隐式关系换回单元}\label{ic:indefinite_special_complete:lambert}\begin{method}{Lambert W：用隐式关系换回单元}\textbf{识别信号} 出现$W(x)$与$x$的乘除，且能用$x=u\e^u$替换。\par \textbf{操作步骤} 设$u=W(x)$，使用$\dd x=(1+u)\e^u\dd u$；计算后保留所取实分支。\par \textbf{疑难点与条件} 分支点$-\frac{1}{\e}$处$W'$不有限；实主支与下支不能混为同一可微函数。\end{method}

\begin{example}{竞赛拓展｜Lambert W两个基本核}{难度：$\star\star\star$}\label{ic:ex:indefinite_special_complete:lambertpair}\textbf{题目} 求$\int \frac{W(x)\dd x}{x}$与$\int\frac{\dd x}{1+W(x)}$，在相应可微支路。\par \textbf{思路} 统一设$x=u\e^u$。\par \textbf{详细解答} 第一式$\int(1+u)\dd u=u+\frac{u^2}{2}+C=W(x)+\frac{W(x)^2}{2}+C$。第二式$\int\e^u\dd u=\e^{W(x)}+C$；若写成$\frac{x}{W(x)}$，$x=0$处须取极限1。\par \textbf{求导验算} 利用$W'=\frac{W}{x(1+W)}$分别验证；在$W=0$处用连续延拓而非代入$\frac{0}{0}$。\par \textbf{疑难点与方法比较} 隐函数积分与特殊函数积分用同一参数法。\end{example}

\subsection{不完全Beta与Gamma的统一幂权}\label{ic:indefinite_special_complete:incomplete}\begin{method}{不完全Beta与Gamma的统一幂权}\textbf{识别信号} 幂函数乘指数，或$x$与$1+x$的分数幂。\par \textbf{操作步骤} 指数权令$t=ax^n$归到$\gamma(s,z)=\int_0^zt^{s-1}\e^{-t}\dd t$；分式幂令$t=\frac{x}{1+x}$归到$B_z(p,q)$。\par \textbf{疑难点与条件} 先保证定义区间；完全积分的端点参数条件通常强于不完全积分。\end{method}

\begin{example}{竞赛拓展｜指数幂权的完整公式}{难度：$\star\star\star$}\label{ic:ex:indefinite_special_complete:incompletegamma}\textbf{题目} 求$\int x^m\e^{-ax^n}\dd x$，$a,n>0,m>-1,x>0$。\par \textbf{思路} 平方换元的推广$t=ax^n$。\par \textbf{详细解答} 设$s=\frac{m+1}{n}>0$，则原函数$\frac{\gamma(s,ax^n)}{n a^s}+C$。因为$x^m\dd x=\frac{a^{-s}t^{s-1}\dd t}{n}$，完整变换直接得到。若$s$为正整数，$\gamma(s,z)$可用分部递推化为指数乘有限多项式；$s=\frac{1}{2}$化为erf。\par \textbf{求导验算} $\gamma_z(s,z)=z^{s-1}\e^{-z}$，链式乘$anx^{n-1}$后所有幂与常数还原为$x^m\e^{-ax^n}$。\par \textbf{疑难点与方法比较} 这个公式统一许多核，不意味着优先跳过可用的初等换元。\end{example}

这里$B_z(p,q)=\int_0^z t^{p-1}(1-t)^{q-1}\dd t$，$p>0,0<z<1$时允许任意实$q$。例如$x>0,p>0$时，$t=\frac{x}{1+x}$给\[\int x^{p-1}(1+x)^{-p-q}\dd x=B_{\frac{x}{1+x}}(p,q)+C,\]因为$\dd x=\frac{\dd t}{(1-t)^2}$，所有幂化为$t^{p-1}(1-t)^{q-1}$。上限若到1，另要求$q>0$才能得到完全Beta积分。

\subsection{三类椭圆标准化与实参数选择}\label{ic:indefinite_special_complete:elliptic}\begin{method}{三类椭圆标准化与实参数选择}\textbf{识别信号} 根式含两个不同二次因子的乘积，简单代换不能消成有理式。\par \textbf{操作步骤} 先移位缩放把分支点化为标准位置，再$x=a\sin\phi$；分别识别$F,E,\Pi$；退化参数另检。\par \textbf{疑难点与条件} 第一、第二、第三类分别有不同的核，不能全部笼统记作一个椭圆函数。\end{method}

\begin{example}{竞赛拓展｜两个二次因子的三类标准式}{难度：$\star\star\star$}\label{ic:ex:indefinite_special_complete:ellipticthree}\textbf{题目} 设$0<a<b,|x|<a$，将三个积分化为标准形：$I_1=\int\frac{\dd x}{\sqrt{(a^2-x^2)(b^2-x^2)}}$，$I_2=\int\frac{\sqrt{b^2-x^2}\dd x}{\sqrt{a^2-x^2}}$，$I_3$为$I_1$另乘$(1-\frac{\nu x^2}{a^2})^{-1}$。\par \textbf{思路} 令$x=a\sin\phi,k=\frac{a}{b}$。\par \textbf{详细解答} $\phi=\arcsin(\frac{x}{a})$，则$I_1=b^{-1}F(\phi,k)+C$、$I_2=bE(\phi,k)+C$、$I_3=b^{-1}\Pi(\nu;\phi,k)+C$。这里$F_\phi=(1-k^2\sin^2\phi)^{-\frac{1}{2}}$，$E_\phi=\sqrt{1-k^2\sin^2\phi}$，$\Pi_\phi=[(1-\nu\sin^2\phi)\sqrt{1-k^2\sin^2\phi}]^{-1}$。\par \textbf{求导验算} 三式分别乘$\phi'=\frac{1}{\sqrt{a^2-x^2}}$并代$\sqrt{1-k^2\sin^2\phi}=\frac{\sqrt{b^2-x^2}}{b}$，精确还原三核。\par \textbf{疑难点与方法比较} 第三类需避开$1-\frac{\nu x^2}{a^2}=0$；$k=0,1$的退化不能无条件当作一般情形。\end{example}

\begin{example}{竞赛拓展｜四次根式化实模数}{难度：$\star\star\star$}\label{ic:ex:indefinite_special_complete:quarticelliptic}\textbf{题目} 表达$\int\frac{\dd x}{\sqrt{1-x^4}}$，$|x|<1$。\par \textbf{思路} 选择$\sin\phi=\frac{\sqrt2x}{\sqrt{1+x^2}}$使模数为实数。\par \textbf{详细解答} 取$\phi=\arcsin(\frac{\sqrt2x}{\sqrt{1+x^2}})$，则原函数$\frac1{\sqrt2}F(\phi,\frac{1}{\sqrt2})+C$。这避免将$1-x^4=(1-x^2)(1+x^2)$直接记成虚模数而未作说明。\par \textbf{求导验算} $\phi'=\frac{\sqrt2}{(1+x^2)\sqrt{1-x^2}}$，$\sqrt{1-\frac12\sin^2\phi}=\frac{1}{\sqrt{1+x^2}}$，相乘得$\frac{1}{\sqrt{1-x^4}}$。\par \textbf{疑难点与方法比较} 本题保留真正椭圆积分；与分子匹配导数的伪椭圆题分开识别。\end{example}

\subsection{初等表达的边界与最终验算}\label{ic:indefinite_special_complete:boundary}\begin{method}{初等表达的边界与最终验算}\textbf{识别信号} 已完成有理化、分部与代数约化，但仍剩标准非初等核。\par \textbf{操作步骤} 给出已定义的特殊函数或确定的归一化积分；核验导数；写出实定义域及分支。计算机的形式输出也按同一标准复核。\par \textbf{疑难点与条件} 不定积分的所有可能技巧不能用有限清单穷举。常用方法可按结构组织为可执行的算法，但并非任意被积函数都有初等原函数。\end{method}

\begin{warning}常见基本非初等核包括$\e^{-x^2}$、$\sin(x^2)$、$\frac{\sin x}{x}$、$\frac{\e^x}{x}$以及一般椭圆微分。表达式复杂不等于非初等，表达式简单也不保证初等。证明某类核不存在初等原函数通常需要Liouville理论等额外工具；“试了几种换元都失败”不能作为不可积证明。\end{warning}

\section{求解后的四项核验}\label{ic:verification}
\begin{method}{从答案反查全部步骤}
\textbf{第一项：微分。}把回代后的$F(x)$逐项求导，化成同一分母，验证$F'=f$；只做数值抽查不能替代恒等式。
\par\textbf{第二项：区间。}检查分母、对数和根式限制，记录在哪个连通区间使用同一个常数。
\par\textbf{第三项：分支。}反三角、反双曲和特殊函数均使用已声明分支；把两式合并时检查是否只在局部相差常数。
\par\textbf{第四项：退化。}检查参数为0、重根和换元分母为0的情况，回到未约化的表达式重新分类。
\end{method}
\begin{warning}疑难点：原函数表达式看似不同，可能只差常数；在连通区间上验证两者导数相同，再选一个点确定差值。不能据此把经过原奇点的不同区间强行合并。\end{warning}

\section{一元积分的组合构型}\label{ic:delta_integral}

\begin{examtip}本节结合部分分式、三角换元、区间再现和指数待定，分析综合积分中的结构选择。\end{examtip}

\subsection{倒数幂分母的通式与双解}\label{ic:delta_integral:reciprocalpower}\begin{method}{倒数幂分母的通式与双解}\textbf{识别信号} $\int\frac{\dd x}{x(x^n+a)}$，$a\ne0,n\ne0$，固定实幂定义区间。\par \textbf{操作步骤} 先取$t=x^n$，得到$\frac{\dd t}{nt(t+a)}$；或先拆成$a^{-1}[\frac{1}{x}-\frac{x^{n-1}}{x^n+a}]$。\par \textbf{疑难点与条件} 对数用绝对值，且$x\ne0,x^n+a\ne0$；实幂支路须可微。\end{method}

\begin{example}{考研风格｜一个通式的两种推导}{难度：$\star\star\star$}\label{ic:ex:delta_integral:reciprocalformula}\textbf{题目} 求$\int\frac{\dd x}{x(x^n+a)}$。\par \textbf{思路} 按上述两条路线推导。\par \textbf{详细解答} 令$t=x^n$，得到$\frac1{an}[\ln|t|-\ln|t+a|]+C=\frac1{an}\ln|\frac{x^n}{x^n+a}|+C$。直接拆分法给$\frac{\ln|x|}{a}-\frac{\ln|x^n+a|}{an}+C$，在$x>0$或相应实支上等价。\par \textbf{验算与条件复核} 对最后两项求导为$\frac{1}{ax}-\frac{x^{n-1}}{a(x^n+a)}=\frac{1}{x(x^n+a)}$。\par \textbf{疑难点与方法比较} $a=0$需另积$x^{-n-1}$，不能把0代进通式。\end{example}

\begin{example}{竞赛风格｜四次分母的对称组合}{难度：$\star\star\star$}\label{ic:ex:delta_integral:quarticcombination}\textbf{题目} 求$\int\frac{(1+x^2)\dd x}{1+x^4}$。\par \textbf{思路} 分母除以$x^2$后与$t=x-\frac{1}{x}$配微分。\par \textbf{详细解答} 在$x\ne0$时，原式$\frac{\dd(x-\frac{1}{x})}{(x-\frac{1}{x})^2+2}$，所以原函数$\frac1{\sqrt2}\arctan(\frac{x-\frac{1}{x}}{\sqrt2})+C$。因为原被积函数在0光滑，这个局部表达需在相邻半轴调整常数后连续连接；二次因子分解还给出全实轴连续的表达$\frac{\arctan(\sqrt2x+1)+\arctan(\sqrt2x-1)}{\sqrt2}+C$。因为令$A=x^2+\sqrt2x+1,B=x^2-\sqrt2x+1$，原式$=\frac12(A^{-1}+B^{-1})$，两项分别配方即可得到该结果。\par \textbf{验算与条件复核} $t'=1+x^{-2}$，分母$t^2+2=x^2+x^{-2}$，比值恢复原式。\par \textbf{疑难点与方法比较} 换元在0的极点不是原被积函数的奇点；写全实原函数需处理反正切分支。\end{example}

\begin{example}{竞赛风格｜指数复合的待定多项式}{难度：$\star\star\star$}\label{ic:ex:delta_integral:exponentialrecurrence}\textbf{题目} 求$\int(1+x-x^{-1})\e^{x+\frac{1}{x}}\dd x$，$x\ne0$。\par \textbf{思路} 设原函数$P(x)\e^{x+\frac{1}{x}}$并比较最高、最低幂。\par \textbf{详细解答} 取$P=x$，则$P'+(1-x^{-2})P=1+x-x^{-1}$，恰好匹配。故原函数$x\e^{x+\frac{1}{x}}+C$。这也可直接看作乘积求导的逆运算。\par \textbf{验算与条件复核} 求导为$\e^{x+\frac{1}{x}}+x\e^{x+\frac{1}{x}}(1-x^{-2})$。\par \textbf{疑难点与方法比较} 不是所有指数复合都无初等原函数，先检查线性微分组合。\end{example}

\section{图解：三角换元的几何关系}
\begin{figure}[htbp]\centering
\begin{tikzpicture}[scale=1.2,>=stealth]
\draw[thick,Ink](0,0)--(4,0)--(4,2.4)--cycle;
\draw(3.72,0)--(3.72,.28)--(4,.28);
\draw[Teal,->](.7,0) arc[start angle=0,end angle=30.96,radius=.7];
\node at (1,.24){$\theta$};
\node[below]at(2,0){$a>0$};\node[right]at(4,1.2){$x$};
\node[above,rotate=31]at(2,1.2){$\sqrt{a^2+x^2}$};
\node[align=left,anchor=west]at(5,.9){$\tan\theta=\frac{x}{a}$\\$x=a\tan\theta$\\$\dd x=a\sec^2\theta\,\dd\theta$};
\end{tikzpicture}
\caption{对$x>0$，直角三角形给$\sqrt{a^2+x^2}=a\sec\theta$。对负$x$用$\theta\in(-\frac{\pi}{2},0)$的有向变量仍有$\sec\theta>0$，故全体实$x$可用$\theta=\arctan(\frac{x}{a})$。图形只帮助识别恒等式，微分因子仍须实际计算。}\label{fig:review:substitution}
\end{figure}

\section{变限积分、临界判敛与几何建模}\label{ci:integrals}
\begin{examtip}
将被积函数的局部性质、积分区间的移动规律和微元的几何意义分别处理。变限积分的可导性须回到差商；含参数的反常积分先求定义域；不等式先寻找可积的点态控制；应用题先确定质量、距离和方向。
\end{examtip}

\subsection{振荡被积函数与平均差商}\label{ci:oscillatory}
\begin{method}{连续性充分条件之外的变限求导}
\textbf{适用条件} 被积函数可积，但在求导点不连续。\par
\textbf{步骤} 写成$\frac{F(x+h)-F(x)}{h}$，考察小区间上的平均值；对快速振荡项，用分部积分取得比区间长度更高阶的估计。\par
\textbf{疑难点} $F'(c)=f(c)$的基本定理以$f$在$c$连续为充分条件，不能仅因该条件不满足就判定$F$不可导。
\end{method}
\begin{example}{考研风格｜不连续被积函数仍给出可导积分}{难度：$\star\star\star$}\label{ci:ex:oscillatory}
\textbf{题目} 定义$f(0)=0$，$x\ne0$时$f(x)=\sin(\frac{1}{x})$，并令$F(x)=\int_0^x f(t)\dd t$。证明$F$在实轴上处处可导，求$F'(0)$。\par
\textbf{思路} 在非零点使用连续性，在零点直接估计积分差商。\par
\textbf{详细解答} $f$有界，任意有界区间上仅在0处可能不连续，因此Riemann可积。对$x\ne0$，连续性给$F'(x)=\sin(\frac{1}{x})$。先设$x>0$。在$[\varepsilon,x]$上用
\[
 \dd\!\left(\cos\dfrac1t\right)=\dfrac{\sin(\frac{1}{t})}{t^2}\dd t
\]
分部积分，再令$\varepsilon\downarrow0$。因$\varepsilon^2\cos(\frac{1}{\varepsilon})\to0$，得到
\[
 F(x)=x^2\cos\dfrac1x-2\int_0^x t\cos\dfrac1t\dd t,
 \qquad |F(x)|\le x^2+2\int_0^x t\dd t=2x^2.
\]
由于$f$为奇函数，$F(-x)=F(x)$，负半轴也满足$|F(x)|\le2x^2$。故
\[
 F'(0)=\lim_{x\to0}\dfrac{F(x)-F(0)}x=0=f(0).
\]
\textbf{方法点评} 振荡积分本身比$f$平滑；关键是$F(x)=O(x^2)$，而非错误地令$\sin(\frac{1}{x})$取极限。\par
\textbf{疑难点} 不能在0处直接套用要求被积函数连续的求导公式。改变$f(0)$不会改变$F$，因而也不会改变$F'(0)$。
\end{example}

\subsection{反常积分参数：先排极点，再谈数值}\label{ci:parameterdomain}
\begin{method}{参数积分的定义域先行}
\textbf{适用条件} 分母零点随参数移动，或端点包含幂、对数等临界因子。\par
\textbf{步骤} 求每个端点的等价式，同时列出区间内部的零点；各处分别判敛，取参数条件的交集；最后才使用原函数和代数方程。\par
\textbf{疑难点} 对数绝对值能写成实数，不代表跨越简单极点的反常积分存在。
\end{method}
\begin{example}{考研风格｜积分方程中的伪参数根}{难度：$\star\star$}\label{ci:ex:parameterpole}
\textbf{题目} 求实数$a$，使通常意义的反常积分
\[
 I(a)=\int_1^\infty\dfrac{a}{x(2x+a)}\dd x=\ln2.
\]
\textbf{思路} 先确定积分存在的参数范围，再计算。\par
\textbf{详细解答} 若$a<-2$，则$x=-\frac{a}{2}>1$为简单极点，积分发散；$a=-2$在左端点出现不可积的一阶极点。$a>-2$时分母在$[1,\infty)$上无零点，且被积式在无穷远为$O(x^{-2})$，积分收敛。此时
\begin{align*}
 I(a)&=\lim_{R\to\infty}
 \left[\ln x-\ln(2x+a)\right]_1^R\\
 &=\ln(2+a)-\ln2=\ln\left(1+\dfrac a2\right).
\end{align*}
于是$I(a)=\ln2$等价于$1+\frac{a}{2}=2$，唯一解为$a=2$。\par
\textbf{方法点评} 若未限制参数而先使用$\ln|2+a|$，会误得$a=-6$；该值在$x=3$处产生极点。\par
\textbf{疑难点} 简单极点两侧的对数发散不能靠相减抵消。通常反常积分要求两侧分别收敛。
\end{example}
\begin{example}{竞赛风格｜两个端点的对数条件相互排斥}{难度：$\star\star\star$}\label{ci:ex:logphase}
\textbf{题目} 对$m,n>0$，完整判定
\[
 J(m,n)=\int_1^\infty\dfrac{\dd x}{(\ln x)^m(1+x^n)}
\]
的收敛范围。\par
\textbf{思路} 左端点用$\ln x\sim x-1$，右端点用幂与对数的分层判别。\par
\textbf{详细解答} 在$x\downarrow1$时，
\[
 \dfrac1{(\ln x)^m(1+x^n)}\sim\dfrac1{2(x-1)^m},
\]
所以左端收敛当且仅当$m<1$。在无穷远，被积式等价于$x^{-n}(\ln x)^{-m}$。若$n>1$，充分大时$(\ln x)^{-m}\le1$，可直接比较$x^{-n}$得收敛。若$0<n<1$，由$(\ln x)^m\le x^{\frac{1-n}{2}}$最终成立，得到
\[
 x^{-n}(\ln x)^{-m}\ge x^{-\frac{n+1}{2}},
\]
而$\frac{n+1}{2}<1$，故发散。若$n=1$，令$t=\ln x$，尾部等价积分为$\int^\infty t^{-m}\dd t$，当且仅当$m>1$收敛。此条件与左端要求$m<1$不相容。综上
\[
 J(m,n)\text{收敛}\quad\Longleftrightarrow\quad 0<m<1,\quad n>1.
\]
\textbf{方法点评} 临界线$n=1$无法通过调大$m$挽救全积分，因为它同时破坏左端的可积性。\par
\textbf{疑难点} 必须对同一个参数同时施加两个端点的条件，不能把它们写成“或”。
\end{example}

\subsection{含绝对值核与移动窗口}\label{ci:kernelwindow}
\begin{method}{先找分界点，再对积分求导}
\textbf{适用条件} 核含$|t-x|$，或积分上下限同步移动。\par
\textbf{步骤} 绝对值题按$x$相对固定区间的位置分类；移动窗口题先用变限公式求导，把积分关系化成差分关系。\par
\textbf{疑难点} 随参数移动的分界点会带来端点项，须完整使用Leibniz公式后检查是否抵消。
\end{method}
\begin{example}{考研风格｜绝对值积分的最优中心}{难度：$\star\star\star$}\label{ci:ex:absolutekernel}
\textbf{题目} 求$\Phi(x)=\int_0^1t|t-x|\dd t$，并求其在$\R$上的最小值点。\par
\textbf{思路} $x$在$[0,1]$内时，于$t=x$处分开；区间外直接去绝对值。\par
\textbf{详细解答} 若$x\le0$，$|t-x|=t-x$，故$\Phi(x)=\frac{1}{3}-\frac{x}{2}$；若$x\ge1$，则$\Phi(x)=\frac{x}{2}-\frac{1}{3}$。对$0<x<1$，
\begin{align*}
 \Phi(x)
 &=\int_0^xt(x-t)\dd t+\int_x^1t(t-x)\dd t\\
 &=\dfrac{x^3}6+\left[\dfrac{t^3}3-\dfrac{xt^2}2\right]_x^1
 =\dfrac{x^3}3-\dfrac x2+\dfrac13.
\end{align*}
三段表达在0、1处函数值与一阶导数均相接。因此
\[
 \Phi'(x)=
 \begin{cases}
 -\dfrac{1}{2},&x<0,\\
 x^2-\dfrac{1}{2},&0\le x\le1,\\
 \dfrac{1}{2},&x>1.
 \end{cases}
\]
唯一极小点为$x_*=\frac{1}{\sqrt2}$，并为全局最小点，其值是$\frac{1}{3}-\frac{1}{3\sqrt2}$。第二种解释：$\Phi'=0$等价于$\int_0^{x_*}t\dd t=\frac12\int_0^1t\dd t$，最优中心是权密度$t$的中位点。\par
\textbf{方法点评} 绝对偏差最小对应中位点，平方偏差最小才对应加权平均点。\par
\textbf{疑难点} $\Phi$在1处没有二阶导数，不能因一阶导数连续就误判为处处二阶可导。
\end{example}
\begin{example}{考研风格｜差分关系决定滑动积分}{难度：$\star\star$}\label{ci:ex:window}
\textbf{题目} 连续函数$f$满足$f(x+2)-f(x)=x$且$\int_0^2f(x)\dd x=0$。求任意实数$a$时$I(a)=\int_a^{a+2}f(x)\dd x$。\par
\textbf{思路} 对窗口位置求导，无需求出$f$的全部表达式。\par
\textbf{详细解答} 由变限求导，
\[
 I'(a)=f(a+2)-f(a)=a,\qquad I(0)=0,
\]
所以$I(a)=\frac{a^2}{2}$。另一法先将区间拆成
\[
 I(a)-I(0)=\int_2^{a+2}f(x)\dd x-\int_0^af(x)\dd x
 =\int_0^a[f(t+2)-f(t)]\dd t=\dfrac{a^2}2.
\]
定向积分的可加性使第二法对$a<0$也成立。\par
\textbf{方法点评} 更一般地，若$f(x+T)-f(x)=q(x)$，则$I(a)-I(0)=\int_0^aq(t)\dd t$。\par
\textbf{疑难点} “积分窗口长度相同”不推出积分相等；需要周期性，或像本题这样掌握精确漂移。
\end{example}

\subsection{积分比较：权函数、累计量和导数包络}\label{ci:inequality}
\begin{method}{把积分差变成同号乘积}
\textbf{适用条件} 区间反射交换两个被积因子，或目标不等式含累计积分的幂。\par
\textbf{步骤} 对称比较先配$x$与$b-x$；累计量设$A(x)=\int_0^xf$，逐次求导寻找非负因子。\par
\textbf{疑难点} 严格不等式须确认乘积在正长度区间上严格为正；只有单点严格通常还需连续性。
\end{method}
\begin{example}{竞赛风格｜单调权下正弦与余弦的比较}{难度：$\star\star\star$}\label{ci:ex:trigweight}
\textbf{题目} 若$w$在$[0,\frac{\pi}{2}]$上连续且严格递减，证明
\[
 \int_0^{\frac{\pi}{2}}w(x)\cos x\dd x>
 \int_0^{\frac{\pi}{2}}w(x)\sin x\dd x.
\]
另比较$\int_0^1\frac{x}{1+x}\sin(\frac{\pi x}{2})\dd x$与$\int_0^1\frac{x}{1+x}\cos(\frac{\pi x}{2})\dd x$。\par
\textbf{思路} 前者用区间反射配对，后者识别权的单调方向相反。\par
\textbf{详细解答} 记第一式左右之差为$D$。将后半区间令$x=\frac{\pi}{2}-t$，得到
\[
 D=\int_0^{\frac{\pi}{4}}
 \left[w(t)-w\!\left(\dfrac\pi2-t\right)\right](\cos t-\sin t)\dd t>0.
\]
两因子在$(0,\frac{\pi}{4})$均严格为正。第二法用Chebyshev比较：$w,\cos$同向，$w,\sin$反向，且$\int_0^{\frac{\pi}{2}}\cos x\dd x=\int_0^{\frac{\pi}{2}}\sin x\dd x=1$，故前者大于共同平均基准，后者小于它。\par
对第二问，令$u=\frac{\pi x}{2}$，权函数变为$\widetilde w(u)=\frac{\frac{2u}{\pi}}{1+\frac{2u}{\pi}}$，它严格递增，因此比较方向反转：
\[
 \int_0^1\dfrac{x\sin(\frac{\pi x}{2})}{1+x}\dd x>
 \int_0^1\dfrac{x\cos(\frac{\pi x}{2})}{1+x}\dd x.
\]
\textbf{方法点评} 配对法直接显示差的符号；平均不等式能同时比较两个量与同一基准。\par
\textbf{疑难点} 不能逐点比较$\sin$与$\cos$后直接积分，它们在区间中点交换大小。
\end{example}
\begin{example}{竞赛风格｜非线性累计积分的两次单调化}{难度：$\star\star\star$}\label{ci:ex:cumulative}
\textbf{题目} 设$f\in C^1[0,1]$，$f(0)=0$且$0\le f'(x)\le1$。证明
\[
 \left(\int_0^1f(x)\dd x\right)^2\ge\int_0^1f(x)^3\dd x.
\]
\textbf{思路} 先证明累计面积控制$f^2$，再把这一控制作为另一个差函数的导数。\par
\textbf{详细解答} 由$f'\ge0$可知$f\ge0$。令$A(x)=\int_0^xf(t)\dd t$及$G(x)=2A(x)-f(x)^2$，则
\[
 G(0)=0,\qquad G'(x)=2f(x)[1-f'(x)]\ge0.
\]
故$2A(x)\ge f(x)^2$。再设$H(x)=A(x)^2-\int_0^xf(t)^3\dd t$，有
\[
 H(0)=0,\qquad H'(x)=2A(x)f(x)-f(x)^3=f(x)G(x)\ge0.
\]
于是$H(1)\ge0$，即所求。作为第二种观察，由$f'\le1$与$f(0)=0$可得$f(x)\le x$；对固定$x$，又有$f(t)\ge\max\{0,f(x)-(x-t)\}$，积分此三角形下界同样得到$A(x)\ge \frac{f(x)^2}{2}$。\par
\textbf{方法点评} 目标中的平方提示先对累计积分求导，三次幂则提示配成$f[2A-f^2]$。$f(x)=x$时两边均为$\frac{1}{4}$，常数不能加强。\par
\textbf{疑难点} $0\le f'\le1$不是$0\le f\le1$的替代写法；只给后者无法得到本题的三角形下界。
\end{example}
\begin{method}{双端点约束产生尖顶包络}
\textbf{适用条件} 两端点函数值已知，且有一致导数界。\par
\textbf{步骤} 从左端点和右端点分别应用中值定理，取较强的上界；积分分界点通常是区间中点。\par
\textbf{疑难点} 最优包络可能不可导，因此最优常数可以是上确界而没有光滑函数达到。
\end{method}
\begin{example}{竞赛风格｜导数约束下的最优面积常数}{难度：$\star\star\star$}\label{ci:ex:tent}
\textbf{题目} 设$f\in C^1[a,b]$，$a<b$，$M\ge0$，$f(a)=f(b)=0$且$|f'|\le M$。证明
\[
 \int_a^b|f(x)|\dd x\le\dfrac{M(b-a)^2}4,
\]
并说明常数$\frac{1}{4}$不可减小。\par
\textbf{思路} 用两端点给出的线性界夹住函数，得到尖顶形的上包络。\par
\textbf{详细解答} 由中值定理或积分估计，
\[
 |f(x)|\le M(x-a),\qquad |f(x)|\le M(b-x).
\]
取较小者，在中点$m=\frac{a+b}{2}$处分开积分：
\[
 \int_a^b|f|\le
 M\int_a^m(x-a)\dd x+M\int_m^b(b-x)\dd x
 =\dfrac{M(b-a)^2}4.
\]
为证明最优，令$h=\frac{b-a}{2}$，对$\varepsilon>0$取
\[
 f_\varepsilon(x)=M\!\left[
 \sqrt{h^2+\varepsilon^2}-\sqrt{(x-m)^2+\varepsilon^2}\right].
\]
其端点为0，且$|f_\varepsilon'|=\frac{M|x-m|}{\sqrt{(x-m)^2+\varepsilon^2}}\le M$。当$\varepsilon\downarrow0$时一致趋于$M[h-|x-m|]$，故积分趋$\frac{M(b-a)^2}{4}$。\par
\textbf{方法点评} 导数界是几何斜率约束；从两端同时估计，比只从一个端点估计所得的$\frac{1}{2}$常数更精确。图\ref{ci:fig:tent}给出另一条数形结合的计算：上包络是底为$b-a$、高为$\frac{M(b-a)}{2}$的三角形，面积正是$\frac{M(b-a)^2}{4}$。\par
\textbf{疑难点} 若$M>0$，达到上界须与尖顶包络处处相同，但其在中点不可导，故非零的$C^1$函数不能达到该上界。
\end{example}
\begin{figure}[htbp]\centering
\begin{tikzpicture}[x=8cm,y=4.4cm]
\fill[Teal!8] (0,0)--(.5,.5)--(1,0)--cycle;
\draw[->] (-.04,0)--(1.09,0) node[right]{$x$};
\draw[->] (0,-.03)--(0,.63) node[above]{$y$};
\draw[Teal,thick] (0,0)--(.5,.5)--(1,0);
\draw[Ink,thick,domain=0:1,samples=81]
 plot(\x,{sqrt(.25+.0064)-sqrt((\x-.5)*(\x-.5)+.0064)});
\draw[dashed,Gold] (.5,0)--(.5,.5);
\node[below] at (0,0){$0$};\node[below] at (.5,0){$\frac{1}{2}$};
\node[below] at (1,0){$1$};\node[left] at (0,.5){$\frac{1}{2}$};
\node[Teal,anchor=west] at (.6,.48){尖顶上包络};
\node[Ink,anchor=west] at (.62,.26){光滑逼近};
\end{tikzpicture}
\caption{在$[0,1]$、$M=1$时，双端点斜率界给出$\min(x,1-x)$。阴影三角形面积为$\frac{1}{4}$；光滑曲线取$\varepsilon=0.08$，逼近该面积但不达到。}\label{ci:fig:tent}
\end{figure}
\begin{example}{竞赛拓展｜零边界与二阶导数的积分代价}{难度：$\star\star\star$}\label{ci:ex:lyapunov}
\textbf{题目} 设$f\in C^2[a,b]$，$f(a)=f(b)=0$且$f(x)\ne0$对$a<x<b$成立。证明
\[
 \int_a^b\left|\dfrac{f''(x)}{f(x)}\right|\dd x\ge\dfrac4{b-a}.
\]
若左端为无穷，按扩展值理解。\par
\textbf{思路} 取$|f|$最大点，左右分别用中值定理，再把两处斜率差写成$f''$的积分。\par
\textbf{详细解答} 设$|f(c)|=K>0$，$a<c<b$。在$[a,c]$和$[c,b]$应用中值定理，存在$\xi\in(a,c)$、$\eta\in(c,b)$使
\[
 f'(\xi)=\dfrac{f(c)}{c-a},\qquad f'(\eta)=-\dfrac{f(c)}{b-c}.
\]
由于$|f(x)|\le K$，在内部非零点有$|\frac{f''}{f}|\ge\frac{|f''|}{K}$。非负积分的截断极限及基本定理给
\begin{align*}
 \int_a^b\left|\dfrac{f''}{f}\right|\dd x
 &\ge\dfrac1K\int_\xi^\eta|f''(x)|\dd x
 \ge\dfrac{|f'(\eta)-f'(\xi)|}K\\
 &=\dfrac1{c-a}+\dfrac1{b-c}
 \ge\dfrac4{b-a}.
\end{align*}
最后一步由$(c-a)(b-c)\le\frac{(b-a)^2}{4}$得到。\par
\textbf{方法点评} 若$q$连续，非平凡解满足$y''+q(x)y=0$且相邻零点为$a,b$，则上式给$\int_a^b|q|\ge\frac{4}{b-a}$，可用于排除过近的零点。\par
\textbf{疑难点} $f$在端点为0，商式积分须按反常积分理解；不能把端点的$\frac{0}{0}$当成普通函数值。
\end{example}

\subsection{分层微元与斜轴旋转}\label{ci:geometry}
\begin{method}{重量、提升距离和截面面积同时入式}
\textbf{适用条件} 准静态抽水、竖直液体压力等连续介质模型，密度与重力加速度取常数。\par
\textbf{步骤} 先固定深度坐标；切片质量为$\rho A(y)\dd y$；抽水乘每片的提升距离，压力乘深度后对受压面积积分。\par
\textbf{疑难点} 抽水的距离和压力的深度是不同物理量，坐标原点改变后两者表达必须同步改变。
\end{method}
\begin{example}{考研风格｜圆锥抽水与半椭圆水压}{难度：$\star\star\star$}\label{ci:ex:water}
\textbf{题目} 各长度参数、密度$\rho$及重力加速度$g$均为正。（1）倒圆锥容器高$H$、口半径$R$，装满水，求抽到容器口所需的功。（2）半椭圆竖直挡板的直径在水面，其水平半轴为$b$、向下的竖直半轴为$a$，求一侧水压力。\par
\textbf{思路} 前者用水平水层，后者用水平窄条。\par
\textbf{详细解答} （1）以锥顶为$y=0$、容器口为$y=H$，截面半径$\frac{Ry}{H}$，每片水须提升$H-y$，所以
\begin{align*}
 W&=\dfrac{\rho g\pi R^2}{H^2}\int_0^Hy^2(H-y)\dd y\\
 &=\dfrac{\rho g\pi R^2}{H^2}
 \left(\dfrac{H^4}3-\dfrac{H^4}4\right)
 =\dfrac{\rho g\pi R^2H^2}{12}.
\end{align*}
也可用锥体水的总重$\frac{\rho g\pi R^2H}{3}$乘质心到口的距离$\frac{H}{4}$核对。\par
（2）以水面为深度$z=0$，宽度为$2b\sqrt{1-\frac{z^2}{a^2}}$，压力$p=\rho gz$。故
\[
 F=2\rho gb\int_0^az\sqrt{1-\frac{z^2}{a^2}}\dd z
 =\rho gba^2\int_0^1u^{\frac{1}{2}}\dd u
 =\dfrac23\rho gba^2,
\]
其中$u=1-\frac{z^2}{a^2}$，换元时上下限随负微分反转。\par
\textbf{方法点评} 同为水平分层，功需“重量乘位移”，压力需“压强乘面积”；两式的量纲分别是能量与力。\par
\textbf{疑难点} 若出口高出容器口$d$，第一问应另加总重乘$d$；忽略这一项相当于改变了题设的终点。
\end{example}
\begin{method}{先转正坐标，再绕斜轴旋转}
\textbf{适用条件} 平面曲线或区域绕平面内直线旋转，区域内部不跨越旋转轴；曲面参数化应避免重复覆盖。\par
\textbf{步骤} 取沿轴坐标$u$与带符号的垂直坐标$v$。半径为$|v|$，面积微元为$2\pi|v|\dd s$；区域的一次矩给旋转体积。\par
\textbf{疑难点} “到轴距离”不能用竖直距离代替。曲线多次生成同一曲面时，直接积分会计重。
\end{method}
\begin{example}{竞赛风格｜斜轴旋转曲面的参数面积}{难度：$\star\star\star$}\label{ci:ex:obliquearea}
\textbf{题目} 求曲线$y=\frac{x^3}{3}+2x$，$0\le x\le1$，绕直线$4x-3y=0$旋转一周生成曲面的面积。\par
\textbf{思路} 到斜轴的垂距配弧长，不能使用$2\pi y\dd s$。\par
\textbf{详细解答} 取正交坐标
\[
 u=\dfrac{3x+4y}5,\qquad v=\dfrac{-4x+3y}5.
\]
轴为$v=0$；曲线上$v=\frac{x(x^2+2)}{5}\ge0$，且$u'=\frac{11+4x^2}{5}>0$，故无回折重复覆盖。由$y'=x^2+2$，
\begin{align*}
 S&=2\pi\int_0^1v\sqrt{1+(y')^2}\dd x\\
 &=\dfrac{2\pi}5\int_0^1x(x^2+2)\sqrt{1+(x^2+2)^2}\dd x\\
 &\overset{t=x^2+2}{=}\dfrac\pi5\int_2^3t\sqrt{1+t^2}\dd t\\
 &=\dfrac\pi{15}\left(10\sqrt{10}-5\sqrt5\right).
\end{align*}
面积公式可独立推导：转正后曲面参数为$(u(x),v(x)\cos\theta,v(x)\sin\theta)$，两个切向量的叉积模为$v(x)\sqrt{u'(x)^2+v'(x)^2}$；正交变换保持弧长，积分$\theta\in[0,2\pi]$即得。\par
\textbf{方法点评} 正交变换同时解释旋转半径和面积元；最终的换元由$x(x^2+2)$与根式内层配合决定。\par
\textbf{疑难点} 这里只求曲线生成的侧面；若题目要求封闭旋转体的总表面积，还应判断是否有端面。
\end{example}
\begin{example}{竞赛风格｜一次矩化解斜轴旋转体积}{难度：$\star\star\star$}\label{ci:ex:pappus}
\textbf{题目} 区域$D$位于第一象限，由直线$x+y=1$与曲线$\sqrt x+\sqrt y=1$围成。求$D$绕$x+y=1$旋转一周生成的体积。\par
\textbf{思路} 旋转轴是边界，区域全部在同一侧，可将体积写成到轴距离的一次矩。\par
\textbf{详细解答} 区域可写为$0\le x\le1$、$(1-\sqrt x)^2\le y\le1-x$。其面积为
\[
 A=\int_0^1[1-x-(1-\sqrt x)^2]\dd x
 =\int_0^1(2\sqrt x-2x)\dd x=\dfrac13.
\]
关于$x,y$交换的对称性给$\bar x=\bar y$，并且
\[
 \int_Dx\dd A
 =\int_0^1x(2\sqrt x-2x)\dd x
 =\dfrac45-\dfrac23=\dfrac2{15},
 \qquad \bar x=\bar y=\dfrac25.
\]
正交转正后，区域上的微元在绕轴旋转时给体积$2\pi v\dd A$，其中$v=\frac{1-x-y}{\sqrt2}\ge0$。因此
\begin{align*}
 V&=\dfrac{2\pi}{\sqrt2}\int_D(1-x-y)\dd A\\
 &=\dfrac{2\pi}{\sqrt2}\,A(1-\bar x-\bar y)
 =\dfrac{\sqrt2\pi}{15}.
\end{align*}
这也证明了本题使用的质心旋转公式$V=2\pi dA$，其中$d$为形心到轴距离。\par
\textbf{方法点评} 先算平面区域的一次矩，避免为斜轴逐段求圆盘半径；区域对称使两个一次矩相同。\par
\textbf{疑难点} 若轴穿过区域内部，两侧旋转可能重叠，不能把带符号的一次矩或两侧绝对距离积分直接当成并集体积。
\end{example}

\section{双重自变量的变限积分：端点项与内部参数项}\label{r2:leibniz}
遇到 $G(x)=\int_{a(x)}^{b(x)}\varphi(x,t)\dd t$，先区分外部变量 $x$ 与积分哑变量 $t$。常用充分条件是 $a,b\in C^1$，且 $\varphi,\varphi_x$ 在包含相关积分区间的紧矩形上连续。此时
\[
 G'(x)=\varphi(x,b(x))b'(x)-\varphi(x,a(x))a'(x)
       +\int_{a(x)}^{b(x)}\varphi_x(x,t)\dd t.
\]
可以合法地对内部参数求偏导再积分，但必须同时保留移动端点项；“绝不能对被积函数求偏导”是错误的。拆项或换元常能降低条件要求或简化计算，并不是唯一合法方法。反常积分的换序、求导须另核查一致可积控制。

\begin{example}{核函数含 $x-t$：拆项避开额外光滑要求}{基础训练}
若 $f$ 连续，$G(x)=\int_0^x(x-t)f(t)\dd t$。拆为
\[
 G(x)=x\int_0^xf(t)\dd t-\int_0^xtf(t)\dd t,
\]
于是 $G'(x)=\int_0^xf(t)\dd t+xf(x)-xf(x)=\int_0^xf(t)\dd t$，再得 $G''(x)=f(x)$。直接用 Leibniz 公式也得到相同结果：上端点的核为零，内部偏导为 $f(t)$。把整个被积函数简单代入 $t=x$ 会错误地得到零。
\end{example}

\begin{example}{核函数含 $xt$：换元与参数法交叉核对}{综合训练}
令 $H(x)=\int_0^x f(xt)\dd t$，$f$ 在涉及区间连续。对 $x\ne0$，代换 $s=xt$ 得
\[
 H(x)=\frac1x\int_0^{x^2}f(s)\dd s,\qquad
 H'(x)=2f(x^2)-\frac1{x^2}\int_0^{x^2}f(s)\dd s.
\]
这个表达对负 $x$ 也成立，换元方向由 $1/x$ 自动反映。在 $x=0$，$H(0)=0$，连续性给
\[
 \lim_{h\to0}\frac{H(h)}h=\lim_{h\to0}\frac1{h^2}\int_0^{h^2}f(s)\dd s=f(0),
\]
故 $H'(0)=f(0)$。若另外 $f\in C^1$，也可先写 $H'=f(x^2)+\int_0^x t f'(xt)\dd t$，再分部积分恢复同式；换元法不需要先假设 $f'$ 存在。
\end{example}

\paragraph{自测与解答}
求 $J(x)=\int_x^{x^2}(x+t^2)\dd t$ 的导数。端点法给
\[
 J'(x)=2x(x+x^4)-(x+x^2)+(x^2-x)
       =2x^5+2x^2-2x.
\]
先积分得到 $J=x^3-x^2+x^6/3-x^3/3$ 后求导同样成立。即使 $x^2<x$，有向积分公式仍可用，不应自行取绝对值。


\section{积分应用训练：从微元到总量}\label{r2:physics}
应用题先画图、定义坐标和单位，再写“微元量=密度或作用强度$\times$几何微元”。物理场景用于训练积分建模，具体哪些类别属于所考卷种，依卷首分流和当年大纲选择，不把本节全部题型标成每卷共同必考。

\subsection{旋转曲面侧面积：半径乘弧长微元}
若 $y=f(x)\ge0$、$f\in C^1[a,b]$，其图像绕 $x$ 轴旋转得到的侧面积为
\[
 S=2\pi\int_a^b f(x)\sqrt{1+[f'(x)]^2}\dd x.
\]
它不包含两端盖面。绕 $y$ 轴且 $x\ge0$ 时，半径改为 $x$。若曲线跨过旋转轴，半径应取绝对值，还要检查不同曲线段是否覆盖同一片曲面，避免重计。
例如 $y=x,0\le x\le1$ 绕 $x$ 轴旋转，$S=2\pi\sqrt2\int_0^1x\dd x=\pi\sqrt2$，与圆锥的 $\pi r l$（$r=1,l=\sqrt2$）一致。

\subsection{变力做功与抽水：距离不能选错}
沿 $x$ 方向的力分量为 $F(x)$，从 $a$ 到 $b$ 的功为 $W=\int_a^bF(x)\dd x$。力与位移反向时功为负。弹簧拉伸 $x\ge0$ 时外力 $kx$，从 $0$ 拉到 $L$ 需做功 $kL^2/2$；弹力本身做负功。

\begin{example}{直立圆柱容器抽水}{应用训练}
圆柱底面积为 $A$、高度 $H$，水深 $h\le H$，水的质量密度为 $\rho$，将水抽到顶部，忽略损耗。以底面为 $y=0$，厚度 $\dd y$ 的水片重量为 $\rho gA\dd y$，提升距离为 $H-y$。所以
\[
 W=\rho gA\int_0^h(H-y)\dd y
   =\rho gA(Hh-h^2/2).
\]
满水时为 $\rho gAH^2/2$，也等于总重乘质心提升距离 $H/2$。若容器变截面，应把常量 $A$ 换成 $A(y)$，不能沿用圆柱结果。
\end{example}

\subsection{液体静压力：压强乘竖直面面积微元}
以自由液面向下的深度 $y$ 为坐标，表压为 $\rho gy$；竖直闸门水平宽度为 $w(y)$，合力大小
\[
 F=\rho g\int_{y_1}^{y_2}y\,w(y)\dd y.
\]
这是竖直平面上各处平行压力的合力，不是液体的重量。宽度 $b$、从液面延伸到深度 $h$ 的矩形闸门，$F=\rho gbh^2/2$，作用深度为
\[
 y_{\mathrm{cp}}=\frac{\int_0^hy^2b\dd y}{\int_0^hyb\dd y}=\frac{2h}{3}.
\]
若两面都有同样大气压，求净力时可用表压；若题目要求单面绝对压力，应另加 $p_0$ 的贡献。

\subsection{细杆引力、薄板质心与转动惯量}
均匀细杆位于 $[0,L]$，线密度 $\lambda$，质量为 $m$ 的质点位于 $x=-a$（$a>0$）。杆上微元质量为 $\lambda\dd x$、距离为 $a+x$，引力沿同一方向，大小为
\[
 F=Gm\lambda\int_0^L\frac{\dd x}{(a+x)^2}
  =Gm\lambda\left(\frac1a-\frac1{a+L}\right).
\]
若 $a=0$，此理想线密度模型的积分发散，不能把有限公式直接延伸到接触点。

薄板在区域 $D$ 上有面密度 $\sigma(x,y)\ge0$、总质量 $M>0$，则
\[
 M=\iint_D\sigma\dd A,\quad
 \bar x=\frac{\iint_Dx\sigma\dd A}{M},\quad
 \bar y=\frac{\iint_Dy\sigma\dd A}{M}.
\]
绕 $x$ 轴、$y$ 轴及垂直板面过原点的轴，转动惯量分别为
\[
 I_x=\iint_Dy^2\sigma\dd A,\quad
 I_y=\iint_Dx^2\sigma\dd A,\quad I_O=I_x+I_y.
\]
惯量中的平方距离与质心的一次矩不能混淆。若 $D=[0,a]\times[0,b]$ 且密度常量 $\sigma$，则 $M=\sigma ab$，质心 $(a/2,b/2)$，$I_x=Mb^2/3,I_y=Ma^2/3$。

\paragraph{自测与解答}
\begin{enumerate}
\item 均匀圆盘半径 $R$，质量 $M$，求中心垂直轴惯量。答：极坐标给 $I=\sigma\int_0^{2\pi}\int_0^Rr^3\dd r\dd\theta=\sigma\pi R^4/2=MR^2/2$。
\item 闸门宽度 $b$、高度 $h$、上缘位于深度 $d$。答：$F=\rho gb\int_d^{d+h}y\dd y=\rho gbh(d+h/2)$，可用形心深度乘面积核对。
\end{enumerate}


\section*{章末分层习题}\addcontentsline{toc}{section}{章末分层习题}
\refstepcounter{exercise}\label{q:3:1}\par\noindent\textbf{习题\theexercise\quad 基础题} 求$\int \frac{x}{1+x^2}\dd x$。\quad\hyperref[a:3:1]{[答案]}\par
\refstepcounter{exercise}\label{q:3:2}\par\noindent\textbf{习题\theexercise\quad 基础题} 求$\int\ln x\dd x$，$x>0$。\quad\hyperref[a:3:2]{[答案]}\par
\refstepcounter{exercise}\label{q:3:3}\par\noindent\textbf{习题\theexercise\quad 基础题} 求$\int_0^1x^3\dd x$。\quad\hyperref[a:3:3]{[答案]}\par
\refstepcounter{exercise}\label{q:3:4}\par\noindent\textbf{习题\theexercise\quad 考研风格} 求$\int_0^\pi x\sin^2x\dd x$。\quad\hyperref[a:3:4]{[答案]}\par
\refstepcounter{exercise}\label{q:3:5}\par\noindent\textbf{习题\theexercise\quad 考研风格} 求$\int\frac{\dd x}{x^2+4}$。\quad\hyperref[a:3:5]{[答案]}\par
\refstepcounter{exercise}\label{q:3:6}\par\noindent\textbf{习题\theexercise\quad 考研风格} 判定$\int_0^1x^{-\frac{2}{3}}\dd x$。\quad\hyperref[a:3:6]{[答案]}\par
\refstepcounter{exercise}\label{q:3:7}\par\noindent\textbf{习题\theexercise\quad 竞赛拓展} 求$\int_0^{\frac{\pi}{2}}\sin^5x\dd x$。\quad\hyperref[a:3:7]{[答案]}\par
\refstepcounter{exercise}\label{q:3:8}\par\noindent\textbf{习题\theexercise\quad 竞赛拓展} 求$\int_0^1\ln x\dd x$。\quad\hyperref[a:3:8]{[答案]}\par
\refstepcounter{exercise}\label{q:3:9}\par\noindent\textbf{习题\theexercise\quad 竞赛拓展} 设$a<b,f\in C[a,b]$，证明$(\int_a^bf)^2\le(b-a)\int_a^bf^2$。\quad\hyperref[a:3:9]{[答案]}\par
\section*{章末方法速查}\addcontentsline{toc}{section}{章末方法速查}
\par\noindent\begin{tabularx}{\textwidth}{p{23mm}XXX}\toprule 方法&适用条件&典型信号&易错点\\\midrule
换元&链式结构&内层导数&漏微分因子\\分部&乘积&多项式降阶&循环移项符号\\对称&区间与权重对称&$a+b-x$&变限负号\\含参求导&可积支配&参数消分母&未查一致控制\\
\bottomrule\end{tabularx}
